% Leçon 2 — Grammaire : le comparatif;……跟/和/与......一样/相同;……跟/和/与......不一样/不同。 — Chinois Tle A — Cours signé OnBuch+
% Programme officiel MINESEC. Compiler avec : tectonic ZH02-grammaire-le-comparatif.tex
\newcommand{\DOCMATIERE}{Chinois}
\newcommand{\DOCNIVEAU}{Tle A}
\newcommand{\DOCMODULE}{Module 1 : Vie familiale et intégration sociale}
\newcommand{\DOCLECON}{ZH02}
\newcommand{\DOCTITRE}{Leçon 2 — Grammaire : le comparatif;……跟/和/与......一样/相同;……跟/和/与......不一样/不同。}
% OnBuch+ — préambule « cours signés » ENRICHI (compilé avec tectonic / XeLaTeX)
% Système visuel « Pop » d'OnBuch+ : fond crème (jamais blanc), bordures encre
% épaisses, ombres « dures » décalées, titres Archivo Black en capitales.
% Version MATHÉMATIQUES : graphes (pgfplots/tikz), boîtes pédagogiques
% (définition, propriété, méthode, attention, le savais-tu, exercice résolu,
% exercice, TP, bilan), page de garde et signature OnBuch+ ancrée en bas.
%
% Macros attendues dans main.tex AVANT \input{preamble} :
%   \DOCMATIERE  \DOCNIVEAU  \DOCMODULE  \DOCLECON (numéro)  \DOCTITRE
\documentclass[11pt]{article}

\usepackage{fontspec}
\usepackage[french]{babel} % français = langue principale ; le chinois via \zh{..} / textebox (xeCJK)
\usepackage{xeCJK}
\usepackage{pifont} % \ding{51} \ding{55} utilisés par les modèles
\usepackage[a4paper,margin=2cm,top=3.4cm,bottom=2.8cm,headheight=1.4cm,footskip=1.3cm]{geometry}
\usepackage{amsmath,amssymb,amsfonts}
\usepackage{mathtools} % \xmapsto, \coloneqq... souvent utilisés par les modèles
\usepackage{cancel} % simplifications barrées (bilans de forces)
% \abs{z} (module, valeur absolue) et \norm{u} (norme), utilisés par les modèles
\providecommand{\abs}{}\let\abs\relax\DeclarePairedDelimiter{\abs}{\lvert}{\rvert}
\providecommand{\norm}{}\let\norm\relax\DeclarePairedDelimiter{\norm}{\lVert}{\rVert}
\usepackage[table]{xcolor} % [table] : fournit aussi \rowcolors (alternance de lignes)
\usepackage{pagecolor}
\usepackage{tikz}
\usetikzlibrary{arrows.meta,positioning,calc,shapes.geometric,shapes.misc,shapes.arrows,decorations.pathmorphing,decorations.pathreplacing,decorations.markings,patterns,shadows,mindmap,trees,fit,backgrounds,matrix,intersections,angles,quotes}
\usepackage{pgfplots}
\usepackage[version=4]{mhchem} % équations nucléaires \ce{^{235}_{92}U}
\usepackage[european,siunitx]{circuitikz} % schémas électriques
\pgfplotsset{compat=1.18}
\usepgfplotslibrary{fillbetween,groupplots} % aires entre courbes (calcul intégral)
\usepackage{enumitem}
\usepackage{fancyhdr}
% [most] charge toutes les bibliothèques tcolorbox (listings, minted, raster,
% documentation...) et gonfle inutilement la mémoire TeX sur un document long
% (~300 boîtes) : on ne charge que ce qui est réellement utilisé.
\usepackage[skins,breakable]{tcolorbox}
\usepackage{graphicx}
\usepackage{colortbl}
\usepackage{booktabs}
\usepackage{tabularx}
\usepackage{diagbox} % case d'en-tête diagonale des tableaux à double entrée
% Colonnes extensibles centrée / gauche / droite (C, L, R), souvent utilisées par les modèles
\newcolumntype{C}{>{\centering\arraybackslash}X}
\newcolumntype{L}{>{\raggedright\arraybackslash}X}
\newcolumntype{R}{>{\raggedleft\arraybackslash}X}
\usepackage{array}
\usepackage{multicol}
\usepackage{multirow}
\usepackage{siunitx}
% parse-numbers=false : accepte aussi \SI{2,5 \times 10^{-3}}{...}
\sisetup{locale=FR,output-decimal-marker={,},group-minimum-digits=4,parse-numbers=false}
\DeclareSIUnit\ohm{\ensuremath{\symup{\Omega}}} % U+2126 absent de la police
\DeclareSIUnit\division{div} % réglages d'oscilloscope (ms/div, V/div)
\DeclareSIUnit\tr{tr} % tours (vitesse de rotation, ex. \SI{1500}{\tr\per\minute})
\DeclareSIUnit\molar{M} % concentration molaire (ex. \SI{3}{\molar}, \SI{1}{\micro\molar})
\DeclareSIUnit\an{an} % année (le modèle confond parfois avec \year, un compteur TeX) — ex. \SI{5}{\kilo\meter\per\an}
\DeclareSIUnit\FCFA{FCFA} % franc CFA (ex. \SI{500}{\FCFA\per\kilo\gram})
\DeclareSIUnit\degreeBrix{\textdegree Brix} % teneur en sucre d'un sirop/jus (réfractométrie)
\DeclareSIUnit\month{mois} % le modèle utilise parfois \month, un compteur TeX, comme unité de durée
\usepackage{qrcode}
% \degree : commande fréquemment utilisée par les modèles pour un angle ou une
% température en dehors de \SI{}{} (non fournie par les paquets chargés).
\providecommand{\degree}{^{\circ}}
% \qed (fin de démonstration), utilisé par les modèles sans amsthm
\providecommand{\qed}{\hfill$\square$}
% \barre : double barre des valeurs interdites dans un tableau de variations (tkz-tab)
\providecommand{\barre}{\ensuremath{\|}}
% environnement proof (amsthm non chargé) utilisé par les modèles
\ifdefined\proof\else\newenvironment{proof}[1][Démonstration]{\par\noindent\textit{#1.}\ }{\qed\par}\fi
\usepackage{lastpage}
\usepackage{hyperref}
\hypersetup{hidelinks}

% ============================================================
% Palette « Pop » OnBuch+ — mêmes hex que la classe P (plus_ui.dart)
% ============================================================
\definecolor{poporange}{HTML}{F59321}
\definecolor{popdark}{HTML}{E07A0C}
\definecolor{popcream}{HTML}{FFF6E8}
\definecolor{popink}{HTML}{1C1714}
\definecolor{poppurple}{HTML}{7A5AE0}
\definecolor{popgreen}{HTML}{2E9E6A}
\definecolor{poppink}{HTML}{E0457B}
\definecolor{popblue}{HTML}{2F7DE1}
\definecolor{popgold}{HTML}{C98A12}
\definecolor{popmuted}{HTML}{8A7F76}
\definecolor{popline}{HTML}{EADFD2}
\definecolor{popgray}{HTML}{8A7F76}
\colorlet{popgrey}{popgray}
% Alias tolérants (le modèle nomme parfois les couleurs différemment)
\colorlet{popwhite}{white}
\colorlet{popblack}{popink}
\colorlet{popred}{poppink}
\colorlet{popyellow}{popgold}
% Teintes claires (fonds de tableaux/encadrés) — le modèle nomme parfois
% les couleurs avec un suffixe L (ex. popblueL) sans les avoir définies.
\colorlet{poporangeL}{poporange!20}
\colorlet{popdarkL}{popdark!20}
\colorlet{poppurpleL}{poppurple!20}
\colorlet{popgreenL}{popgreen!20}
\colorlet{poppinkL}{poppink!20}
\colorlet{popblueL}{popblue!20}
\colorlet{popgoldL}{popgold!20}
\colorlet{popredL}{popred!20}
\colorlet{popgrayL}{popgray!20}
% Teintes claires pour fonds de tableaux / zones
\colorlet{popblueL}{popblue!10!white}
\colorlet{poporangeL}{poporange!14!white}
\colorlet{popgreenL}{popgreen!12!white}
\colorlet{poppurpleL}{poppurple!12!white}
\colorlet{poppinkL}{poppink!12!white}
\colorlet{popgoldL}{popgold!14!white}

\pagecolor{popcream}

% ============================================================
% Typographie
% ============================================================
\defaultfontfeatures{Ligatures=TeX}
\setmainfont{PlusJakartaSans-Medium.ttf}[Path=../../../fonts/,BoldFont=PlusJakartaSans-Bold.ttf]
% Symboles, API et emojis absents de la police principale : repli sur DejaVu Sans (caractères actifs)
\newfontface\symfont{DejaVuSans.ttf}[Path=../../../fonts/]
\makeatletter
\newcommand\symactive[1]{\catcode#1=13 \begingroup\lccode`\~=#1 \lowercase{\endgroup\def~}{{\symfont\char#1}}}
\newcommand\symblank[1]{\catcode#1=13 \begingroup\lccode`\~=#1 \lowercase{\endgroup\def~}{}}
\newcommand\symhyph[1]{\catcode#1=13 \begingroup\lccode`\~=#1 \lowercase{\endgroup\def~}{\mbox{-}}}
\newcount\symc
\newcommand\symrange[2]{\symc=#1 \@whilenum\symc<#2 \do{\expandafter\symactive\expandafter{\the\symc}\advance\symc by 1 }}
\newcommand\symblankrange[2]{\symc=#1 \@whilenum\symc<#2 \do{\expandafter\symblank\expandafter{\the\symc}\advance\symc by 1 }}
\makeatother
\symrange{"0250}{"02FF}\symactive{"2713}\symactive{"2714}\symactive{"2717}\symactive{"2718}\symactive{"2610}\symactive{"2611}\symactive{"2612}
\symrange{"2600}{"27BF}
\catcode"2705=13 \begingroup\lccode`\~="2705 \lowercase{\endgroup\def~}{{\symfont\char"2713}}\symblank{"26BD}\symblank{"26A1}
\symrange{"2460}{"257F}\symactive{"223C}\symactive{"2248}\symactive{"2260}
\symhyph{"2011}\symhyph{"2E1A}\symblankrange{"1F300}{"1FAFF}\symblank{"FE0F}
% Caractères chinois : WenQuanYi Zen Hei (sous-ensemble GB2312 + pinyin), gras/italique simulés
\setCJKmainfont{WenQuanYiZenHei-subset.ttf}[Path=../../../fonts/,AutoFakeBold=3,AutoFakeSlant=0.2]
\xeCJKsetup{CJKspace=false}
% Pinyin : voyelles à caron (ǎ ǐ ǒ ǔ ǖ ǘ ǚ ǜ) absentes de la police principale -> caractères actifs
\newfontface\pyfont{WenQuanYiZenHei-subset.ttf}[Path=../../../fonts/]
\newcommand\pyactive[1]{\catcode#1=13 \begingroup\lccode`\~=#1 \lowercase{\endgroup\def~}{{\pyfont\char#1}}}
\pyactive{"01CD}\pyactive{"01CE}\pyactive{"01CF}\pyactive{"01D0}\pyactive{"01D1}\pyactive{"01D2}\pyactive{"01D3}\pyactive{"01D4}\pyactive{"01D5}\pyactive{"01D6}\pyactive{"01D7}\pyactive{"01D8}\pyactive{"01D9}\pyactive{"01DA}\pyactive{"01DB}\pyactive{"01DC}
% Maths en sans-serif (Fira Math) avec chiffres et lettres droites de Plus
% Jakarta, pour que les formules se fondent dans le texte.
\usepackage[math-style=ISO,bold-style=ISO]{unicode-math}
\setmathfont{FiraMath-Regular.otf}[Path=../../../fonts/]
\setmathfont{PlusJakartaSans-Medium.ttf}[Path=../../../fonts/,range={up/num,up/{Latin,latin},\mathup}]
\usepackage{icomma} % virgule décimale sans espace en mode math
% Caractères mathématiques Unicode absents des polices de texte (Plus Jakarta,
% Archivo Black) : rendus par leur équivalent mathématique, en texte comme en
% formule — sinon ils s'affichent en carré vide (ex. « Arithmétique dans ℤ »).
\usepackage{newunicodechar}
\newunicodechar{ℝ}{\ensuremath{\mathbb{R}}}
\newunicodechar{ℤ}{\ensuremath{\mathbb{Z}}}
\newunicodechar{ℂ}{\ensuremath{\mathbb{C}}}
\newunicodechar{ℕ}{\ensuremath{\mathbb{N}}}
\newunicodechar{ℚ}{\ensuremath{\mathbb{Q}}}
\newunicodechar{𝔻}{\ensuremath{\mathbb{D}}}
\newunicodechar{α}{\ensuremath{\alpha}}
\newunicodechar{β}{\ensuremath{\beta}}
\newunicodechar{γ}{\ensuremath{\gamma}}
\newunicodechar{δ}{\ensuremath{\delta}}
\newunicodechar{ε}{\ensuremath{\varepsilon}}
\newunicodechar{θ}{\ensuremath{\theta}}
\newunicodechar{λ}{\ensuremath{\lambda}}
\newunicodechar{μ}{\ensuremath{\mu}}
\newunicodechar{σ}{\ensuremath{\sigma}}
\newunicodechar{τ}{\ensuremath{\tau}}
\newunicodechar{φ}{\ensuremath{\varphi}}
\newunicodechar{ψ}{\ensuremath{\psi}}
\newunicodechar{ω}{\ensuremath{\omega}}
\newunicodechar{Σ}{\ensuremath{\Sigma}}
% majuscules produites par \MakeUppercase dans les titres (α-aminés -> Α)
\newunicodechar{Α}{\ensuremath{\alpha}}
\newunicodechar{Β}{\ensuremath{\beta}}
\newunicodechar{Γ}{\ensuremath{\Gamma}}
\newunicodechar{Δ}{\ensuremath{\Delta}}
\newunicodechar{Ε}{\ensuremath{\varepsilon}}
\newunicodechar{Θ}{\ensuremath{\Theta}}
\newunicodechar{Λ}{\ensuremath{\Lambda}}
\newunicodechar{Μ}{\ensuremath{\mu}}
\newunicodechar{Τ}{\ensuremath{\tau}}
\newunicodechar{Φ}{\ensuremath{\Phi}}
\newunicodechar{Ψ}{\ensuremath{\Psi}}
\newunicodechar{Ω}{\ensuremath{\Omega}}
\newunicodechar{↦}{\ensuremath{\mapsto}}
\newunicodechar{∈}{\ensuremath{\in}}
\newunicodechar{∉}{\ensuremath{\notin}}
\newunicodechar{∧}{\ensuremath{\wedge}}
\newunicodechar{⇔}{\ensuremath{\Leftrightarrow}}
\newunicodechar{⟺}{\ensuremath{\Longleftrightarrow}}
\newunicodechar{⇒}{\ensuremath{\Rightarrow}}
\newunicodechar{⟹}{\ensuremath{\Longrightarrow}}
\newunicodechar{∘}{\ensuremath{\circ}}
\newunicodechar{∪}{\ensuremath{\cup}}
\newunicodechar{∩}{\ensuremath{\cap}}
\newunicodechar{≡}{\ensuremath{\equiv}}
\newunicodechar{⊂}{\ensuremath{\subset}}
\newunicodechar{⊥}{\ensuremath{\perp}}
\newunicodechar{∃}{\ensuremath{\exists}}
\newunicodechar{∀}{\ensuremath{\forall}}
\newunicodechar{∛}{\ensuremath{\sqrt[3]{\,}}}
\newunicodechar{∜}{\ensuremath{\sqrt[4]{\,}}}
\newunicodechar{⃗}{\ensuremath{{}^{\to}}}
\newunicodechar{ⁿ}{\ifmmode^{n}\else\textsuperscript{\ensuremath{n}}\fi}
\newunicodechar{ˣ}{\ifmmode^{x}\else\textsuperscript{\ensuremath{x}}\fi}
\newunicodechar{ᵇ}{\ifmmode^{b}\else\textsuperscript{\ensuremath{b}}\fi}
\newunicodechar{ʳ}{\ifmmode^{r}\else\textsuperscript{\ensuremath{r}}\fi}
\newunicodechar{ᵘ}{\ifmmode^{u}\else\textsuperscript{\ensuremath{u}}\fi}
\newunicodechar{ʸ}{\ifmmode^{y}\else\textsuperscript{\ensuremath{y}}\fi}
\newunicodechar{ᵃ}{\ifmmode^{a}\else\textsuperscript{\ensuremath{a}}\fi}
\newunicodechar{ᵉ}{\ifmmode^{e}\else\textsuperscript{\ensuremath{e}}\fi}
\newunicodechar{ᵏ}{\ifmmode^{k}\else\textsuperscript{\ensuremath{k}}\fi}
\newunicodechar{ᵖ}{\ifmmode^{p}\else\textsuperscript{\ensuremath{p}}\fi}
\newunicodechar{ᴺ}{\ifmmode^{N}\else\textsuperscript{\ensuremath{N}}\fi}
\newunicodechar{ᶜ}{\ifmmode^{c}\else\textsuperscript{\ensuremath{c}}\fi}
\newunicodechar{ʰ}{\ifmmode^{h}\else\textsuperscript{\ensuremath{h}}\fi}
\newunicodechar{ᵠ}{\ifmmode^{q}\else\textsuperscript{\ensuremath{q}}\fi}
\newunicodechar{ₙ}{\ifmmode_{n}\else\textsubscript{\ensuremath{n}}\fi}
\newunicodechar{ₐ}{\ifmmode_{a}\else\textsubscript{\ensuremath{a}}\fi}
\newunicodechar{ᵢ}{\ifmmode_{i}\else\textsubscript{\ensuremath{i}}\fi}
\newunicodechar{ᵣ}{\ifmmode_{r}\else\textsubscript{\ensuremath{r}}\fi}
\newunicodechar{ₖ}{\ifmmode_{k}\else\textsubscript{\ensuremath{k}}\fi}
\newunicodechar{ᵦ}{\ifmmode_{\beta}\else\textsubscript{\ensuremath{\beta}}\fi}
\newunicodechar{ₓ}{\ifmmode_{x}\else\textsubscript{\ensuremath{x}}\fi}
\newfontfamily\popdisplay{ArchivoBlack-Regular.ttf}[Path=../../../fonts/]
\newfontfamily\poplogo{SpaceGrotesk-Bold.ttf}[Path=../../../fonts/]
\newfontfamily\popsemi{PlusJakartaSans-SemiBold.ttf}[Path=../../../fonts/]
\newfontfamily\popxbold{PlusJakartaSans-ExtraBold.ttf}[Path=../../../fonts/]
\newfontfamily\popxboldls{PlusJakartaSans-ExtraBold.ttf}[Path=../../../fonts/,LetterSpace=3.0]

\setlist[enumerate]{leftmargin=1.7em,itemsep=5pt,topsep=4pt}
\setlist[itemize]{leftmargin=1.7em,itemsep=4pt,topsep=4pt,label={\color{poporange}\textbullet}}
\setlength{\parindent}{0pt}
\setlength{\parskip}{5pt}
\renewcommand{\arraystretch}{1.25}

% pgfplots : style Pop par défaut
\pgfplotsset{
  every axis/.append style={axis line style={popink,line width=1pt},tick style={popink},
    grid style={popline,line width=0.5pt},label style={font=\small},tick label style={font=\footnotesize}},
  popcurve/.style={poporange,line width=1.8pt,no markers,smooth},
}
% Même style disponible dans un \draw TikZ simple (hors pgfplots)
\tikzset{popcurve/.style={poporange,line width=1.8pt}}
\tikzset{popgrid/.style={popline,very thin}}
\colorlet{popcurve}{poporange} % le nom du style est parfois utilisé comme couleur

% ============================================================
% Pill / squiggle
% ============================================================
\newcommand{\poppill}[3]{%
  \tikz[baseline=-0.55ex]\node[draw=popink,line width=1.2pt,rounded corners=2.6pt,%
    fill=#2,inner xsep=6.5pt,inner ysep=3pt]{%
    \popxboldls\fontsize{7}{7}\selectfont\color{#3}\MakeUppercase{#1}};%
}
\newcommand{\popsquiggle}[1]{%
  \tikz[baseline=-0.4ex]{
    \draw[#1,line width=2.6pt,line cap=round]
      plot [smooth] coordinates {(0,0.09) (0.34,0.01) (0.68,0.21) (1.02,0.01) (1.36,0.10)};
  }%
}
% Mot-clé surligné (marqueur orange sous le texte)
\newcommand{\cle}[1]{{\popxbold\color{popdark}#1}}

% ============================================================
% En-tête / pied
% ============================================================
\pagestyle{fancy}
\fancyhf{}
\renewcommand{\headrulewidth}{0pt}
\renewcommand{\footrulewidth}{0pt}
\lhead{\raisebox{-0.15ex}{\poplogo\fontsize{13}{13}\selectfont\color{popink}OnBuch\color{poporange}+}}
\rhead{\poppill{\DOCMATIERE\ \textbullet\ \DOCNIVEAU\ \textbullet\ Leçon \DOCLECON}{white}{popink}}
\lfoot{\popxbold\fontsize{7.6}{7.6}\selectfont\color{popmuted}\MakeUppercase{Cours signé OnBuch+}\ \textbullet\ {\popsemi\DOCTITRE}}
\rfoot{\tikz[baseline=-0.6ex]\node[rounded corners=8.5pt,fill=popink,minimum height=17pt,minimum width=17pt,inner xsep=5pt,inner sep=0]{\popxbold\fontsize{8}{8}\selectfont\color{popcream}\thepage\,/\,\pageref*{LastPage}};}

% ============================================================
% Titres
% ============================================================
\newcommand{\chaptitre}[2]{%
  \begin{tcolorbox}[enhanced,colback=poporange,colframe=popink,boxrule=2.6pt,arc=11pt,
      drop shadow=popink,left=15pt,right=15pt,top=12pt,bottom=11pt,before skip=3pt,after skip=18pt]
    \poppill{#2}{popink}{popcream}\par\vspace{9pt}
    {\raggedright\hyphenpenalty=10000\exhyphenpenalty=10000\popdisplay\fontsize{22}{24}\selectfont\color{popcream}\MakeUppercase{#1}\par}
  \end{tcolorbox}
}
% Section numérotée : pastille orange avec le numéro + titre Archivo
\newcounter{popsec}
\newcommand{\coursec}[1]{%
  \stepcounter{popsec}\setcounter{popsub}{0}%
  \par\vspace{16pt}\noindent
  \begin{minipage}{\linewidth}
  \tikz[baseline=-0.7ex]\node[draw=popink,line width=1.6pt,fill=poporange,rounded corners=4pt,minimum size=22pt,inner sep=0]{\popdisplay\fontsize{12}{12}\selectfont\color{popcream}\thepopsec};%
  \hspace{8pt}{\popdisplay\fontsize{15}{17}\selectfont\color{popink}\MakeUppercase{#1}}\par
  \vspace{4pt}\noindent{\color{poporange}\rule{36pt}{3.6pt}}
  \end{minipage}\par\nopagebreak\vspace{8pt}
}
\newcounter{popsub}
\newcommand{\courssub}[1]{%
  \stepcounter{popsub}%
  \par\vspace{9pt}\noindent{\popxbold\fontsize{11.5}{13}\selectfont\color{popink}{\color{poporange}\thepopsec.\thepopsub}\hspace{6pt}#1}\par\nopagebreak\vspace{4pt}
}

% ============================================================
% Boîtes pédagogiques — même recette : fond blanc, bordure épaisse colorée,
% ombre dure encre, badge plein. Titre optionnel [#1].
% ============================================================
\tcbset{popbase/.style={breakable,enhanced,colback=white,boxrule=2.3pt,arc=9pt,
  shadow={3pt}{-3pt}{0pt}{popink},left=14pt,right=14pt,top=11pt,bottom=11pt,before skip=11pt,after skip=15pt,
  fontupper=\popsemi\fontsize{10.6}{14.5}\selectfont\color{popink},
  fontlower=\popsemi\fontsize{10.6}{14.5}\selectfont\color{popink}}}
% Coche dessinée (le glyphe ✓ n'existe pas dans les polices du document)
\newcommand{\popcheck}[1]{\tikz[baseline=-0.5ex]\draw[#1,line width=1.6pt,line cap=round,line join=round](-0.13,0.0)--(-0.04,-0.09)--(0.13,0.1);}
\providecommand{\coche}{\popcheck{popgreen}}
\providecommand{\resultat}[1]{\par\smallskip\noindent\fcolorbox{popgreen}{popgreenL}{\parbox{\dimexpr\linewidth-2\fboxsep-2\fboxrule\relax}{\textbf{Résultat.} #1}}\par\smallskip}
\newcommand{\popboxhead}[3]{\poppill{#1}{#2}{white}\ifx\relax#3\relax\else\hspace{6pt}{\popxbold\fontsize{10.6}{12}\selectfont\color{popink}#3}\fi\par\vspace{7pt}}

\newtcolorbox{aretenir}[1][]{popbase,colframe=popgreen,before upper={\popboxhead{À retenir}{popgreen}{#1}}}
% Texte en chinois (document de lecture, dialogue, extrait) : cadre
% d'étude. Un titre optionnel (source, auteur) s'affiche dans l'en-tête.
\newtcolorbox{textebox}[1][]{popbase,colframe=popblue,before upper={\popboxhead{文本 · Texte}{popblue}{#1}},after upper={}}
% Marques ✓ / ✗ (forme correcte / fautive) souvent utilisées par les modèles
\providecommand{\cross}{\textcolor{poppink}{\ensuremath{\times}}}
\providecommand{\xmark}{\cross}\providecommand{\crossmark}{\cross}\providecommand{\cmark}{\ensuremath{\checkmark}}
% Ligne de réponse / justification à compléter (inventée par les modèles)
\providecommand{\justification}[1]{\par\noindent\textit{Justification~:} \dotfill\par}
% \textipa (package tipa non chargé) : la police ne couvre pas l'API -> simple passage
\providecommand{\textipa}[1]{#1}
% Soulignement (ulem non chargé) : \uline, \ul
\providecommand{\uline}[1]{\underline{#1}}\providecommand{\ul}[1]{\underline{#1}}
% \glossaire{\item ...} inventée par les modèles : liste de glossaire
\providecommand{\glossaire}[1]{\begin{itemize}#1\end{itemize}}
% \alert (beamer) inventée par les modèles : texte mis en évidence
\providecommand{\alert}[1]{\textbf{\textcolor{poppink}{#1}}}
% variantes ASCII d'ordinaux écrites en macro
\providecommand{\iere}{\textsuperscript{re}}\providecommand{\ieme}{\textsuperscript{e}}\providecommand{\ere}{\textsuperscript{re}}\providecommand{\eme}{\textsuperscript{e}}\providecommand{\er}{\textsuperscript{er}}\providecommand{\ie}{\textsuperscript{e}}
% Ordinaux français écrits en macro par les modèles : 1\ière, 3\ième
\newcommand{\ière}{\textsuperscript{re}}\newcommand{\ième}{\textsuperscript{e}}
% \exemple (inventée par les modèles) : étiquette « Exemple : »
\providecommand{\exemple}{\textbf{Exemple~:}\ }
% \square utilisable aussi hors mode math (cases à cocher)
\AtBeginDocument{\let\squareMath\square\renewcommand{\square}{\ensuremath{\squareMath}}}
% Mot ou phrase en chinois à l'intérieur d'un paragraphe français
\newcommand{\zh}[1]{\ifmmode\text{#1}\else{#1}\fi}
\let\eng\zh \let\esp\zh \let\all\zh % alias tolérés
% flèches d'intonation inventées par les modèles
\providecommand{\textsearrow}{}\renewcommand{\textsearrow}{\ensuremath{\searrow}}\providecommand{\textnearrow}{}\renewcommand{\textnearrow}{\ensuremath{\nearrow}}
% \fr{..} : mot français (inventé par les modèles)
\providecommand{\fr}[1]{\textit{#1}}
% zhbox : encadré de texte chinois inventé par les modèles
\newenvironment{zhbox}{\begin{quote}}{\end{quote}}
% \vs : « versus » inventé par les modèles
\providecommand{\vs}{\textit{vs}\ }
% \blank : case à compléter inventée par les modèles
\providecommand{\blank}[1][]{\underline{\hspace{1.6em}}}
\newtcolorbox{exemplebox}[1][]{popbase,colframe=poporange,before upper={\popboxhead{Exemple}{poporange}{#1}}}
\newtcolorbox{definition}[1][]{popbase,colframe=popblue,before upper={\popboxhead{Définition}{popblue}{#1}}}
\newtcolorbox{propriete}[1][]{popbase,colframe=poppurple,before upper={\popboxhead{Propriété}{poppurple}{#1}}}
\newtcolorbox{methode}[1][]{popbase,colframe=popgold,before upper={\popboxhead{Méthode}{popgold}{#1}}}
\newtcolorbox{attention}[1][]{popbase,colframe=poppink,before upper={\popboxhead{Attention}{poppink}{#1}}}
\newtcolorbox{experience}[1][]{popbase,colframe=popblue,colback=popblueL,before upper={\popboxhead{Expérience}{popblue}{#1}}}
% « Le savais-tu ? » : carte violette pleine (contexte, culture, Cameroun)
\newtcolorbox{savaistu}[1][]{popbase,colback=poppurple,colframe=popink,
  fontupper=\popsemi\fontsize{10.4}{14.2}\selectfont\color{white},
  before upper={\poppill{Le savais-tu\,?}{white}{poppurple}\ifx\relax#1\relax\else\hspace{6pt}{\popxbold\color{white}#1}\fi\par\vspace{7pt}}}
% Exercice résolu : énoncé en haut, \tcblower puis la solution sur fond crème
\newcounter{popexo}\newcounter{popexr}
\newtcolorbox{exoresolu}[1][]{popbase,colframe=poporange,colbacklower=popcream,
  segmentation style={popink,line width=1.2pt,dashed},
  before upper={\stepcounter{popexr}\popboxhead{Exercice résolu \thepopexr}{poporange}{#1}},
  before lower={\poppill{Solution}{popink}{popcream}\par\vspace{6pt}}}
% Exercice (non résolu en place) — la correction est dans la partie Corrigés
\newtcolorbox{exercice}[1][]{popbase,colframe=popink,
  before upper={\stepcounter{popexo}\popboxhead{Exercice \thepopexo}{popink}{#1}}}
% Corrigé
\newtcolorbox{corrige}[1][]{popbase,colframe=popgreen,colback=popgreenL,
  before upper={\popboxhead{Corrigé}{popgreen}{#1}}}
% Objectifs / prérequis / bilan
\newtcolorbox{objectifs}{popbase,colframe=popink,colback=poporangeL,
  before upper={\popboxhead{Objectifs de la leçon}{popink}{}}}
\newtcolorbox{prerequis}{popbase,colframe=popink,colback=popblueL,
  before upper={\popboxhead{Prérequis}{popblue}{}}}
\newtcolorbox{bilan}{popbase,colframe=popink,colback=white,boxrule=2.8pt,
  before upper={\popboxhead{Fiche bilan}{poporange}{L'essentiel à connaître pour l'examen}}}
% Figure encadrée avec légende
\newtcolorbox{popfigure}[1][]{enhanced,breakable=false,colback=white,colframe=popline,boxrule=1.4pt,arc=8pt,
  left=10pt,right=10pt,top=10pt,bottom=8pt,before skip=10pt,after skip=12pt,halign=center,
  lower separated=false,halign lower=center,
  fontlower=\popsemi\fontsize{9}{11.5}\selectfont\color{popmuted},
  #1}
\newcommand{\legende}[1]{\par\vspace{6pt}{\popsemi\fontsize{9}{11.5}\selectfont\color{popmuted}#1}}

% ============================================================
% Signature OnBuch+
% ============================================================
\newcommand{\poplogoinline}[1]{{\poplogo\fontsize{#1}{#1}\selectfont\color{popink}OnBuch\color{poporange}+}}
\newcommand{\signatureonbuch}{%
  \begin{tcolorbox}[enhanced,colback=popink,colframe=popink,boxrule=2.2pt,arc=10pt,
      drop shadow=poporange,left=14pt,right=14pt,top=10pt,bottom=10pt,before skip=0pt,after skip=0pt]
    \tikz[baseline=-0.7ex]\node[circle,fill=popgreen,draw=popcream,line width=1.3pt,minimum size=22pt,inner sep=0]{\popcheck{white}};%
    \hspace{9pt}\begin{minipage}[c]{0.62\linewidth}
      {\popxbold\fontsize{12}{13}\selectfont\color{popcream}Signé\ \textbullet\ {\poplogo OnBuch\color{poporange}+}}\par\vspace{2pt}
      {\raggedright\popsemi\fontsize{8}{10.5}\selectfont\color{popline}\MakeUppercase{Équipe pédagogique OnBuch+}\\\MakeUppercase{Conforme au programme officiel MINESEC}\par}
    \end{minipage}\hfill
    {\poplogo\fontsize{22}{22}\selectfont\color{popcream}OnBuch\color{poporange}+}
  \end{tcolorbox}%
}

% ============================================================
% Page de garde
% #1 titre · #2 matière · #3 niveau · #4 module · #5 n° leçon · #6 durée ·
% #7 description / objectifs (texte ou itemize)
% ============================================================
\newcommand{\pagegarde}[7]{%
  \thispagestyle{empty}
  \newgeometry{a4paper,margin=1.7cm,top=1.6cm,bottom=1.4cm}%
  \begin{tikzpicture}[remember picture,overlay]
    \fill[poporange!18] ([xshift=-3.2cm,yshift=-3.6cm]current page.north east) circle (4.6cm);
    \fill[poppurple!14] ([xshift=2.4cm,yshift=7.6cm]current page.south west) circle (3.8cm);
  \end{tikzpicture}%
  {\poplogo\fontsize{30}{30}\selectfont\color{popink}OnBuch\color{poporange}+}\hfill
  \raisebox{0.6ex}{\poppill{Cours signé \textbullet\ Premium}{popink}{popcream}}\par
  \vspace{1pt}\popsquiggle{poppurple}\par
  \vspace{8pt}
  \begin{tcolorbox}[enhanced,colback=poporange,colframe=popink,boxrule=2.8pt,arc=13pt,
      drop shadow=popink,left=18pt,right=18pt,top=12pt,bottom=13pt,after skip=10pt]
    \poppill{#2 \textbullet\ #3}{popink}{popcream}\hspace{5pt}\poppill{#4}{white}{popink}\par\vspace{8pt}
    {\popdisplay\fontsize{14}{14}\selectfont\color{popink}LEÇON #5}\par\vspace{4pt}
    {\raggedright\hyphenpenalty=10000\exhyphenpenalty=10000\popdisplay\fontsize{25}{27}\selectfont\color{popcream}\MakeUppercase{#1}\par}
    \vspace{8pt}
    {\popxbold\fontsize{9.5}{9.5}\selectfont\color{popink}\MakeUppercase{Durée conseillée : #6}}
  \end{tcolorbox}
  \begin{tcolorbox}[enhanced,colback=white,colframe=popink,boxrule=2.2pt,arc=10pt,
      drop shadow=popink,left=16pt,right=16pt,top=10pt,bottom=10pt,after skip=8pt]
    \poppill{Dans cette leçon}{poppurple}{white}\par\vspace{5pt}
    \popsemi\fontsize{9.3}{12.6}\selectfont\color{popink}#7
  \end{tcolorbox}
  \noindent\begin{minipage}[t]{0.487\linewidth}
    \begin{tcolorbox}[enhanced,colback=popgreen,colframe=popink,boxrule=2.2pt,arc=10pt,
        drop shadow=popink,left=11pt,right=11pt,top=10pt,bottom=10pt,height=3.1cm,valign=center]
      \tcbox[colback=white,colframe=popink,boxrule=1.3pt,arc=5pt,boxsep=0pt,nobeforeafter,
        left=3pt,right=3pt,top=3pt,bottom=3pt,valign=center]{\qrcode[height=1.9cm]{https://play.google.com/store/apps/details?id=com.luvvix.onbuch&hl=fr}}%
      \hspace{8pt}\begin{minipage}[c]{0.5\linewidth}
        {\popdisplay\fontsize{11}{12.5}\selectfont\color{white}\MakeUppercase{Télécharge OnBuch}}\par\vspace{4pt}
        {\raggedright\popsemi\fontsize{8.4}{11}\selectfont\color{white}Gratuit \textbullet\ Google Play\\Tuteur IA, annales, concours, résultats\par}
      \end{minipage}
    \end{tcolorbox}
  \end{minipage}\hfill
  \begin{minipage}[t]{0.487\linewidth}
    \begin{tcolorbox}[enhanced,colback=poppurple,colframe=popink,boxrule=2.2pt,arc=10pt,
        drop shadow=popink,left=11pt,right=11pt,top=10pt,bottom=10pt,height=3.1cm,valign=center]
      \tikz[baseline=-0.7ex]\node[circle,fill=white,draw=popink,line width=1.3pt,minimum size=1.6cm,inner sep=0]{\popdisplay\fontsize{24}{24}\selectfont\color{poppurple}+};%
      \hspace{8pt}\begin{minipage}[c]{0.6\linewidth}
        {\popdisplay\fontsize{11}{12.5}\selectfont\color{white}\MakeUppercase{Passe à OnBuch+}}\par\vspace{4pt}
        {\raggedright\popsemi\fontsize{8.4}{11}\selectfont\color{white}Tous les cours signés, Tuteur IA illimité, fiches PDF — onglet OnBuch+ de l'app\par}
      \end{minipage}
    \end{tcolorbox}
  \end{minipage}
  \vspace{8pt}
  \popcontenu
  \vfill
  \signatureonbuch
  \restoregeometry
  \newpage
}

% Rangée « Ce cours contient » (page de garde)
\newcommand{\poptile}[3]{% #1 couleur #2 gros texte #3 légende
  \begin{tcolorbox}[enhanced,colback=#1,colframe=popink,boxrule=2pt,arc=9pt,shadow={2.5pt}{-2.5pt}{0pt}{popink},
      left=6pt,right=6pt,top=9pt,bottom=9pt,halign=center,valign=center,height=1.75cm,width=\linewidth,nobeforeafter]
    {\popdisplay\fontsize{12.5}{13}\selectfont\color{white}\MakeUppercase{#2}}\par\vspace{3pt}
    {\centering\popsemi\fontsize{7.8}{9.5}\selectfont\color{white}#3\par}
  \end{tcolorbox}}
\newcommand{\popcontenu}{%
  \par\noindent{\popxboldls\fontsize{7.5}{7.5}\selectfont\color{popmuted}CE COURS SIGNÉ CONTIENT}\par\vspace{7pt}
  \noindent\begin{minipage}[t]{0.235\linewidth}\poptile{popblue}{Cours}{complet, illustré, pas à pas}\end{minipage}\hfill
  \begin{minipage}[t]{0.235\linewidth}\poptile{poporange}{Méthodes}{et exercices résolus}\end{minipage}\hfill
  \begin{minipage}[t]{0.235\linewidth}\poptile{poppink}{Exercices}{type examen + corrigés}\end{minipage}\hfill
  \begin{minipage}[t]{0.235\linewidth}\poptile{popgreen}{Fiche bilan}{l'essentiel à retenir}\end{minipage}\par}

% Sommaire
\newtcolorbox{sommaire}{enhanced,colback=white,colframe=popink,boxrule=2.2pt,arc=10pt,
  drop shadow=popink,left=18pt,right=18pt,top=13pt,bottom=13pt,before skip=6pt,after skip=20pt,
  before upper={\poppill{Sommaire}{poppurple}{white}\par\vspace{9pt}\popsemi\fontsize{10.6}{14.5}\selectfont\color{popink}}}

% Signature de fin de cours — poussée tout en bas de la dernière page
\newcommand{\findecours}{%
  \par\vfill\nopagebreak
  \begin{center}\popsquiggle{poporange}\end{center}\vspace{4pt}
  \signatureonbuch
}
\AtBeginDocument{\def\llbracket{\mathopen{[\mkern-3.5mu[}}\def\rrbracket{\mathclose{]\mkern-3.5mu]}}} % intervalles d'entiers (glyphe ⟦ absent de la police)
% Glyphes absents de Fira Math : redéfinis à partir de glyphes disponibles
\AtBeginDocument{%
  \def\setminus{\mathbin{\backslash}}\let\smallsetminus\setminus
  \def\vdots{\mathord{\vbox{\baselineskip3.2pt\lineskiplimit0pt\kern4pt\hbox{.}\hbox{.}\hbox{.}}}}%
  \def\ddots{\mathinner{\mkern1mu\raise6.5pt\hbox{.}\mkern2mu\raise3.6pt\hbox{.}\mkern2mu\raise0.7pt\hbox{.}\mkern1mu}}%
  \def\triangle{\mathord{\scalebox{0.9}{$\symup{\Delta}$}}}%
  \def\nabla{\mathord{\raisebox{\depth}{\scalebox{1}[-1]{$\symup{\Delta}$}}}}%
  \def\checkmark{\popcheck{popdark}}%
  \def\star{\mathbin{\ast}}\def\bigstar{\mathord{\ast}}%
  \def\diamond{\mathbin{\scalebox{0.6}{$\Diamond$}}}%
  \def\bigcup{\mathop{\mathchoice{\raisebox{-0.3em}{\scalebox{1.7}{$\cup$}}}{\scalebox{1.25}{$\cup$}}{\cup}{\cup}}\displaylimits}%
  \def\bigcap{\mathop{\mathchoice{\raisebox{-0.3em}{\scalebox{1.7}{$\cap$}}}{\scalebox{1.25}{$\cap$}}{\cap}{\cap}}\displaylimits}%
  \def\lesseqqgtr{\mathrel{\lessgtr}}\def\bigoplus{\mathop{\oplus}\displaylimits}%
}
\providecommand{\ssi}{si et seulement si} % abréviation fréquente
\AtBeginDocument{\providecommand{\wideparen}{\overparen}} % arc de cercle

%% FIGFIX-BEGIN v3 — correctifs globaux des figures (docs/onbuch-generation/tools/figfix)
\usepackage{adjustbox}\usepackage{etoolbox}
% toute figure TikZ/pgfplots plus large que la ligne est réduite à la largeur disponible
\BeforeBeginEnvironment{tikzpicture}{\begin{adjustbox}{max width=\linewidth}}
\AfterEndEnvironment{tikzpicture}{\end{adjustbox}}
% légendes pgfplots lisibles par-dessus les courbes
\pgfplotsset{every axis/.append style={legend style={fill=white,fill opacity=0.92,text opacity=1,draw=black!15,inner sep=3pt}}}
% étiquettes d'arêtes (syntaxe edge["..."]) avec fond blanc
\tikzset{every edge quotes/.append style={fill=white,inner sep=1.2pt}}
% bulles de cartes mentales : texte à la ligne dans la bulle, bulle un peu plus grande
\tikzset{every concept/.append style={text width=2.0cm,align=center,minimum size=2.3cm,inner sep=2pt}}
%% FIGFIX-END
\begin{document}

\pagegarde{Leçon 2 — Grammaire : le comparatif;……跟/和/与......一样/相同;……跟/和/与......不一样/不同。}{Chinois}{Tle A}{Module 1 : Vie familiale et intégration sociale}{ZH02}{2 h}{Cette leçon de grammaire enseigne les structures chinoises de comparaison d'égalité et de différence : ……跟/和/与……一样/相同 et ……跟/和/与……不一样/不同. L'élève apprend à construire, transformer et utiliser ces structures dans des phrases et de courts textes, en évitant les calques du français.
\vspace{4pt}
\begin{multicols}{2}\small
\begin{itemize}
\item À la fin de cette leçon, je sais identifier dans un texte les structures ……跟/和/与……一样/相同 et ……跟/和/与……不一样/不同.
\item À la fin de cette leçon, je sais construire une phrase d'égalité avec 跟/和/与 + 一样/相同 en respectant l'ordre des mots chinois.
\item À la fin de cette leçon, je sais construire une phrase de différence avec 不一样/不同 et la négation 不.
\item À la fin de cette leçon, je sais choisir correctement entre 跟, 和 et 与 selon le registre (oral / écrit).
\item À la fin de cette leçon, je sais ajouter un adjectif ou un verbe après 一样 pour préciser le point de comparaison (一样高, 跟……一样喜欢).
\item À la fin de cette leçon, je sais transformer une phrase affirmative en phrase négative et en question.
\end{itemize}\end{multicols}}

\begin{sommaire}
\begin{multicols}{2}
\begin{enumerate}
\item Observation guidée d'exemples authentiques
\item Schéma de la structure d'égalité : A 跟/和/与 B 一样/相同
\item Schéma de la différence : A 跟/和/与 B 不一样/不同
\item Comparaison avec le français et pièges des francophones
\item Emplois étendus et exceptions (complément utile à l'examen)
\item Exercices d'application et de transformation
\item Activité d'intégration
\item Méthodes et exercices résolus
\item Exercices
\item Corrigés des exercices
\item Fiche bilan
\end{enumerate}
\end{multicols}
\end{sommaire}

\begin{objectifs}
\begin{itemize}
\item À la fin de cette leçon, je sais identifier dans un texte les structures ……跟/和/与……一样/相同 et ……跟/和/与……不一样/不同.
\item À la fin de cette leçon, je sais construire une phrase d'égalité avec 跟/和/与 + 一样/相同 en respectant l'ordre des mots chinois.
\item À la fin de cette leçon, je sais construire une phrase de différence avec 不一样/不同 et la négation 不.
\item À la fin de cette leçon, je sais choisir correctement entre 跟, 和 et 与 selon le registre (oral / écrit).
\item À la fin de cette leçon, je sais ajouter un adjectif ou un verbe après 一样 pour préciser le point de comparaison (一样高, 跟……一样喜欢).
\item À la fin de cette leçon, je sais transformer une phrase affirmative en phrase négative et en question.
\item À la fin de cette leçon, je sais éviter les erreurs fréquentes des francophones (place de 一样, calque de « aussi... que », oubli de 跟/和).
\item À la fin de cette leçon, je sais comparer des personnes, des objets et des situations du quotidien camerounais dans un court texte de 5 à 8 phrases.
\end{itemize}
\end{objectifs}

\begin{prerequis}
\begin{itemize}
\item Les pronoms personnels et la structure de base Sujet + Verbe + Complément (Première)
\item L'adjectif qualificatif en chinois (很 + adjectif) et la négation 不 (Première)
\item La conjonction 和 « et » reliant deux noms (Première)
\item Le vocabulaire de la famille, de l'école et de la description physique (Module 1, Tle)
\item La formation des questions avec 吗 et 是不是 (Première)
\end{itemize}
\end{prerequis}

\newpage

\coursec{Leçon 2 — Grammaire : Le comparatif d'égalité et de différence}

\courssub{Observation guidée d'exemples authentiques}

\courssub{Situation d'accroche et problématique}

En classe de Terminale A au lycée de Yaoundé, deux amis comparent leurs résultats au Probatoire : \zh{「我哥哥跟我一样高，可是他英语不如我好。」} (Wǒ gēge gēn wǒ yíyàng gāo, kěshì tā Yīngyǔ bùrú wǒ hǎo.) — « Mon frère aîné est aussi grand que moi, mais il ne parle pas anglais aussi bien que moi. » Au marché Mokolo, une cliente compare deux pagnes : \zh{「这块布与那块布相同，但价格不一样。」} (Zhè kuài bù yǔ nà kuài bù xiāngtóng, dàn jiàgé bù yíyàng.) — « Ce pagne est identique à celui-là, mais le prix n'est pas le même. »

En français, nous disons « aussi... que », « identique à », « différent de » en plaçant l'adjectif \emph{entre} les deux termes de la comparaison. En chinois, la logique est radicalement différente : les deux éléments comparés (A et B) sont \emph{encadrés} par une préposition (\zh{跟}, \zh{和}, \zh{与}) et le mot de liaison (\zh{一样}, \zh{相同}, \zh{不一样}, \zh{不同}) se place \textbf{après} le second terme, souvent suivi de l'adjectif ou du verbe. Comment repérer et construire ces structures sans calquer le français ? C'est l'objet de cette première section d'observation.

\courssub{Lecture à voix haute des six phrases modèles}

Lisez chaque phrase à haute voix en prêtant attention aux tons (marqués sur les voyelles : ā á ǎ à). Repérez visuellement les « mots-outils » de la comparaison.

\begin{textebox}[Six phrases modèles pour l'observation]
\zh{1. 我跟我哥哥一样高。} \\
Wǒ gēn wǒ gēge yíyàng gāo. \\
« Je suis aussi grand que mon grand frère. » \\[6pt]
\zh{2. 杜阿拉和雅温得不一样大。} \\
Duālā hé Yǎwēndé bù yíyàng dà. \\
« Douala et Yaoundé ne sont pas aussi grandes l'une que l'autre / n'ont pas la même taille. » \\[6pt]
\zh{3. 这本书与那本书相同。} \\
Zhè běn shū yǔ nà běn shū xiāngtóng. \\
« Ce livre est identique à celui-là. » (registre écrit / formel) \\[6pt]
\zh{4. 我的学校跟你的学校不同。} \\
Wǒ de xuéxiào gēn nǐ de xuéxiào bùtóng. \\
« Mon école est différente de la tienne. » \\[6pt]
\zh{5. 妈妈和爸爸一样喜欢喝茶。} \\
Māma hé bàba yíyàng xǐhuan hē chá. \\
« Maman et papa aiment tous les deux boire du thé / Maman aime le thé autant que papa. » \\[6pt]
\zh{6. 这个手机跟那个手机不一样贵。} \\
Zhège shǒujī gēn nàge shǒujī bù yíyàng guì. \\
« Ce téléphone n'est pas aussi cher que celui-là / n'a pas le même prix que celui-là. »
\end{textebox}

\courssub{Repérage guidé : code couleur et analyse}

Appliquez le code couleur suivant sur vos cahiers (ou sur l'affichage projeté) :
\begin{itemize}
    \item \textbf{Rouge} : les prépositions de comparaison \zh{跟} (gēn), \zh{和} (hé), \zh{与} (yǔ).
    \item \textbf{Bleu} : les mots de liaison d'égalité \zh{一样} (yíyàng), \zh{相同} (xiāngtóng) et de différence \zh{不一样} (bù yíyàng), \zh{不同} (bùtóng).
    \item \textbf{Entourez} (ou soulignez en vert) les deux éléments comparés : le terme A (sujet) et le terme B (objet de la préposition).
\end{itemize}

\begin{popfigure}
\begin{center}
\begin{tabularx}{\linewidth}{|X|X|X|}
\hline
\rowcolor{popblueL} \textbf{Phrase chinoise (Caractères + Pinyin)} & \textbf{Mots-outils repérés (Couleur)} & \textbf{Traduction française} \\
\hline
\zh{我跟我哥哥一样高。} (Wǒ \textcolor{red}{gēn} wǒ gēge \textcolor{blue}{yíyàng} gāo.) & \textcolor{red}{跟} \dots \textcolor{blue}{一样} & Je suis aussi grand que mon grand frère. \\
\hline
\zh{杜阿拉和雅温得不一样大。} (Duālā \textcolor{red}{hé} Yǎwēndé \textcolor{blue}{bù yíyàng} dà.) & \textcolor{red}{和} \dots \textcolor{blue}{不一样} & Douala et Yaoundé n'ont pas la même taille. \\
\hline
\zh{这本书与那本书相同。} (Zhè běn shū \textcolor{red}{yǔ} nà běn shū \textcolor{blue}{xiāngtóng}.) & \textcolor{red}{与} \dots \textcolor{blue}{相同} & Ce livre est identique à celui-là. \\
\hline
\zh{我的学校跟你的学校不同。} (Wǒ de xuéxiào \textcolor{red}{gēn} nǐ de xuéxiào \textcolor{blue}{bùtóng}.) & \textcolor{red}{跟} \dots \textcolor{blue}{不同} & Mon école est différente de la tienne. \\
\hline
\zh{妈妈和爸爸一样喜欢喝茶。} (Māma \textcolor{red}{hé} bàba \textcolor{blue}{yíyàng} xǐhuan hē chá.) & \textcolor{red}{和} \dots \textcolor{blue}{一样} & Maman aime le thé autant que papa. \\
\hline
\zh{这个手机跟那个手机不一样贵。} (Zhège shǒujī \textcolor{red}{gēn} nàge shǒujī \textcolor{blue}{bù yíyàng} guì.) & \textcolor{red}{跟} \dots \textcolor{blue}{不一样} & Ce téléphone n'est pas aussi cher que celui-là. \\
\hline
\end{tabularx}
\end{center}
\legende{Tableau d'observation : repérage des mots-outils (rouge) et des liaisons (bleu) dans six phrases modèles.}
\end{popfigure}

\courssub{Carte mentale : la logique binaire de la comparaison}

La structure chinoise sépare nettement l'égalité de la différence. Observez la carte mentale ci-dessous : le nœud central est la préposition (\zh{跟/和/与}), les deux branches principales sont « ÉGALITÉ » et « DIFFÉRENCE », chacune avec ses deux variantes (oral/écrit).

\begin{popfigure}
\begin{center}
\begin{tikzpicture}[
    mindmap,
    grow cyclic,
    every node/.style={concept, execute at begin node=\hskip0pt},
    concept color=popblue!80,
    level 1/.append style={level distance=4.5cm, sibling angle=90},
    level 2/.append style={level distance=3.2cm, sibling angle=45},
    text=white,
    font=\small
]
\node[align=center]{COMPARAISON\\A \zh{跟/和/与} B}
    child [concept color=popgreen!70] {
        node {ÉGALITÉ}
        child [concept color=popgreen!50] { node[align=center]{\zh{一样}\\yíyàng\\(oral / courant)} }
        child [concept color=popgreen!50] { node[align=center]{\zh{相同}\\xiāngtóng\\(écrit / formel)} }
    }
    child [concept color=poporange!70] {
        node {DIFFÉRENCE}
        child [concept color=poporange!50] { node[align=center]{\zh{不一样}\\bù yíyàng\\(oral / courant)} }
        child [concept color=poporange!50] { node[align=center]{\zh{不同}\\bùtóng\\(écrit / formel)} }
    };
\end{tikzpicture}
\end{center}
\legende{Carte mentale : deux prépositions (A 跟/和/与 B) → deux sens (Égalité / Différence) × deux registres (Courant / Formel).}
\end{popfigure}

\begin{methode}[Comment repérer une structure de comparaison en chinois]
Contrairement au français où l'on cherche « aussi... que » ou « plus... que », en chinois la méthode infaillible est \textbf{descendante} :
\begin{enumerate}
    \item \textbf{Cherchez le mot de liaison final} (le « pivot ») : \zh{一样}, \zh{相同}, \zh{不一样}, \zh{不同}. Il se trouve \emph{après} le second terme et \emph{avant} l'adjectif/verbe éventuel.
    \item \textbf{Remontez vers la gauche} pour trouver la préposition : \zh{跟}, \zh{和} ou \zh{与}.
    \item \textbf{Identifiez les deux termes} :
    \begin{itemize}
        \item Terme A = Sujet de la phrase (avant la préposition).
        \item Terme B = Nom/Groupe nominal situé \emph{entre} la préposition et le mot de liaison.
    \end{itemize}
    \item \textbf{Déduisez le sens} : \zh{一样/相同} = égalité ; \zh{不一样/不同} = différence.
\end{enumerate}
\emph{Exemple appliqué à la phrase 2 :} Je vois \zh{不一样大} $\rightarrow$ je remonte : \zh{和} $\rightarrow$ Terme A = \zh{杜阿拉}, Terme B = \zh{雅温得}. Sens : différence de taille.
\end{methode}

\begin{savaistu}[Registres et géographie : 跟 vs 和 vs 与]
\begin{itemize}
    \item \zh{跟} (gēn) : \textbf{Le plus fréquent à l'oral}, surtout dans le nord de la Chine (Pékin, Tianjin, Dongbei). Il implique souvent une relation personnelle ou une action commune. Ex : \zh{我跟你去} (J'y vais avec toi / Je suis comme toi pour y aller).
    \item \zh{和} (hé) : \textbf{Neutre, universel}. Utilisé à l'oral comme à l'écrit, dans le nord comme dans le sud. C'est le choix le plus sûr pour un apprenant. Il coordonne simplement deux noms.
    \item \zh{与} (yǔ) : \textbf{Registre écrit, formel, administratif, littéraire}. Rare à l'oral courant. On le voit dans les journaux (\zh{人民日报}), les contrats, les discours officiels, la littérature classique. Ex : \zh{中非合作与发展} (Coopération et développement sino-africains).
\end{itemize}
\textbf{Astuce examen (Bac / HSK) :} Si la phrase contient \zh{与}, le contexte est formel $\rightarrow$ privilégiez \zh{相同} ou \zh{不同} plutôt que \zh{一样}/\zh{不一样}. Si la phrase contient \zh{跟}, le contexte est familier $\rightarrow$ \zh{一样}/\zh{不一样} sont naturels.
\end{savaistu}

\courssub{Question déclencheuse : position de \zh{一样} vs « aussi... que »}

\begin{attention}[Piège n°1 : L'ordre des mots inversé par rapport au français]
En français, la structure est : \textbf{A + aussi + adjectif + que + B}. L'adjectif est \emph{au milieu}.
En chinois, la structure est : \textbf{A + 预位 + B + 一样 + adjectif}. L'adjectif est \emph{à la fin}.

\begin{center}
\begin{tabularx}{\linewidth}{|X|X|X|}
\hline
\rowcolor{popblueL} \textbf{Langue} & \textbf{Structure schématique} & \textbf{Exemple} \\
\hline
Français & A \textbf{aussi} Adj \textbf{que} B & Je suis \textbf{aussi grand} \textbf{que} mon frère. \\
\hline
Chinois & A \zh{跟/和/与} B \textbf{\zh{一样}} Adj & \zh{我 \textcolor{red}{跟} 我哥哥 \textcolor{blue}{一样} 高。} \\
\hline
\end{tabularx}
\end{center}

\textbf{Conséquence majeure :} Ne traduisez jamais \zh{一样} par « aussi » placé devant l'adjectif. \zh{一样} signifie « de la même manière / pareil » et fonctionne comme un adverbe de manière placé \textbf{avant} l'adjectif/verbe qu'il modifie.
\begin{itemize}
    \item \zh{一样高} = « de la même manière grand » $\rightarrow$ « aussi grand ».
    \item \zh{一样喜欢} = « de la même manière aimer » $\rightarrow$ « aimer autant ».
\end{itemize}
\end{attention}

\courssub{Exercice résolu : Repérage et traduction}

\begin{exoresolu}[Analyse de phrases nouvelles]
\textbf{Énoncé :} Pour chacune des phrases ci-dessous : 
1. Entourez la préposition (\zh{跟/和/与}). 
2. Soulignez le mot de liaison (\zh{一样/相同/不一样/不同}). 
3. Identifiez A et B. 
4. Traduisez en français naturel.

\begin{enumerate}
    \item \zh{喀麦隆与中国不一样大。} (Kāmàilóng yǔ Zhōngguó bù yíyàng dà.)
    \item \zh{妹妹跟姐姐一样漂亮。} (Mèimei gēn jiějie yíyàng piàoliang.)
    \item \zh{这道菜和那道菜相同。} (Zhè dào cài hé nà dào cài xiāngtóng.)
    \item \zh{我的中文老师跟你的不同。} (Wǒ de Zhōngwén lǎoshī gēn nǐ de bùtóng.)
\end{enumerate}

\tcblower
\textbf{Solution détaillée :}

\begin{center}
\begin{tabularx}{\linewidth}{|X|X|X|X|X|}
\hline
\rowcolor{popblueL} \textbf{Phrase} & \textbf{Préposition} & \textbf{Liaison} & \textbf{A / B} & \textbf{Traduction} \\
\hline
1. \zh{喀麦隆与中国不一样大。} & \zh{与} (formel) & \zh{不一样} & A=\zh{喀麦隆} / B=\zh{中国} & Le Cameroun n'est pas aussi grand que la Chine (différence de taille). \\
\hline
2. \zh{妹妹跟姐姐一样漂亮。} & \zh{跟} (oral) & \zh{一样} & A=\zh{妹妹} / B=\zh{姐姐} & La cadette est aussi jolie que l'aînée. \\
\hline
3. \zh{这道菜和那道菜相同。} & \zh{和} (neutre) & \zh{相同} (formel) & A=\zh{这道菜} / B=\zh{那道菜} & Ce plat est identique à celui-là. \\
\hline
4. \zh{我的中文老师跟你的不同。} & \zh{跟} (oral) & \zh{不同} (formel) & A=\zh{我的中文老师} / B=\zh{你的} (ellipse de \zh{中文老师}) & Mon prof de chinois est différent du tien. \\
\hline
\end{tabularx}
\end{center}

\textbf{Observations pédagogiques :}
\begin{itemize}
    \item Phrase 1 : Mélange registre formel (\zh{与}) + courant (\zh{不一样}). Possible à l'oral soigné, mais \zh{不同} serait plus cohérent avec \zh{与}.
    \item Phrase 3 : \zh{和} + \zh{相同} = registre standard/écrit. Notez l'absence d'adjectif après \zh{相同} (verbe d'état « être identique »).
    \item Phrase 4 : Ellipse fréquente : \zh{你的} = \zh{你的中文老师}. La structure \zh{A 跟 B 不同} suffit, pas besoin d'adjectif.
\end{itemize}
\end{exoresolu}

\courssub{Synthèse de l'observation : ce qu'il faut retenir avant la règle formelle}

\begin{aretenir}[Bilan de l'observation (Section 1)]
\begin{itemize}
    \item La comparaison en chinois utilise une \textbf{préposition binaire} : \zh{A 跟/和/与 B}.
    \item Le \textbf{mot de liaison} (\zh{一样}, \zh{相同}, \zh{不一样}, \zh{不同}) se place \textbf{après B}, jamais entre A et B.
    \item L'\textbf{adjectif ou le verbe} (s'il y en a un) se place \textbf{après} le mot de liaison : \zh{A 跟 B 一样 Adj/V}.
    \item Choix du registre : \zh{跟/一样/不一样} (oral/courant) \textbullet\ \zh{和} (neutre) \textbullet\ \zh{与/相同/不同} (écrit/formel).
    \item \textbf{Erreur fatale à bannir :} \zh{*我一样高跟我哥哥} (calque du français « Je aussi grand que mon frère »). L'ordre chinois est immuable : \zh{Sujet - Préposition - Objet - Liaison - Adj/Verbe}.
\end{itemize}
\end{aretenir}

\coursec{Schéma de la structure d'égalité : A 跟/和/与 B 一样/相同}

\begin{propriete}[Structure de l'égalité (A = B)]
\textbf{Schéma de base :}
\begin{center}
\textbf{A + 跟/和/与 + B + 一样/相同 (+ Adjectif / Verbe)}
\end{center}

\textbf{Règles de construction :}
\begin{enumerate}
    \item \textbf{Préposition} : \zh{跟} (oral), \zh{和} (neutre), \zh{与} (formel). Invariables.
    \item \textbf{Mot de liaison} :
    \begin{itemize}
        \item \zh{一样} (yíyàng) : courant, oral. \textbf{Exige généralement un adjectif ou un verbe derrière lui} (sauf réponse courte : \zh{他们一样} « Ils sont pareils »).
        \item \zh{相同} (xiāngtóng) : formel, écrit. \textbf{Peut s'employer seul comme prédicat} (verbe d'état « être identique ») sans adjectif suivant.
    \end{itemize}
    \item \textbf{Adjectif/Verbe (optionnel)} : Se place \textbf{après} \zh{一样}. Après \zh{相同}, l'adjectif est rare (on dit \zh{相同} tout court ou \zh{内容相同} « contenu identique »).
\end{enumerate}

\textbf{Tableau d'exemples (Égalité) :}

\begin{center}
\begin{tabularx}{\linewidth}{|X|X|X|}
\hline
\rowcolor{popblueL} \textbf{Structure} & \textbf{Exemple (Caractères + Pinyin)} & \textbf{Traduction} \\
\hline
A \zh{跟} B \zh{一样} Adj & \zh{这碗米粉跟那碗一样好吃。} (Zhè wǎn mǐfěn gēn nà wǎn yíyàng hǎochī.) & Ce bol de nouilles de riz est aussi bon que celui-là. \\
\hline
A \zh{和} B \zh{一样} V & \zh{弟弟和哥哥一样喜欢足球。} (Dìdi hé gēge yíyàng xǐhuan zúqiú.) & Le petit frère aime le foot autant que le grand frère. \\
\hline
A \zh{与} B \zh{相同} & \zh{合同条款与草案相同。} (Hétóng tiáokuǎn yǔ cǎo'àn xiāngtóng.) & Les clauses du contrat sont identiques au projet. \\
\hline
A \zh{跟} B \zh{一样} (seul) & \zh{你的想法跟我的一样。} (Nǐ de xiǎngfǎ gēn wǒ de yíyàng.) & Ton idée est la même que la mienne. \\
\hline
\end{tabularx}
\end{center}
\end{propriete}

\coursec{Schéma de la structure de différence : A 跟/和/与 B 不一样/不同}

\begin{propriete}[Structure de la différence (A ≠ B)]
\textbf{Schéma de base :}
\begin{center}
\textbf{A + 跟/和/与 + B + 不一样/不同 (+ Adjectif / Verbe)}
\end{center}

\textbf{Règles de construction :}
\begin{enumerate}
    \item \textbf{Préposition} : Identiques à l'égalité (\zh{跟}, \zh{和}, \zh{与}).
    \item \textbf{Mot de liaison} :
    \begin{itemize}
        \item \zh{不一样} (bù yíyàng) : courant, oral. \textbf{Exige généralement un adjectif/verbe derrière}. \zh{不} (4e ton) + \zh{一} (change en 2e ton \zh{jí} devant 4e ton) + \zh{样} (4e ton) $\rightarrow$ \textbf{bù yí yàng}.
        \item \zh{不同} (bùtóng) : formel, écrit. \textbf{Peut s'employer seul comme prédicat} (« être différent »).
    \end{itemize}
    \item \textbf{Nuance sémantique} :
    \begin{itemize}
        \item \zh{不一样 + Adj} : « pas aussi [Adj] que » (différence de degré). Ex : \zh{不一样贵} « pas aussi cher ».
        \item \zh{不同} (seul) : « différent de » (différence de nature/identité).
    \end{itemize}
\end{enumerate}

\textbf{Tableau d'exemples (Différence) :}

\begin{center}
\begin{tabularx}{\linewidth}{|X|X|X|}
\hline
\rowcolor{popblueL} \textbf{Structure} & \textbf{Exemple (Caractères + Pinyin)} & \textbf{Traduction} \\
\hline
A \zh{跟} B \zh{不一样} Adj & \zh{今年的芒果跟去年不一样甜。} (Jīnnián de mángguǒ gēn qùnián bù yíyàng tián.) & Les mangues de cette année ne sont pas aussi sucrées que celles de l'an dernier. \\
\hline
A \zh{和} B \zh{不一样} V & \zh{我和你不一样看这个问题。} (Wǒ hé nǐ bù yíyàng kàn zhège wèntí.) & Je ne vois pas ce problème de la même façon que toi. \\
\hline
A \zh{与} B \zh{不同} & \zh{喀麦隆的气候与中国北方的不同。} (Kāmàilóng de qìhòu yǔ Zhōngguó běifāng de bùtóng.) & Le climat du Cameroun est différent de celui du nord de la Chine. \\
\hline
A \zh{跟} B \zh{不同} (seul) & \zh{这双鞋跟那双不同。} (Zhè shuāng xié gēn nà shuāng bùtóng.) & Cette paire de chaussures est différente de celle-là. \\
\hline
\end{tabularx}
\end{center}
\end{propriete}

\coursec{Comparaison avec le français et pièges des francophones}

\begin{attention}[Piège n°2 : Oublier la préposition (\zh{跟/和/与})]
En français, « que » fait office de marqueur unique. En chinois, la préposition est \textbf{obligatoire}.
\begin{itemize}
    \item \textbf{Faux :} \zh{*我哥哥一样高。} (Il manque \zh{跟我}).
    \item \textbf{Vrai :} \zh{我哥哥\textcolor{red}{跟我}一样高。}
\end{itemize}

\end{attention}

\begin{attention}[Piège n°3 : Confondre \zh{不一样} (degré) et \zh{不同} (nature)]
\begin{itemize}
    \item \zh{这两本书不一样贵。} = « Ces deux livres n'ont pas le même prix » (différence de degré sur l'adjectif \zh{贵}).
    \item \zh{这两本书不同。} = « Ces deux livres sont différents » (contenu, auteur, couverture... différence globale).
    \item \zh{这两本书不一样。} = « Ces deux livres ne sont pas pareils » (courant, peut porter sur le prix, la couleur, le contenu...).
\end{itemize}

\end{attention}

\begin{attention}[Piège n°4 : Le ton de \zh{一} dans \zh{不一样}]
\zh{一} (yī) change de ton devant un 4e ton : \zh{不} (bù, 4e) \zh{一} $\rightarrow$ \zh{jí} (2e) \zh{样} (yàng, 4e) $\rightarrow$ \textbf{bù yí yàng}. Ne dites pas *bù yī yàng.
\end{attention}

\begin{attention}[Piège n°5 : \zh{和} (hé) vs \zh{和} (huò/huó/hué)]
Le caractère \zh{和} a plusieurs lectures :
\begin{itemize}
    \item \textbf{hé} (2e ton) : préposition « avec », conjonction « et », comparatif « que » (notre leçon).
    \item \textbf{huò} (4e ton) : « mélanger », « médicament » (\zh{和药}).
    \item \textbf{huó} (2e ton) : « pétrir » (\zh{和面} « pétrir la pâte »).
    \item \textbf{hué} (2e ton, familier) : « mélange » (\zh{搅和}).
\end{itemize}
\textbf{En grammaire de comparaison, c'est toujours \zh{hé} (2e ton).}
\end{attention}

\coursec{Emplois étendus et exceptions (Complément utile à l'examen)}

\begin{propriete}[Cas particuliers et extensions]
\textbf{1. Ellipse du terme B (contexte clair) :}
\zh{这个颜色跟那个（颜色）一样。} $\rightarrow$ \zh{这个颜色跟那个一样。} (Ce couleur est la même que celle-là.)
\textit{Règle :} On garde le démonstratif (\zh{那个} $\rightarrow$ \zh{那个}) ou le pronom possessif (\zh{你的}, \zh{我的}).

\textbf{2. Comparaison avec un verbe d'action (même manière) :}
\zh{请跟我一样写。} (Qǐng gēn wǒ yíyàng xiě.) « Écrivez de la même manière que moi / comme moi. »
\zh{他不跟大家一样思考。} (Tā bù gēn dàjiā yíyàng sīkǎo.) « Il ne réfléchit pas comme tout le monde. »

\textbf{3. Structure \zh{A 与 B 相同/不同} comme attribut (devant nom) :}
Nécessite \zh{的} : \zh{与去年相同的错误} (les mêmes erreurs que l'an dernier), \zh{与众不同的风格} (un style différent de la foule / unique).

\textbf{4. Négation de l'égalité = Différence :}
\zh{不跟...一样} = \zh{跟...不一样}.
\zh{他不跟我一样去。} = \zh{他跟我不一样去。} (Il n'y va pas comme moi / Il y va différemment de moi).

\textbf{5. Comparatif d'infériorité (aperçu, hors programme strict mais fréquent) :}
\zh{A 不如 B + Adj} (A est moins [Adj] que B). Ex : \zh{我哥哥英语不如我好。} (Mon frère est moins bon en anglais que moi). \textit{Serra traité en Leçon 3}.
\end{propriete}

\coursec{Vocabulaire clé de la leçon}

\begin{center}
\begin{tabularx}{\linewidth}{|X|X|X|X|}
\hline
\rowcolor{popblueL} \textbf{Caractères} & \textbf{Pinyin} & \textbf{Nature} & \textbf{Traduction} \\
\hline
跟 & gēn & prép. & avec, que (comparatif), suivre \\
\hline
和 & hé & prép./conj. & avec, et, que (comparatif) \\
\hline
与 & yǔ & prép. & avec, que (comparatif, formel), donner \\
\hline
一样 & yíyàng & adv./adj. & pareil, de la même manière, aussi \\
\hline
相同 & xiāngtóng & verbe d'état & identique, même (formel) \\
\hline
不一样 & bù yíyàng & adv./adj. & pas pareil, pas aussi... que \\
\hline
不同 & bùtóng & verbe d'état & différent (formel) \\
\hline
比较 & bǐjiào & adv. & relativement, assez (pour nuancer) \\
\hline
人口 & rénkǒu & nom & population \\
\hline
文化 & wénhuà & nom & culture \\
\hline
词典 & cídiǎn & nom & dictionnaire \\
\hline
图书馆 & túshūguǎn & nom & bibliothèque \\
\hline
\end{tabularx}
\end{center}

\coursec{Écriture des caractères clés : radicaux et ordre des traits}

\begin{definition}[Analyse graphémique des mots-outils]
\begin{itemize}
    \item \zh{\textbf{跟} (gēn)} : Radical \zh{足} (zú, pied, 7 traits). Phonétique \zh{艮} (gèn). Ordre : pied (haut→bas, gauche→droite) puis \zh{艮}. Sens : suivre à pied $\rightarrow$ avec, que.
    \item \zh{\textbf{和} (hé)} : Radical \zh{禾} (hé, grain/riz, 5 traits). Phonétique \zh{口} (kǒu). Ordre : \zh{禾} (haut→bas) puis \zh{口}. Sens originel : mélanger le grain $\rightarrow$ harmonie, avec.
    \item \zh{\textbf{与} (yǔ)} : Radical \zh{一} (yī, un, 1 trait). Structure ancienne : main tenant un objet. Ordre : trait horizontal, puis le reste (haut→bas). Sens : donner, accorder $\rightarrow$ avec (formel).
    \item \zh{\textbf{同} (tóng)} : Radical \zh{口} (kǒu, bouche). Haut : \zh{冂} (jiōng) + \zh{一}. Ordre : haut (\zh{冂}, \zh{一}) puis \zh{口}. Sens : même bouche/parole $\rightarrow$ même, identique.
    \item \zh{\textbf{样} (yàng)} : Radical \zh{木} (mù, bois/arbre). Phonétique \zh{羊} (yáng). Ordre : \zh{羊} (haut→bas) puis \zh{木} (bas). Sens : modèle en bois $\rightarrow$ manière, aspect.
    \item \zh{\textbf{不} (bù)} : Radical \zh{一}. 4 traits. Ordre : haut→bas, gauche→droite. Attention : ton change (bù $\rightarrow$ bú devant 4e ton).
\end{itemize}
\end{definition}

\coursec{Expression orale : Mise en situation (Cameroun / Chine)}

\begin{experience}[Débat comparatif : « La vie ici et là-bas »]
\textbf{Consigne :} En binômes. L'un joue un élève camerounais, l'autre un élève chinois en échange à Yaoundé. Comparez 5 aspects de la vie quotidienne en utilisant \textbf{obligatoirement} les 4 structures : \zh{跟...一样}, \zh{和...相同}, \zh{跟...不一样}, \zh{与...不同}. Alternez les registres (oral/écrit).

\textbf{Pistes de comparaison (vocabulaire fourni) :}
\begin{itemize}
    \item Climat : \zh{气候} (qìhòu), \zh{热} (rè), \zh{湿} (shī), \zh{干燥} (gānzào).
    \item Transport : \zh{交通} (jiāotōng), \zh{堵车} (dǔchē), \zh{方便} (fāngbiàn), \zh{地铁} (dìtiě).
    \item Alimentation : \zh{饮食} (yǐnshí), \zh{辣} (là), \zh{主食} (zhǔshí - \zh{米饭} vs \zh{玉米糊/ndolé}), \zh{筷子} (kuàizi) vs \zh{主食} (mains/couverts).
    \item École : \zh{学校} (xuéxiào), \zh{作业} (zuòyè), \zh{校服} (xiàofú), \zh{升旗} (shēngqí).
    \item Fêtes : \zh{节日} (jiérì), \zh{春节} (Chūnjié) vs \zh{Fête nationale 20 mai}, \zh{放假} (fàngjià).
\end{itemize}

\textbf{Modèle de démarrage :}
\zh{A: 咱们喀麦隆的气候跟中国南方的一样热，可是不一样湿。} \\
\zh{B: 是的。喀麦隆的饮食与中国的饮食不同，我们吃主食用筷子，你们用手或勺子。} \\
\zh{A: 不过，足球的热情跟中国的一样高！}
\end{experience}

\coursec{Exercices d'application et de transformation}

\begin{exercice}[Exercice 1 : Construction guidée (Thème)]
Traduisez en chinois (caractères + pinyin). Choisissez la préposition et le mot de liaison adaptés au contexte indiqué.
\begin{enumerate}
    \item Douala est aussi peuplée que Yaoundé. (contexte oral, entre amis)
    \item Ce dictionnaire est identique à celui de la bibliothèque. (contexte écrit, description formelle)
    \item Mon portable n'est pas aussi cher que le tien. (contexte oral)
    \item La culture camerounaise est différente de la culture chinoise. (contexte écrit, essai)
    \item Ma sœur cadette parle chinois aussi bien que moi. (contexte neutre)
    \item Ce problème n'est pas le même que celui d'hier. (contexte oral)
\end{enumerate}
\end{exercice}

\begin{exercice}[Exercice 2 : Transformation de structures]
Transformez les phrases suivantes selon l'indication. Donnez le résultat en caractères + pinyin.
\begin{enumerate}
    \item \zh{这件衣服跟那件一样贵。} $\rightarrow$ Remplacez \zh{跟} par \zh{与} et \zh{一样} par \zh{相同} (adaptez la fin de phrase).
    \item \zh{我的观点与老师的相同。} $\rightarrow$ Passez au registre oral (\zh{跟} + \zh{一样}) et ajoutez l'adjectif \zh{正确} (zhèngquè, correct).
    \item \zh{北京的冬天和哈尔滨的不一样冷。} $\rightarrow$ Reformulez avec \zh{不同} (attention : structure change).
    \item \zh{这道题跟那道题不同。} $\rightarrow$ Reformulez avec \zh{不一样} + adjectif \zh{难} (nán, difficile).
\end{enumerate}
\end{exercice}

\begin{exercice}[Exercice 3 : Expression écrite guidée (Version + Thème court)]
\textbf{A.} Traduisez en français :
\begin{enumerate}
    \item \zh{喀麦隆的足球水平跟非洲强队一样高。}
    \item \zh{雅温得的交通与十年前不同，现在有公交车了。}
    \item \zh{学中文不一样难，只要方法跟别人一样好。}
\end{enumerate}

\textbf{B.} Rédigez un court paragraphe (5-6 phrases) en chinois pour présenter votre ville/village à un correspondant chinois. Utilisez au moins 3 structures comparatives étudiées (égalité et différence). Comparez : taille, climat, ambiance, nourriture, transport.
\end{exercice}

\begin{corrige}[Exercice 1 : Construction guidée]
\begin{enumerate}
    \item \zh{杜阿拉跟雅温得一样人口多。} (Duālā gēn Yǎwēndé yíyàng rénkǒu duō.) \emph{Note : « peuplée » $\rightarrow$ \zh{人口多} ou \zh{人多}.}
    \item \zh{这本词典与图书馆的那本相同。} (Zhè běn cídiǎn yǔ túshūguǎn de nà běn xiāngtóng.)
    \item \zh{我的手机跟你的不一样贵。} (Wǒ de shǒujī gēn nǐ de bù yíyàng guì.)
    \item \zh{喀麦隆的文化与中国的文化不同。} (Kāmàilóng de wénhuà yǔ Zhōngguó de wénhuà bùtóng.)
    \item \zh{妹妹和我说中文一样好。} (Mèimei hé wǒ shuō Zhōngwén yíyàng hǎo.)
    \item \zh{这个问题跟昨天的不一样。} (Zhège wèntí gēn zuótiān de bù yíyàng.)
\end{enumerate}
\end{corrige}

\begin{corrige}[Exercice 2 : Transformation de structures]
\begin{enumerate}
    \item \zh{这件衣服与那件相同。} (Zhè jiàn yīfu yǔ nà jiàn xiāngtóng.) \emph{Note : \zh{相同} s'emploie seul, on supprime \zh{贵}.}
    \item \zh{我的观点跟老师的一样正确。} (Wǒ de guāndiǎn gēn lǎoshī de yíyàng zhèngquè.)
    \item \zh{北京的冬天与哈尔滨的不同。} (Běijīng de dōngtiān yǔ Hā'ěrbīn de bùtóng.) \emph{Note : \zh{不同} s'emploie seul, on supprime \zh{冷}.}
    \item \zh{这道题跟那道题不一样难。} (Zhège tí gēn nàge tí bù yíyàng nán.)
\end{enumerate}
\end{corrige}

\begin{corrige}[Exercice 3 : Expression écrite guidée]
\textbf{A. Version :}
\begin{enumerate}
    \item Le niveau du football camerounais est aussi élevé que celui des grandes équipes africaines.
    \item La circulation de Yaoundé est différente de celle d'il y a dix ans, maintenant il y a des bus.
    \item Apprendre le chinois n'est pas aussi difficile [qu'on le croit], pourvu que la méthode soit aussi bonne que celle des autres.
\end{enumerate}

\textbf{B. Production écrite (exemple de réponse attendue) :}
\zh{我的家乡是布埃亚。布埃亚跟雅温得不一样大，人口比较少。可是布埃亚的气候和昆明的一样舒服，不太热也不太冷。这里的交通与大城市不同，不太堵车。吃饭方面，布埃亚的主食跟杜阿拉的一样，都是木薯粉和植物香蕉。欢迎你来布埃亚参观！} \\
(Wǒ de jiāxiāng shì Bù'āiyà. Bù'āiyà gēn Yǎwēndé bù yíyàng dà, rénkǒu bǐjiào shǎo. Kěshì Bù'āiyà de qìhòu hé Kūnmíng de yíyàng shūfu, bù tài rè yě bù tài lěng. Zhèlǐ de jiāotōng yǔ dà chéngshì bùtóng, bù tài dǔchē. Chīfàn fāngmiàn, Bù'āiyà de zhǔshí gēn Duālā de yíyàng, dōu shì mùshǔ fěn hé zhíwù xiāngjiāo. Huānyíng nǐ lái Bù'āiyà cānguān!)
\end{corrige}

\begin{aretenir}[Récapitulatif final pour le Bac / HSK 4]
\begin{itemize}
    \item \textbf{Ordre immuable :} \zh{Sujet + 跟/和/与 + Objet + 一样/相同/不一样/不同 + (Adj/Verbe)}.
    \item \textbf{Registres :} \zh{跟/一样/不一样} (Oral) | \zh{和} (Neutre) | \zh{与/相同/不同} (Écrit/Formel).
    \item \textbf{Avec adjectif/verbe :} \zh{一样}, \zh{不一样} (quasi obligatoire). \zh{相同}, \zh{不同} (souvent seuls).
    \item \textbf{Ton de \zh{一} :} \zh{一样} (yíyàng), \zh{不一样} (bù yí yàng).
    \item \textbf{Pièges :} Pas de préposition $\rightarrow$ phrase fausse. Calque français \zh{*一样 Adj 跟...} $\rightarrow$ phrase fausse.
\end{itemize}
\end{aretenir}

\coursec{Schéma de la structure d'égalité : A 跟/和/与 B 一样/相同}

Dans la section précédente, nous avons observé des phrases authentiques comme \zh{我跟他一样高。} (\emph{Wǒ gēn tā yíyàng gāo.} « Je suis aussi grand que lui. ») et \zh{我们的想法相同。} (\emph{Wǒmen de xiǎngfa xiāngtóng.} « Nos idées sont identiques. »). Nous avons remarqué que le chinois ne dit pas « aussi... que » avec deux mots séparés comme le français, mais avec une construction en blocs qui encadre les deux termes comparés. Il est temps maintenant de dégager la règle générale, de la formaliser et de l'entraîner. C'est l'objet de cette section.

\courssub{Le schéma canonique : A + 跟/和/与 + B + 一样 (+ adjectif)}

Observons d'abord trois exemples gradués, du plus simple au plus complet :

\begin{exemplebox}[Trois exemples gradués à observer]
\zh{我跟他一样。} \emph{Wǒ gēn tā yíyàng.} « Je suis comme lui. / Nous sommes pareils, lui et moi. »\\[4pt]
\zh{我跟他一样高。} \emph{Wǒ gēn tā yíyàng gāo.} « Je suis aussi grand que lui. »\\[4pt]
\zh{我跟他一样喜欢汉语。} \emph{Wǒ gēn tā yíyàng xǐhuan Hànyǔ.} « J'aime le chinois autant que lui. »
\end{exemplebox}

Dans les trois phrases, on retrouve exactement les mêmes pièces, dans le même ordre : d'abord le premier terme comparé (\zh{我}), puis le mot-outil (\zh{跟}), puis le second terme (\zh{他}), puis le mot \cle{一样} « pareil, de la même façon », et enfin, éventuellement, un adjectif (\zh{高}) ou un verbe (\zh{喜欢汉语}) qui précise \emph{sur quel point} les deux termes sont égaux.

\begin{propriete}[Règle de la structure d'égalité avec 一样]
Pour dire « A est aussi... que B » ou « A est comme B », le chinois utilise la structure :

\textbf{A + 跟/和/与 + B + 一样 (+ adjectif / verbe)}

\begin{itemize}
\item \cle{跟}, \cle{和} ou \cle{与} jouent le rôle du « que / comme » français de comparaison ; ils se placent \textbf{entre A et B}.
\item Le mot \cle{一样} (yíyàng, « pareil, identique ») vient \textbf{après B}, jamais entre A et B.
\item L'adjectif (ou le verbe) qui précise le point de comparaison se place \textbf{après 一样} : \zh{一样高} « aussi grand », \zh{一样贵} « aussi cher », \zh{一样喜欢} « aimer autant ».
\item Sans adjectif, \zh{A 跟 B 一样} signifie simplement « A est comme B, A et B sont pareils ».
\end{itemize}
\end{propriete}

La figure suivante visualise cet enchaînement de blocs, dans l'ordre fixe et immuable de la phrase chinoise.

\begin{popfigure}
\begin{center}
\begin{tikzpicture}[
  bloc/.style={rectangle, rounded corners=3pt, draw=popblue, fill=popblueL, minimum height=0.9cm, align=center, font=\small},
  blocb/.style={rectangle, rounded corners=3pt, draw=poporange, fill=poporangeL, minimum height=0.9cm, align=center, font=\small},
  blocc/.style={rectangle, rounded corners=3pt, draw=popgreen, fill=popgreenL, minimum height=0.9cm, align=center, font=\small},
  fl/.style={-{Stealth[length=2.5mm]}, thick, popdark}
]
\node[bloc, align=center] (a) at (0,0){A\\\zh{我}};
\node[blocb, right=0.7cm of a, align=center] (gen){跟 / 和 / 与\\« que, comme »};
\node[bloc, right=0.7cm of gen, align=center] (b){B\\\zh{他}};
\node[blocc, right=0.7cm of b, align=center] (yy){一样\\« pareil »};
\node[blocc, right=0.7cm of yy, align=center] (adj){(adjectif / verbe)\\\zh{高} / \zh{喜欢汉语}};
\draw[fl] (a) -- (gen);
\draw[fl] (gen) -- (b);
\draw[fl] (b) -- (yy);
\draw[fl] (yy) -- (adj);
\node[below=0.5cm of gen, align=center, font=\itshape\small, text=popmuted] {« A est aussi ... que B »};
\end{tikzpicture}
\end{center}
\legende{Schéma de la structure d'égalité : les blocs s'enchaînent dans un ordre fixe.}
\end{popfigure}

\begin{aretenir}[L'ordre des mots, clé de toute la structure]
En français, « aussi » se place \emph{avant} l'adjectif et « que » \emph{avant} B : « je suis \textbf{aussi} grand \textbf{que} lui ». En chinois, c'est l'inverse logique : le mot-outil \zh{跟} précède B, et \zh{一样} suit B. On peut mémoriser la formule : \textbf{le « que » est avant B, le « pareil » est après B}. Les deux termes A et B sont comme encadrés : A \zh{跟} B \zh{一样}.
\end{aretenir}

\courssub{Préciser le point de comparaison : 一样 + adjectif}

Le plus souvent, on ne dit pas seulement que deux choses « sont pareilles » : on précise \emph{en quoi} elles se ressemblent. On ajoute alors un adjectif après \zh{一样}. Tous les adjectifs qualificatifs conviennent : \zh{高} « grand », \zh{大} « grand (taille) », \zh{漂亮} « beau », \zh{贵} « cher », \zh{热} « chaud », \zh{远} « loin », \zh{忙} « occupé », \zh{好吃} « bon (au goût) ».

\begin{exemplebox}[一样 + adjectif : exemples variés]
\zh{Yaoundé 跟 Douala 一样热吗？} \emph{Yǎwēndé gēn Duālā yíyàng rè ma?} « Yaoundé est-elle aussi chaude que Douala ? »\\[4pt]
\zh{这本书和那本书一样贵。} \emph{Zhè běn shū hé nà běn shū yíyàng guì.} « Ce livre est aussi cher que celui-là. »\\[4pt]
\zh{我妹妹跟我一样大。} \emph{Wǒ mèimei gēn wǒ yíyàng dà.} « Ma petite sœur a le même âge que moi (est aussi grande que moi). »\\[4pt]
\zh{今天的作业跟昨天的一样多。} \emph{Jīntiān de zuòyè gēn zuótiān de yíyàng duō.} « Les devoirs d'aujourd'hui sont aussi nombreux que ceux d'hier. »\\[4pt]
\zh{她跟她妈妈一样漂亮。} \emph{Tā gēn tā māma yíyàng piàoliang.} « Elle est aussi belle que sa mère. »
\end{exemplebox}

Remarquez la question de la première phrase : on ajoute simplement la particule interrogative \zh{吗} à la fin, sans rien changer à l'ordre des mots. La structure d'égalité reste intacte à l'interrogatif.

\courssub{Comparer des actions : 一样 + verbe}

La structure ne sert pas qu'avec les adjectifs. Pour dire que deux personnes font une action « de la même façon, autant l'une que l'autre », on place le groupe verbal après \zh{一样} :

\textbf{A + 跟/和/与 + B + 一样 + verbe (+ complément)}

\begin{exemplebox}[Comparaison d'actions]
\zh{他跟我一样喜欢足球。} \emph{Tā gēn wǒ yíyàng xǐhuan zúqiú.} « Il aime le football comme moi / autant que moi. »\\[4pt]
\zh{我弟弟跟我一样想学汉语。} \emph{Wǒ dìdi gēn wǒ yíyàng xiǎng xué Hànyǔ.} « Mon petit frère veut apprendre le chinois comme moi. »\\[4pt]
\zh{喀麦隆学生跟中国学生一样努力。} \emph{Kāmàilóng xuésheng gēn Zhōngguó xuésheng yíyàng nǔlì.} « Les élèves camerounais travaillent aussi dur que les élèves chinois. »
\end{exemplebox}

Notez que dans la dernière phrase, \zh{努力} « faire des efforts » fonctionne comme un adjectif-verbe : la grammaire est exactement la même. Le chinois ne distingue pas rigoureusement adjectif et verbe dans cette position ; ce qui compte, c'est que le mot placé après \zh{一样} exprime la qualité ou l'action comparée.

\courssub{La variante avec 相同 : « identique »}

\begin{definition}[相同 (xiāngtóng)]
\cle{相同} est un adjectif qui signifie « identique, pareil, semblable ». Contrairement à \zh{一样}, il s'emploie \textbf{seul, sans adjectif après lui} : on ne dit pas que deux choses sont « identiquement grandes », on dit qu'elles sont identiques, point.

\textbf{A + 跟/和/与 + B + 相同} ou bien \textbf{A + 和 + B + 相同} : « A et B sont identiques ».
\end{definition}

\begin{exemplebox}[相同 en contexte]
\zh{我们的想法相同。} \emph{Wǒmen de xiǎngfa xiāngtóng.} « Nos idées sont identiques. »\\[4pt]
\zh{这个字的发音跟那个字相同。} \emph{Zhège zì de fāyīn gēn nàge zì xiāngtóng.} « La prononciation de ce caractère est identique à celle de celui-là. »\\[4pt]
\zh{他们的答案与我们的相同。} \emph{Tāmen de dá'àn yǔ wǒmen de xiāngtóng.} « Leurs réponses sont identiques aux nôtres. »
\end{exemplebox}

Différence d'emploi à retenir : \zh{一样} est le mot de la conversation quotidienne et accepte un adjectif ou un verbe après lui (\zh{一样高}, \zh{一样喜欢}) ; \zh{相同} est un peu plus soutenu, plus « constat d'identité », et se termine la phrase (\zh{相同。}). À l'examen, si un adjectif suit, choisissez \zh{一样} ; si la phrase s'arrête sur « identique », les deux sont possibles, \zh{相同} étant alors très naturel.

\courssub{Intensifier : 完全一样 « tout à fait pareil »}

On peut renforcer la comparaison en plaçant l'adverbe \cle{完全} (wánquán, « complètement ») devant \zh{一样} :

\textbf{A + 跟/和/与 + B + 完全 + 一样 (+ adjectif)} : « A est tout à fait comme B ».

\begin{exemplebox}[Intensification avec 完全]
\zh{我的看法跟你的完全一样。} \emph{Wǒ de kànfǎ gēn nǐ de wánquán yíyàng.} « Mon point de vue est tout à fait identique au tien. »\\[4pt]
\zh{这两个词的意思完全一样。} \emph{Zhè liǎng ge cí de yìsi wánquán yíyàng.} « Le sens de ces deux mots est exactement le même. »
\end{exemplebox}

\courssub{Choisir le mot-outil : 跟, 和 ou 与 ?}

Les trois mots-outils \zh{跟}, \zh{和} et \zh{与} construisent \emph{exactement la même grammaire}. Seul le \textbf{registre de langue} change :

\begin{itemize}
\item \cle{跟} (gēn) : le plus courant \textbf{à l'oral}, dans la conversation de tous les jours ;
\item \cle{和} (hé) : \textbf{neutre}, convient à l'oral comme à l'écrit ; c'est le choix de sécurité ;
\item \cle{与} (yǔ) : \textbf{écrit et formel} (articles, dissertations, documents officiels).
\end{itemize}

\begin{popfigure}
\begin{center}
\begin{tikzpicture}
\draw[-{Stealth[length=3mm]}, very thick, popline] (0,0) -- (12,0);
\node[circle, draw=popgreen, fill=popgreenL, minimum size=1.1cm, align=center, font=\small] at (1.5,0) {\zh{跟}\\gēn};
\node[circle, draw=popblue, fill=popblueL, minimum size=1.1cm, align=center, font=\small] at (6,0) {\zh{和}\\hé};
\node[circle, draw=poppurple, fill=poppurpleL, minimum size=1.1cm, align=center, font=\small] at (10.5,0) {\zh{与}\\yǔ};
\node[below=0.75cm of {1.5,0}, align=center, font=\small] at (1.5,-0.8) {oral courant};
\node[below=0.75cm of {6,0}, align=center, font=\small] at (6,-0.8) {neutre};
\node[below=0.75cm of {10.5,0}, align=center, font=\small] at (10.5,-0.8) {écrit, formel};
\node[above=0.2cm, align=center, font=\itshape\small, text=popmuted] at (6,0.9) {même grammaire, registre différent};
\end{tikzpicture}
\end{center}
\legende{Frise de registre : de la conversation familière au style écrit formel.}
\end{popfigure}

Ainsi, les trois phrases suivantes ont \emph{exactement le même sens}, seul le niveau de langue diffère :

\begin{exemplebox}[Trois registres, un seul sens]
\zh{他跟我一样高。} \emph{Tā gēn wǒ yíyàng gāo.} (oral) « Il est aussi grand que moi. »\\[4pt]
\zh{他和我一样高。} \emph{Tā hé wǒ yíyàng gāo.} (neutre) « Il est aussi grand que moi. »\\[4pt]
\zh{他与我一样高。} \emph{Tā yǔ wǒ yíyàng gāo.} (formel) « Il est aussi grand que moi. »
\end{exemplebox}

\courssub{Tableau récapitulatif et comparaison avec le français}

\begin{center}
\begin{tabularx}{\linewidth}{|X|X|X|}
\hline
\rowcolor{popblueL} Structure & Exemple (caractères + pinyin) & Traduction \\
\hline
A + 跟/和/与 + B + 一样 & \zh{我跟他一样。} Wǒ gēn tā yíyàng. & Je suis comme lui. \\
\hline
A + 跟/和/与 + B + 一样 + Adj. & \zh{我跟他一样高。} Wǒ gēn tā yíyàng gāo. & Je suis aussi grand que lui. \\
\hline
A + 跟/和/与 + B + 一样 + V. & \zh{他跟我一样喜欢足球。} Tā gēn wǒ yíyàng xǐhuan zúqiú. & Il aime le football comme moi. \\
\hline
A + 跟/和/与 + B + 完全一样 & \zh{我的想法跟你的完全一样。} Wǒ de xiǎngfa gēn nǐ de wánquán yíyàng. & Mon idée est tout à fait la même que la tienne. \\
\hline
A + 跟/和/与 + B + 相同 & \zh{我们的想法相同。} Wǒmen de xiǎngfa xiāngtóng. & Nos idées sont identiques. \\
\hline
\end{tabularx}
\end{center}

\begin{center}
\begin{tabularx}{\linewidth}{|X|X|X|X|}
\hline
\rowcolor{popblueL} Caractères & Pinyin & Nature & Traduction \\
\hline
一样 & yíyàng & adverbe / adjectif & pareil, de la même façon \\
\hline
相同 & xiāngtóng & adjectif & identique, semblable \\
\hline
跟 & gēn & préposition / conjonction & avec, que (comparaison), oral \\
\hline
和 & hé & conjonction / préposition & et, avec, que (neutre) \\
\hline
与 & yǔ & préposition / conjonction & avec, et (écrit, formel) \\
\hline
完全 & wánquán & adverbe & complètement, tout à fait \\
\hline
\end{tabularx}
\end{center}

Comparaison des deux langues : le français dit « A est \textbf{aussi} grand \textbf{que} B » avec deux mots répartis autour de l'adjectif ; le chinois dit « A \zh{跟} B \zh{一样} grand » avec deux mots répartis autour de B. L'adjectif chinois se place tout à la fin, sans mot de liaison devant lui. Autre différence : le français exige l'accord et parfois un pronom disjoint (« que \emph{lui}, que \emph{moi} ») ; le chinois utilise simplement le pronom ou le nom après \zh{跟}, sans aucune modification.

\begin{attention}[Piège n°1 : ne jamais traduire « aussi » par 也]
L'erreur la plus fréquente des francophones est de traduire « aussi... que » mot à mot avec \zh{也} : \zh{我也跟他一样高} pour « je suis aussi grand que lui ». Or \cle{也} (yě) signifie « également, aussi » au sens d'addition (« moi aussi »), et \textbf{ne construit pas la comparaison}. La comparaison d'égalité se construit uniquement avec \zh{跟/和/与 ... 一样}. La phrase correcte est : \zh{我跟他一样高。} (\emph{Wǒ gēn tā yíyàng gāo.}) — sans \zh{也}. Si vous ajoutez \zh{也}, la phrase signifie « moi aussi, je suis aussi grand que lui », ce qui suppose qu'une troisième personne vient d'être mentionnée.
\end{attention}

\begin{attention}[Piège n°2 : la place de 一样]
Ne placez jamais \zh{一样} entre A et B : \zh{我一样跟他高} est incorrect. Le mot \zh{一样} vient toujours \textbf{après B}. De même, l'adjectif ne se met jamais avant \zh{一样} : on dit \zh{一样高}, jamais \zh{高一样}. Mémorisez le bloc insécable \zh{一样} + adjectif.
\end{attention}

\begin{exoresolu}[Construire une comparaison d'égalité]
Énoncé : traduisez en chinois la phrase « Ma chambre est aussi grande que celle de mon frère. » (vocabulaire : \zh{房间} fángjiān « chambre » ; \zh{哥哥} gēge « frère aîné » ; \zh{大} dà « grand »).

\tcblower

Solution détaillée, étape par étape :

\begin{enumerate}
\item Identifier les deux termes comparés : A = « ma chambre » \zh{我的房间} ; B = « celle de mon frère » \zh{我哥哥的} (le pronom possessif \zh{的} suffit, comme « celle de » en français).
\item Identifier le point de comparaison : « grand » \zh{大}.
\item Choisir le mot-outil : contexte courant, on prend \zh{跟} (ou \zh{和}).
\item Assembler dans l'ordre fixe : A + \zh{跟} + B + \zh{一样} + adjectif.
\end{enumerate}

Phrase finale : \zh{我的房间跟我哥哥的一样大。}

\emph{Wǒ de fángjiān gēn wǒ gēge de yíyàng dà.}

Vérification : \zh{一样} est bien après B (\zh{我哥哥的}), l'adjectif \zh{大} est bien après \zh{一样}, et aucun \zh{也} ne s'est glissé dans la phrase. La structure est correcte.
\end{exoresolu}

\begin{exercice}[À vous de jouer]
Traduisez en chinois, en utilisant la structure d'égalité : 1. « Douala est aussi grande que Yaoundé ? » (forme interrogative). 2. « Mon père aime le football autant que moi. » 3. « Leurs réponses sont identiques. » (employez \zh{相同}). 4. « Ce marché est tout à fait comme celui de mon village. » (employez \zh{完全}).
\end{exercice}

\begin{corrige}[Exercice « À vous de jouer »]
1. \zh{Douala 跟 Yaoundé 一样大吗？} \emph{Duālā gēn Yǎwēndé yíyàng dà ma?} — la particule \zh{吗} suffit pour la question.\\
2. \zh{我爸爸跟我一样喜欢足球。} \emph{Wǒ bàba gēn wǒ yíyàng xǐhuan zúqiú.} — verbe \zh{喜欢} après \zh{一样}.\\
3. \zh{他们的答案相同。} \emph{Tāmen de dá'àn xiāngtóng.} — \zh{相同} seul, sans adjectif.\\
4. \zh{这个市场跟我家乡的完全一样。} \emph{Zhège shìchǎng gēn wǒ jiāxiāng de wánquán yíyàng.} — \zh{完全} devant \zh{一样}.
\end{corrige}

La structure d'égalité est désormais posée : un ordre de blocs fixe, trois mots-outils interchangeables selon le registre, et deux mots de comparaison, \zh{一样} (qui accepte un adjectif ou un verbe) et \zh{相同} (qui conclut la phrase). La section suivante exploitera ce même schéma pour exprimer, au contraire, la \emph{différence}, grâce à la négation \zh{不一样} et à l'adjectif \zh{不同}.

\coursec{Le comparatif : égalité et différence}

\courssub{Schéma de l'égalité : A 跟/和/与 B 一样/相同}

\begin{textebox}[Dialogue d'observation : deux camarades de classe]
\zh{— 你的中文跟我的中文一样好吗？} \emph{Nǐ de Zhōngwén gēn wǒ de Zhōngwén yíyàng hǎo ma?}\\
\zh{— 是的，我们的水平相同。} \emph{Shì de, wǒmen de shuǐpíng xiāngtóng.}\\
\zh{— 那我们一起复习吧，方法也一样。} \emph{Nà wǒmen yìqǐ fùxí ba, fāngfǎ yě yíyàng.}
\par\medskip
\emph{— Ton chinois est-il aussi bon que le mien ?}\\
\emph{— Oui, notre niveau est identique.}\\
\emph{— Alors révisons ensemble, la méthode est aussi la même.}
\end{textebox}

\begin{propriete}[Schéma de l'égalité]
\textbf{Forme 1 (courante, avec adjectif) :}
\begin{center}
A + 跟/和/与 + B + \textbf{一样} + adjectif
\end{center}
\textbf{Forme 2 (formelle, sans adjectif) :}
\begin{center}
A + 跟/和/与 + B + \textbf{相同}
\end{center}
\zh{跟} (gēn), \zh{和} (hé), \zh{与} (yǔ) sont des prépositions signifiant « avec » ; \zh{与} est plus soutenu (écrit). \zh{一样} (yíyàng) « pareil, identique » s'emploie à l'oral et à l'écrit ; \zh{相同} (xiāngtóng) « identique » est formel. L'adjectif est facultatif avec \zh{一样}, impossible avec \zh{相同}.
\end{propriete}

\begin{popfigure}
\begin{center}
\begin{tikzpicture}[
  bloc/.style={rectangle, rounded corners=3pt, draw=popblue, fill=popblueL, minimum height=0.9cm, align=center, font=\small},
  blocopt/.style={rectangle, rounded corners=3pt, draw=popgreen, fill=popgreenL, minimum height=0.9cm, align=center, font=\small},
  fleche/.style={-{Stealth[length=2.5mm]}, thick, popdark}
]
\node[bloc, align=center] (a){A\\我};
\node[bloc, right=0.55cm of a, align=center] (gen){跟 / 和 / 与\\gēn / hé / yǔ};
\node[bloc, right=0.55cm of gen, align=center] (b){B\\他};
\node[bloc, right=0.55cm of b, align=center] (same){\textbf{一样}\\yíyàng};
\node[blocopt, right=0.55cm of same, align=center] (adj){(adjectif)\\高 gāo};
\draw[fleche] (a) -- (gen);
\draw[fleche] (gen) -- (b);
\draw[fleche] (b) -- (same);
\draw[fleche] (same) -- (adj);
\node[below=0.35cm of same, align=center, font=\footnotesize\itshape, text=popdark] {我跟他一样高。\\ « Je suis aussi grand que lui. »};
\end{tikzpicture}
\end{center}
\legende{Schéma de la structure d'égalité : A + 跟/和/与 + B + 一样 + (adjectif)}
\end{popfigure}

\begin{exemplebox}[Exemples d'égalité]
\zh{雅温得跟杜阿拉一样热。} \emph{Yǎwēndé gēn Dùālā yíyàng rè.} « Yaoundé est aussi chaud que Douala. »\\
\zh{这道菜和那道菜一样好吃。} \emph{Zhè dào cài hé nà dào cài yíyàng hǎochī.} « Ce plat est aussi bon que celui-là. »\\
\zh{中非两国的友谊与以前相同。} \emph{Zhōng-Fēi liǎngguó de yǒuyì yǔ yǐqián xiāngtóng.} « L'amitié entre la Chine et le Cameroun est identique à avant. »\\
\zh{我们的教科书相同。} \emph{Wǒmen de jiàokēshū xiāngtóng.} « Nos manuels sont identiques. »
\end{exemplebox}

\begin{aretenir}[Nuance : 一样 ou 相同 ?]
\begin{itemize}
\item \zh{一样} : registre courant, admet un adjectif (\zh{一样大}), peut être modifié (\zh{完全一样} « exactement pareil »).
\item \zh{相同} : registre formel/écrit, \textbf{ne prend pas d'adjectif} derrière (✗ \zh{相同大}), s'emploie souvent pour des concepts abstraits (niveau, nature, structure, culture).
\end{itemize}
\end{aretenir}

\courssub{Forme interrogative de l'égalité}
Pour demander si deux choses sont pareilles :
\begin{itemize}
\item \zh{A 跟 B 一样吗？} \emph{A gēn B yíyàng ma?} (neutre)
\item \zh{A 跟 B 是不是一样？} \emph{A gēn B shì bu shì yíyàng?} (attente de confirmation)
\end{itemize}
Réponses : \zh{一样。} / \zh{不一样。} ; \zh{一样高。} / \zh{不一样高。}

\courssub{Schéma de la différence : A 跟/和/与 B 不一样/不同}

\begin{textebox}[Transition : constater des différences]
\zh{— 喀麦隆北方跟南方的气候一样吗？} \emph{Kāmàilóng běifāng gēn nánfāng de qìhòu yíyàng ma?}\\
\zh{— 不一样。北方与南方的气候不同。} \emph{Bù yíyàng. Běifāng yǔ nánfāng de qìhòu bùtóng.}\\
\zh{— 那杜阿拉和雅温得一样大吗？} \emph{Nà Dùālā hé Yǎwēndé yíyàng dà ma?}\\
\zh{— 不一样大，杜阿拉比雅温得大。} \emph{Bù yíyàng dà, Dùālā bǐ Yǎwēndé dà.}
\par\medskip
\emph{— Le climat du Nord et du Sud du Cameroun est-il pareil ?}\\
\emph{— Pas pareil. Le climat du Nord est différent de celui du Sud.}\\
\emph{— Et Douala et Yaoundé sont-elles aussi grandes ?}\\
\emph{— Pas aussi grandes, Douala est plus grande que Yaoundé.}
\end{textebox}

\begin{propriete}[Schéma de la différence]
\textbf{Forme 1 (courante, degré, avec adjectif possible) :}
\begin{center}
A + 跟/和/与 + B + \textbf{不} + \textbf{一样} + (adjectif)
\end{center}
\textbf{Forme 2 (formelle, catégorique, sans adjectif) :}
\begin{center}
A + 跟/和/与 + B + \textbf{不同} \quad | \quad \textbf{不同} + \textbf{的} + nom
\end{center}
La négation \zh{不} (bù) se place \textbf{uniquement devant 一样}, jamais devant \zh{跟/和/与}. \zh{不同} (bùtóng) est un adjectif verbal signifiant « être différent », plus formel que \zh{不一样}.
\end{propriete}

\begin{popfigure}
\begin{center}
\begin{tikzpicture}[
  bloc/.style={rectangle, rounded corners=3pt, draw=popblue, fill=popblueL, minimum height=0.9cm, align=center, font=\small},
  blocneg/.style={rectangle, rounded corners=3pt, draw=poporange, fill=poporangeL, minimum height=0.9cm, align=center, font=\small},
  blocopt/.style={rectangle, rounded corners=3pt, draw=popgreen, fill=popgreenL, minimum height=0.9cm, align=center, font=\small},
  fleche/.style={-{Stealth[length=2.5mm]}, thick, popdark}
]
\node[bloc, align=center] (a){A\\今天};
\node[bloc, right=0.55cm of a, align=center] (gen){跟 / 和 / 与\\gēn / hé / yǔ};
\node[bloc, right=0.55cm of gen, align=center] (b){B\\昨天};
\node[blocneg, right=0.55cm of b, align=center] (neg){\textbf{不} + 一样\\ou \textbf{不同}};
\node[blocopt, right=0.55cm of neg, align=center] (adj){(adjectif)\\冷 lěng};
\draw[fleche] (a) -- (gen);
\draw[fleche] (gen) -- (b);
\draw[fleche] (b) -- (neg);
\draw[fleche] (neg) -- (adj);
\node[below=0.35cm of neg, align=center, font=\footnotesize\itshape, text=popdark] {今天跟昨天不一样冷。\\ « Aujourd'hui il n'est pas aussi froid qu'hier. »};
\end{tikzpicture}
\end{center}
\legende{Schéma de la structure de différence : A + 跟/和/与 + B + 不一样 / 不同 + (adjectif)}
\end{popfigure}

\begin{propriete}[Règle d'or : la place de la négation]
Le chinois nie l'égalité (\zh{一样}), pas la préposition (\zh{跟}).
\begin{center}
\begin{tabularx}{\linewidth}{|X|X|}
\hline
\rowcolor{poporangeL} \textbf{Correct} & \textbf{Incorrect (faute grave)} \\
\hline
\zh{我跟他不一样高。} \emph{Wǒ gēn tā bù yíyàng gāo.} & \zh{我不跟他一样高。} \emph{Wǒ bù gēn tā yíyàng gāo.} \\
\zh{法语跟英语不一样难。} \emph{Fǎyǔ gēn Yīngyǔ bù yíyàng nán.} & \zh{法语不跟英语一样难。} \\
\hline
\end{tabularx}
\end{center}
\zh{跟 B} forme un bloc inséparable « avec B ». \zh{不一样} forme le bloc « pas pareil ». On ne casse pas le premier bloc.
\end{propriete}

\begin{exemplebox}[不一样 + adjectif : « pas aussi … que »]
\zh{今天跟昨天不一样冷。} \emph{Jīntiān gēn zuótiān bù yíyàng lěng.} « Aujourd'hui il n'est pas aussi froid qu'hier. »\\
\zh{英语跟法语不一样难。} \emph{Yīngyǔ gēn Fǎyǔ bù yíyàng nán.} « L'anglais n'est pas aussi difficile que le français. »\\
\zh{这本书跟那本书不一样贵。} \emph{Zhè běn shū gēn nà běn shū bù yíyàng guì.} « Ce livre n'est pas aussi cher que celui-là. »\\
\zh{汉语跟英语不一样容易。} \emph{Hànyǔ gēn Yīngyǔ bù yíyàng róngyì.} « Le chinois n'est pas aussi facile que l'anglais. »
\end{exemplebox}

\begin{attention}[Nuance de sens et ordre de A et B]
\zh{A 跟 B 不一样 + adjectif} signifie « A n'est pas aussi [adjectif] que B » $\rightarrow$ A possède \textbf{moins} la qualité.
\zh{英语跟法语不一样难} = L'anglais < Le français (difficulté).
\zh{法语跟英语不一样难} = Le français < L'anglais (difficulté).
L'ordre de A et B inverse le sens, exactement comme en français : « Paul n'est pas aussi grand que Pierre » $\neq$ « Pierre n'est pas aussi grand que Paul ».
\end{attention}

\begin{definition}[不同 bùtóng : différence catégorique, registre formel]
\zh{不同} s'emploie de deux façons :
\begin{enumerate}
\item \zh{A 跟/和/与 B 不同} : « A est différent de B » (identité, nature, catégorie).
\item \zh{不同 + 的 + Nom} : « des [Nom] différents ».
\end{enumerate}
\emph{Attention :} on ne dit jamais ✗ \zh{不同 + adjectif} (✗ \zh{不同高}, ✗ \zh{不同难}). Pour exprimer « pas aussi [adj] que », on utilise obligatoirement \zh{不一样 + adjectif}.
\end{definition}

\begin{exemplebox}[不同 à l'écrit et dans les noms]
\zh{北方与南方的气候不同。} \emph{Běifāng yǔ nánfāng de qìhòu bùtóng.} « Le climat du Nord diffère de celui du Sud. »\\
\zh{中国和喀麦隆有不同的文化。} \emph{Zhōngguó hé Kāmàilóng yǒu bùtóng de wénhuà.} « La Chine et le Cameroun ont des cultures différentes. »\\
\zh{每个学生的学习方法不同。} \emph{Měi ge xuésheng de xuéxí fāngfǎ bùtóng.} « La méthode d'apprentissage de chaque élève est différente. »\\
\zh{我们对这个问题的看法不同。} \emph{Wǒmen duì zhè ge wèntí de kànfǎ bùtóng.} « Notre avis sur cette question diverge. »
\end{exemplebox}

\begin{aretenir}[Choisir entre 不一样 et 不同]
\begin{center}
\begin{tabularx}{\linewidth}{|X|X|X|}
\hline
\rowcolor{popblueL} \textbf{Critère} & \textbf{不一样 (bù yíyàng)} & \textbf{不同 (bùtóng)} \\
\hline
Sens & « Pas pareil » (degré, intensité) & « Différent » (nature, catégorie) \\
\hline
Registre & Oral, courant, universel & Écrit, formel, soutenu \\
\hline
+ Adjectif & \textbf{OUI} (\zh{不一样大}) & \textbf{NON} (✗ \zh{不同大}) \\
\hline
+ 的 + Nom & Rare (\zh{不一样的颜色} possible) & \textbf{STANDARD} (\zh{不同的颜色}) \\
\hline
Exemple type & \zh{这两件衣服不一样大。} & \zh{这两件衣服颜色不同。} \\
\hline
\end{tabularx}
\end{center}
\end{aretenir}

\courssub{Forme interrogative de la différence}
Mêmes structures qu'à l'affirmatif, avec \zh{吗} ou \zh{是不是} :
\begin{itemize}
\item \zh{你的学校跟我的学校不一样吗？} \emph{Nǐ de xuéxiào gēn wǒ de xuéxiào bù yíyàng ma?}
\item \zh{这两个汉字是不是不同？} \emph{Zhè liǎng ge Hànzì shì bu shì bùtóng?}
\end{itemize}
Réponses : \zh{不一样。} / \zh{一样。} ; \zh{不同。} / \zh{相同。}

\courssub{Tableau récapitulatif des structures (Égalité vs Différence)}

\begin{center}
\begin{tabularx}{\linewidth}{|X|X|X|}
\hline
\rowcolor{popblueL} \textbf{Structure} & \textbf{Exemple (caractères + pinyin)} & \textbf{Traduction} \\
\hline
A 跟 B \textbf{一样} + adj. & \zh{我跟他一样高。} \emph{Wǒ gēn tā yíyàng gāo.} & Je suis aussi grand que lui. \\
\hline
A 跟 B \textbf{相同} & \zh{我们的观点相同。} \emph{Wǒmen de guāndiǎn xiāngtóng.} & Nos points de vue sont identiques. \\
\hline
A 跟 B \textbf{不一样} + adj. & \zh{我跟他不一样高。} \emph{Wǒ gēn tā bù yíyàng gāo.} & Je ne suis pas aussi grand que lui. \\
\hline
A 跟 B \textbf{不同} & \zh{我们的观点不同。} \emph{Wǒmen de guāndiǎn bùtóng.} & Nos points de vue sont différents. \\
\hline
\textbf{不同} + 的 + Nom & \zh{不同的文化} \emph{bùtóng de wénhuà} & Des cultures différentes \\
\hline
\end{tabularx}
\end{center}

\courssub{Vocabulaire clé de la leçon}

\begin{center}
\begin{tabularx}{\linewidth}{|X|X|X|X|}
\hline
\rowcolor{popblueL} \textbf{Caractères} & \textbf{Pinyin} & \textbf{Nature} & \textbf{Traduction} \\
\hline
跟 & gēn & Préposition & avec (oral, courant) \\
\hline
和 & hé & Préposition & avec (neutre) \\
\hline
与 & yǔ & Préposition & avec (formel, écrit) \\
\hline
一样 & yíyàng & Adj. verbal & pareil, identique, semblable \\
\hline
相同 & xiāngtóng & Adj. verbal & identique (formel) \\
\hline
不一样 & bù yíyàng & Locution & pas pareil, pas aussi … que \\
\hline
不同 & bùtóng & Adj. verbal & différent (formel) \\
\hline
程度 & chéngdù & Nom & degré, niveau \\
\hline
类别 & lèibié & Nom & catégorie, type \\
\hline
比较 & bǐjiào & Verbe & comparer \\
\hline
\end{tabularx}
\end{center}

\courssub{Méthode de transformation et exercice résolu}

\begin{methode}[Transformer égalité $\leftrightarrow$ différence]
\begin{enumerate}
\item Identifier le marqueur d'égalité : \zh{一样} (ou \zh{相同}).
\item \textbf{Affirmatif $\rightarrow$ Négatif} : insérer \zh{不} \textbf{immédiatement devant} \zh{ayama}. (Si \zh{相同} $\rightarrow$ remplacer par \zh{不同}).
\item \textbf{Négatif $\rightarrow$ Affirmatif} : supprimer \zh{不} devant \zh{ayama} (ou remplacer \zh{不同} par \zh{相同} / \zh{一样}).
\item \textbf{Ne jamais déplacer} \zh{不} devant \zh{跟/和/与}.
\item Vérifier la cohérence sémantique : \zh{不一样 + adj} = « moins [adj] que ».
\end{enumerate}
\end{methode}

\begin{exoresolu}[Transformation guidée]
Énoncé : Transformez en phrase de différence.\\
1. \zh{杜阿拉的天气跟雅温得一样热。} \emph{Dùālā de tiānqì gēn Yǎwēndé yíyàng rè.}\\
2. \zh{中非两国的友谊与以前相同。} \emph{Zhōng-Fēi liǎngguó de yǒuyì yǔ yǐqián xiāngtóng.}
\tcblower
\textbf{Solution 1 :} On repère \zh{一样}, on insère \zh{不} devant.\\
\zh{杜阿拉的天气跟雅温得\bfseries 不\mdseries 一样热。} \emph{Dùālā de tiānqì gēn Yǎwēndé bù yíyàng rè.}\\
« Il ne fait pas aussi chaud à Douala qu'à Yaoundé. »\\
\textbf{Solution 2 :} \zh{相同} (formel) $\rightarrow$ \zh{不同}.\\
\zh{中非两国的友谊与以前\bfseries 不同\mdseries 。} \emph{Zhōng-Fēi liǎngguó de yǒuyì yǔ yǐqián bùtóng.}\\
« L'amitié sino-camerounaise est différente d'avant. »\\
\emph{Piège évité :} ✗ \zh{杜阿拉的天气不跟雅温得一样热。}
\end{exoresolu}

\begin{attention}[Pièges fréquents des francophones — Synthèse]
\begin{itemize}
\item \textbf{Négation mal placée :} ✗ \zh{我不跟他一样} $\rightarrow$ ✓ \zh{我跟他不一样}. (Le français nie le verbe « je \textbf{ne} suis \textbf{pas} », le chinois nie l'adjectif « pareil »).
\item \textbf{Mot-outil manquant :} ✗ \zh{我不同他} $\rightarrow$ ✓ \zh{我跟他不一样 / 我跟他不同}. (\zh{跟/和/与} est obligatoire).
\item \textbf{不同 + adjectif :} ✗ \zh{不同难} $\rightarrow$ ✓ \zh{不一样难}.
\item \textbf{Oubli de 的 :} ✗ \zh{不同文化} $\rightarrow$ ✓ \zh{不同的文化} (attribut du nom).
\item \textbf{Ordre A / B :} \zh{A 不一样 adj} = A < B. Ne pas inverser A et B par rapport au français visé.
\end{itemize}
\end{attention}

\begin{experience}[Activité orale : « Le Cameroun et la Chine : pareil ou différent ? »]
\textbf{Consigne :} En binôme, utilisez les structures étudiées pour comparer 5 paires d'éléments (un camerounais, un chinois). Alternez égalité et différence.
\par\medskip
\textbf{Modèle :}
\begin{itemize}
\item A : \zh{喀麦隆的足球跟中国的足球一样受欢迎吗？}
\item B : \zh{不一样受欢迎。喀麦隆的足球更受欢迎。} (ou \zh{不一样。})
\item A : \zh{喀麦隆的米饭跟中国的米饭相同吗？}
\item B : \zh{不一样。喀麦隆的米饭跟中国的米饭不一样香。}
\end{itemize}
\textbf{Pistes :} ndolé / riz cantonnais ; climat Yaoundé / Pékin ; lycée / 高中 ; transport (moto-taxi / métro) ; salutations.
\end{experience}

\begin{savaistu}[Culture : La modestie dans la comparaison chinoise]
En chinois, quand on compare deux personnes (surtout soi-même et autrui), la culture confucéenne valorise la \cle{modestie} (谦虚 qiānxū).
\begin{itemize}
\item On évite souvent de dire \zh{我跟他一样聪明} « Je suis aussi intelligent que lui » (trop direct, vaniteux).
\item On préfère : \zh{我跟他不一样聪明} « Je ne suis pas aussi intelligent que lui » (même si c'est faux, c'est poli).
\item À l'inverse, pour complimenter : \zh{你比我聪明} « Tu es plus intelligent que moi ».
\end{itemize}
C'est l'inverse de l'assertivité française souvent attendue à l'oral (exposé, entretien). Au Bac oral, montrez que vous maîtrisez \zh{不一样} pour nuancer vos propos.
\end{savaistu}

\begin{exercice}[Entraînement complet]
\textbf{1. Transformez en phrase de différence (garde le même A et B) :}
\begin{enumerate}
\item \zh{汉语跟法语一样难。}
\item \zh{我的房间跟你的房间一样大。}
\item \zh{这两个词的意思相同。}
\item \zh{杜阿拉和雅温得一样热。}
\end{enumerate}
\textbf{2. Traduisez en chinois (caractères + pinyin) :}
\begin{enumerate}
\item Le football et le basketball ne sont pas pareils.
\item Le climat de Buea n'est pas aussi chaud que celui de Douala.
\item La Chine et le Cameroun ont des cultures différentes.
\item Ma méthode d'apprentissage n'est pas la même que la tienne.
\end{enumerate}
\textbf{3. Corrigez les erreurs :}
\begin{enumerate}
\item \zh{我不同意他一样高。}
\item \zh{这本书跟那本书不同贵。}
\item \zh{我们有不一样的文化。} (contexte formel/écrit)
\end{enumerate}
\end{exercice}

\begin{corrige}[Exercice 1]
\textbf{1. Transformation :}
\begin{enumerate}
\item \zh{汉语跟法语\bfseries 不\mdseries 一样难。} \emph{Hànyǔ gēn Fǎyǔ bù yíyàng nán.} (Le chinois n'est pas aussi difficile que le français.)
\item \zh{我的房间跟你的房间\bfseries 不\mdseries 一样大。} \emph{Wǒ de fángjiān gēn nǐ de fángjiān bù yíyàng dà.} (Ma chambre n'est pas aussi grande que la tienne.)
\item \zh{这两个词的意思\bfseries 不同。} \emph{Zhè liǎng ge cí de yìsi bùtóng.} (Le sens de ces deux mots est différent.)
\item \zh{杜阿拉和雅温得\bfseries 不\mdseries 一样热。} \emph{Dùālā hé Yǎwēndé bù yíyàng rè.} (Douala n'est pas aussi chaud que Yaoundé.)
\end{enumerate}
\textbf{2. Traduction :}
\begin{enumerate}
\item \zh{足球跟篮球不一样。} \emph{Zúqiú gēn lánqiú bù yíyàng.}
\item \zh{布埃亚的气候跟杜阿拉不一样热。} \emph{Bù'āiyà de qìhòu gēn Dùālā bù yíyàng rè.}
\item \zh{中国和喀麦隆有不同的文化。} \emph{Zhōngguó hé Kāmàilóng yǒu bùtóng de wénhuà.}
\item \zh{我的学习方法跟你的不一样。} \emph{Wǒ de xuéxí fāngfǎ gēn nǐ de bù yíyàng.} (ou \zh{...跟你的不同。})
\end{enumerate}
\textbf{3. Correction :}
\begin{enumerate}
\item ✗ \zh{我不同意他一样高。} (\zh{同意} = être d'accord, confusion phonétique) $\rightarrow$ ✓ \zh{我跟他不一样高。}
\item ✗ \zh{这本书跟那本书不同贵。} (\zh{不同} + adjectif impossible) $\rightarrow$ ✓ \zh{这本书跟那本书不一样贵。}
\item ✗ \zh{我们有不一样的文化。} (registre oral) $\rightarrow$ ✓ \zh{我们有不同的文化。} (registre formel/écrit attendu).
\end{enumerate}
\end{corrige}

\coursec{Comparaison avec le français et pièges des francophones}

Après avoir analysé les schémas canoniques de l'égalité (\zh{A跟/和/与B一样/相同}) et de la différence (\zh{A跟/和/与B不一样/不同}), nous devons maintenant comprendre \emph{pourquoi} le passage du français au chinois génère systématiquement des erreurs d'ordre des mots chez les apprenants camerounais. Le français et le chinois n'organisent pas la comparaison selon la même logique syntaxique : le français est une langue \emph{prépositionnelle} (la relation se marque avant le terme de comparaison), le chinois est une langue \emph{sérielle} (la relation se construit par enchaînement d'éléments dans un ordre fixe). Cette section identifie les correspondances exactes, explique le choc des ordres des mots et dresse la liste des cinq pièges fatals au Bac.

\courssub{Correspondances fondamentales : « aussi… que », « identique à », « différent de »}

Avant d'aborder les pièges, fixons les équivalences fonctionnelles. Le français utilise trois constructions distinctes selon la nuance (égalité de degré, identité de nature, différence) ; le chinois unifie l'outillage (\zh{跟/和/与}) et ne varie que le noyau prédicatif (\zh{一样/相同} vs \zh{不一样/不同}).

\begin{center}
\begin{tabularx}{\linewidth}{|X|X|X|X|}
\hline
\rowcolor{popblueL}
\textbf{Notion française} & \textbf{Structure chinoise type} & \textbf{Exemple (caractères + pinyin)} & \textbf{Traduction} \\
\hline
Aussi \emph{adj} que (égalité de degré) & A 跟/和/与 B 一样 \emph{adj} & \zh{我跟他一样高。 Wǒ gēn tā yíyàng gāo.} & Je suis aussi grand que lui. \\
\hline
Identique à (identité de nature/aspect) & A 跟/和/与 B 相同 & \zh{这本书跟那本相同。 Zhè běn shū gēn nà běn xiāngtóng.} & Ce livre est identique à celui-là. \\
\hline
Différent de (différence de nature/aspect) & A 跟/和/与 B 不同 & \zh{我的想法跟你不同。 Wǒ de xiǎngfǎ gēn nǐ bùtóng.} & Mon opinion est différente de la tienne. \\
\hline
Pas aussi \emph{adj} que / Différent sur un trait & A 跟/和/与 B 不一样 \emph{adj} & \zh{这道菜跟那道不一样辣。 Zhè dào cài gēn nà dào bùyíyàng là.} & Ce plat n'est pas aussi épicé que celui-là / est différemment épicé. \\
\hline
\end{tabularx}
\end{center}

\cle{相同} (xiāngtóng) est plus formel, écrit, et s'emploie souvent sans adjectif suivant (identité globale). \cle{一样} (yíyàng) est la forme passe-partout, orale comme écrite, et \textbf{exige} un adjectif ou un verbe d'état en finale quand on compare un degré (\zh{一样高}, \zh{一样好看}). \cle{不同} (bùtóng) est l'antonyme formel de \cle{相同} ; \cle{不一样} (bùyíyàng) est l'antonyme courant de \cle{一样}.

\begin{attention}[Piège n°0 : Traduire mot à mot « aussi… que »]
En français, « aussi… que » est un \textbf{corrélatif bipartite} : l'adjectif est \emph{encadré} (aussi \underline{grand} que). En chinois, \textbf{il n'existe aucun mot pour « aussi » ni pour « que »} dans cette construction. Le cadre est \zh{跟……一样……} : l'adjectif se place \textbf{après} \zh{一样}, en position de complément de résultat/dégré. Ne cherchez pas à traduire « aussi » par \zh{也} (yě) ou \zh{很} (hěn) : c'est la source de l'erreur n°3 ci-dessous.
\end{attention}

\courssub{L'ordre des mots : le choc syntaxique majeur}

C'est le point unique qui fait perdre le plus de points en traduction thème (français $\to$ chinois) au Baccalauréat.

\begin{propriete}[Règle d'or de l'ordre des mots comparatifs]
\textbf{Français :} Sujet + \emph{aussi} + \textbf{Adjectif} + \emph{que} + Complément de comparaison. \\
\textbf{Chinois :} Sujet (A) + \textbf{跟/和/与} + Complément de comparaison (B) + \textbf{一样/不一样/相同/不同} + \textbf{Adjectif/Verbe d'état}.
\end{propriete}

\emph{Analyse contrastive :}
\begin{itemize}
    \item \textbf{Français :} « Je suis \underline{aussi grand} \underline{que} lui. » $\to$ L'adjectif \emph{grand} est \textbf{avant} l'outil de comparaison \emph{que}.
    \item \textbf{Chinois :} \zh{我跟他一样高。} $\to$ L'adjectif \zh{高} est \textbf{après} le marqueur d'égalité \zh{一样}.
\end{itemize}

\begin{popfigure}
\begin{tikzpicture}[
    node distance=0.8cm and 1.2cm,
    frbox/.style={rectangle, draw=popblue, thick, fill=popblueL, rounded corners, minimum height=1.1cm, text width=3.2cm, align=center, font=\small},
    cnbox/.style={rectangle, draw=popgreen, thick, fill=popgreenL, rounded corners, minimum height=1.1cm, text width=3.2cm, align=center, font=\small},
    errbox/.style={rectangle, draw=poporange, thick, fill=poporangeL, rounded corners, minimum height=1.1cm, text width=3.2cm, align=center, font=\small},
    arrow/.style={-Stealth, thick, popdark},
    labelstyle/.style={font=\bfseries\small, popdark}
]
% En-têtes
\node[labelstyle] (hfr) at (0, 2.5) {Français (source)};
\node[labelstyle] (hcalque) at (4.5, 2.5) {Calque fautif \ding{55}};
\node[labelstyle] (hcn) at (9, 2.5) {Chinois correct \ding{51}};

% Ligne 1 : Égalité simple
\node[frbox] (fr1) at (0, 1.2) {Je suis \textbf{aussi grand} \underline{que} lui.};
\node[errbox, align=center] (er1) at (4.5, 1.2){\zh{我也是高跟他} \\ \emph{Wǒ yě shì gāo gēn tā}};
\node[cnbox, align=center] (cn1) at (9, 1.2){\zh{我跟他一样高} \\ \emph{Wǒ gēn tā yíyàng gāo}};

% Ligne 2 : Identité
\node[frbox] (fr2) at (0, 0) {Ce livre est \textbf{identique à} celui-ci.};
\node[errbox, align=center] (er2) at (4.5, 0){\zh{这本书相同那本} \\ \emph{Zhè běn shū xiāngtóng nà běn}};
\node[cnbox, align=center] (cn2) at (9, 0){\zh{这本书跟那本相同} \\ \emph{Zhè běn shū gēn nà běn xiāngtóng}};

% Ligne 3 : Différence
\node[frbox] (fr3) at (0, -1.2) {Mon avis \textbf{n'est pas le même que} le tien.};
\node[errbox, align=center] (er3) at (4.5, -1.2){\zh{我意见不一样你} \\ \emph{Wǒ yìjiàn bùyíyàng nǐ}};
\node[cnbox, align=center] (cn3) at (9, -1.2){\zh{我意见跟你不一样} \\ \emph{Wǒ yìjiàn gēn nǐ bùyíyàng}};

% Flèches
\draw[arrow] (fr1.east) -- (er1.west) node[midway, above, font=\tiny, poporange] {Traduction mot à mot};
\draw[arrow, popgreen] (fr1.east) .. controls +(right:2cm) and +(left:2cm) .. (cn1.west) node[midway, above, font=\tiny, popgreen] {Restructuration};
\draw[arrow] (fr2.east) -- (er2.west);
\draw[arrow, popgreen] (fr2.east) .. controls +(right:2cm) and +(left:2cm) .. (cn2.west);
\draw[arrow] (fr3.east) -- (er3.west);
\draw[arrow, popgreen] (fr3.east) .. controls +(right:2cm) and +(left:2cm) .. (cn3.west);
\end{tikzpicture}
\legende{Tableau comparatif : le français place l'adjectif \textbf{avant} l'outil de liaison ; le chinois le place \textbf{après} \zh{一样/不同}. Le calque mot à mot (colonne centrale) produit des phrases incompréhensibles.}
\end{popfigure}

\courssub{Les cinq pièges classiques des apprenants camerounais}

Voici les cinq erreurs les plus fréquentes observées dans les copies du MINESEC, classées par ordre d'importance pédagogique. Chacune fait l'objet d'une boîte \texttt{attention} et d'exemples corrigés.

\begin{attention}[Piège n°1 : Oublier le mot-outil \zh{跟/和/与} (✗ \zh{我他一样高})]
En français, « que » suffit à introduire le second terme. En chinois, \textbf{aucun lien syntaxique n'existe entre A et B sans préposition}. \zh{跟} (oral), \zh{和} (neutre), \zh{与} (formel/écrit) sont \textbf{obligatoires}.
\end{attention}
\begin{exemplebox}[Correction Piège 1]
✗ \zh{我弟弟我一样高。} (Wǒ dìdi wǒ yíyàng gāo.) \\
✓ \zh{我跟弟弟一样高。} (Wǒ gēn dìdi yíyàng gāo.) — « Je suis aussi grand que mon petit frère. » \\
✓ \zh{这件衬衫和那件相同。} (Zhè jiàn chènshān hé nà jiàn xiāngtóng.) — « Cette chemise est identique à celle-là. »
\end{exemplebox}

\begin{attention}[Piège n°2 : Placer \zh{一样} avant le terme B (✗ \zh{我一样跟他高})]
L'ordre \zh{Sujet + 一样 + 跟 + B + Adj} est \textbf{grammaticalement impossible}. La structure est une chaîne : \zh{A} $\to$ \zh{跟 B} $\to$ \zh{一样} $\to$ \zh{Adj}. \zh{一样} est le pivot qui ferme la comparaison et ouvre le degré.
\end{attention}
\begin{exemplebox}[Correction Piège 2]
✗ \zh{喀麦隆一样中国很大。} (Kāmàilóng yíyàng Zhōngguó hěn dà.) \\
✓ \zh{喀麦隆跟中国一样大。} (Kāmàilóng gēn Zhōngguó yíyàng dà.) — « Le Cameroun est aussi grand que la Chine. » \\
✗ \zh{他不一样我忙。} (Tā bùyíyàng wǒ máng.) \\
✓ \zh{他跟我不一样忙。} (Tā gēn wǒ bùyíyàng máng.) — « Il n'est pas aussi occupé que moi / Il est différemment occupé. »
\end{exemplebox}

\begin{attention}[Piège n°3 : Calquer « aussi » par \zh{也} (aussi) ou \zh{很} (très) (✗ \zh{我也跟他高})]
\zh{也} signifie « aussi » au sens de « de plus, également » (addition), pas « à un degré égal ». \zh{很} est un marqueur de degré neutre (très), non comparatif. L'égalité de degré est \textbf{exclusivement} portée par \zh{一样}.
\end{attention}
\begin{exemplebox}[Correction Piège 3]
✗ \zh{我也跟他高。} (Wǒ yě gēn tā gāo.) $\to$ Sens : « Moi aussi, je suis grand (par rapport à un tiers), et lui aussi. » Pas de comparaison directe. \\
✗ \zh{我很跟他高。} (Wǒ hěn gēn tā gāo.) $\to$ Ingrammatical. \zh{很} ne précède pas \zh{跟}. \\
✓ \zh{我跟他一样高。} (Wǒ gēn tā yíyàng gāo.) — Seul \zh{一样} exprime l'égalité du degré.
\end{exemplebox}

\begin{attention}[Piège n°4 : Mettre la négation au mauvais endroit (✗ \zh{我不跟他一样高})]
La négation \zh{不} nie \textbf{l'égalité/identité}, pas la préposition \zh{跟}. \zh{不跟} signifierait « ne pas suivre / ne pas être avec ». La place de \zh{不} est \textbf{immédiatement devant \zh{一样} ou \zh{同}}.
\end{attention}
\begin{exemplebox}[Correction Piège 4]
✗ \zh{我不跟他一样高。} (Wǒ bù gēn tā yíyàng gāo.) $\to$ « Je ne suis \textbf{pas avec} lui de la même grandeur. » (Absurde). \\
✓ \zh{他跟我不一样高。} (Tā gēn wǒ \textbf{bù}yíyàng gāo.) — « Je ne suis pas aussi grand que lui. » \\
✗ \zh{这个苹果不和那个相同。} \\
✓ \zh{这个苹果和那个不同。} (Zhè ge píngguǒ hé nà ge \textbf{bù}tóng.) — « Cette pomme est différente de celle-là. »
\end{exemplebox}

\begin{attention}[Piège n°5 : Confondre \zh{一样} (pareil) et \zh{一点儿} (un peu) dans les réponses]
À l'oral, face à « \zh{你跟他一样高吗?} » (Est-tu aussi grand que lui ?), l'élève paniqué répond souvent \zh{不，我一点儿高} (Non, je suis un peu grand) au lieu de \zh{不，我跟他不一样高} (Non, je ne suis pas aussi grand que lui). \zh{一点儿} modifie l'adjectif seul ; \zh{不一样} nie la relation comparative.
\end{attention}
\begin{exemplebox}[Correction Piège 5 — Mini-dialogue]
\zh{— 你跟你哥哥一样高吗？} (Nǐ gēn nǐ gēge yíyàng gāo ma?) — « Es-tu aussi grand que ton grand frère ? » \\
\zh{— 不，我跟他不一样高。我比他矮一点儿。} (Bù, wǒ gēn tā \textbf{bùyíyàng} gāo. Wǒ bǐ tā ǎi yìdiǎnr.) — « Non, je ne suis pas aussi grand que lui. Je suis un peu plus petit que lui. » \\
\emph{Note :} La réponse négative à une question \zh{一样…吗?} exige \zh{不一样} (ou \zh{不同}), jamais \zh{一点儿} seul.
\end{exemplebox}

\courssub{Méthode de vérification et organigramme de décision}

Avant de rendre une copie ou de parler, appliquez ce contrôle qualité en 4 étapes. C'est la « check-list » du candidat au Bac.

\begin{methode}[Vérifier sa phrase comparative en 4 étapes]
\textbf{Étape 1 — Deux termes ?} Ai-je bien un terme A (sujet) et un terme B (objet de comparaison) distincts ? \\
\textbf{Étape 2 — Mot-outil ?} Ai-je mis \zh{跟} (oral), \zh{和} (standard) ou \zh{与} (formel) \textbf{entre} A et B ? (Jamais au début, jamais à la fin). \\
\textbf{Étape 3 — Noyau comparatif ?} Ai-je mis \zh{一样} / \zh{相同} (égalité) OU \zh{不一样} / \zh{不同} (différence) \textbf{immédiatement après B} ? \\
\textbf{Étape 4 — Adjectif final ?} Si j'utilise \zh{一样} ou \zh{不一样} pour comparer un \textbf{degré}, l'adjectif (\zh{高}, \zh{好}, \zh{忙}, \zh{辣}…) est-il bien en \textbf{toute dernière position} ?
\end{methode}

L'organigramme ci-dessous résume le choix des mots-outils selon le registre et le type de relation.

\begin{popfigure}
\begin{tikzpicture}[
    node distance=1.1cm and 1.5cm,
    start/.style={ellipse, draw=popblue, thick, fill=popblueL, minimum height=1cm, text width=2.8cm, align=center, font=\small\bfseries},
    decision/.style={diamond, draw=poporange, thick, fill=poporangeL, aspect=2, minimum height=1cm, text width=2.5cm, align=center, font=\small},
    choice/.style={rectangle, draw=popgreen, thick, fill=popgreenL, rounded corners, minimum height=0.9cm, text width=2.8cm, align=center, font=\small},
    end/.style={rectangle, draw=poppurple, thick, fill=poppurpleL, rounded corners, minimum height=0.9cm, text width=3cm, align=center, font=\small\bfseries},
    arrow/.style={-Stealth, thick, popdark},
    label/.style={font=\tiny, popdark, midway, fill=white, inner sep=1pt}
]
% Nœuds
\node[start] (start) {Je veux comparer A et B};
\node[decision, below of=start, align=center] (eqdiff){Égalité / Identité \\ OU Différence ?};
\node[choice, below left=1.2cm and 0.5cm of eqdiff, align=center] (eq){ÉGALITÉ / IDENTITÉ \\ \zh{一样} / \zh{相同}};
\node[choice, below right=1.2cm and 0.5cm of eqdiff, align=center] (diff){DIFFÉRENCE \\ \zh{不一样} / \zh{不同}};
\node[decision, below of=eq] (rege) {Registre ?};
\node[decision, below of=diff] (regd) {Registre ?};
\node[end, below left=1cm and 0.5cm of rege, align=center] (eqoral){ORAL / COURANT \\ \zh{跟} + \zh{一样}};
\node[end, below right=1cm and 0.5cm of rege, align=center] (eqwrit){ÉCRIT / FORMEL \\ \zh{和} / \zh{与} + \zh{相同}};
\node[end, below left=1cm and 0.5cm of regd, align=center] (diforal){ORAL / COURANT \\ \zh{跟} + \zh{不一样}};
\node[end, below right=1cm and 0.5cm of regd, align=center] (difwrit){ÉCRIT / FORMEL \\ \zh{和} / \zh{与} + \zh{不同}};

% Flèches
\draw[arrow] (start) -- (eqdiff);
\draw[arrow] (eqdiff) -- node[label, left] {Oui} (eq);
\draw[arrow] (eqdiff) -- node[label, right] {Non} (diff);
\draw[arrow] (eq) -- (rege);
\draw[arrow] (diff) -- (regd);
\draw[arrow] (rege) -- node[label, left] {Oral} (eqoral);
\draw[arrow] (rege) -- node[label, right] {Écrit} (eqwrit);
\draw[arrow] (regd) -- node[label, left] {Oral} (diforal);
\draw[arrow] (regd) -- node[label, right] {Écrit} (difwrit);
\end{tikzpicture}
\legende{Organigramme de décision : le choix entre \zh{跟/和/与} et \zh{一样/相同/不一样/不同} dépend de la relation (égalité vs différence) ET du registre (oral vs écrit).}
\end{popfigure}

\courssub{Exercice de diagnostic : corrigez les erreurs typiques}

Appliquez la méthode en 4 étapes aux 5 phrases suivantes, produites par des élèves de Terminale A lors d'un entraînement Bac. Chaque phrase contient \textbf{au moins une} erreur majeure (parfois deux). Corrigez et justifiez.

\begin{exoresolu}[Exercice de diagnostic — 5 phrases à corriger]
\textbf{Énoncé :} Corrigez les phrases ci-dessous. Pour chaque correction, indiquez le(s) piège(s) évité(s) (n°1 à n°5).

1. \zh{我妈妈我一样喜欢吃椒盐鱼。} \\
2. \zh{这个手机和那个相同贵。} \\
3. \zh{他不跟我一样忙，他更忙。} \\
4. \zh{喀麦隆足球队也中国队一样强。} \\
5. \zh{— 你跟老师一样高吗？ — 不，我一点儿高。} \\

\tcblower

\textbf{Solution détaillée :}

\textbf{1.} ✗ \zh{我妈妈我一样喜欢吃椒盐鱼。} \\
✓ \zh{我跟妈妈一样喜欢吃椒盐鱼。} (Wǒ \textbf{gēn} māma \textbf{yíyàng} xǐhuan chī jiāoyán yú.) \\
\textit{Justification :} \textbf{Piège n°1} (absence de \zh{跟}) + \textbf{Piège n°2} (\zh{一样} placé avant B). Le verbe \zh{喜欢} étant un verbe d'état mental, on peut le suivre directement après \zh{一样} sans adjectif supplémentaire.

\textbf{2.} ✗ \zh{这个手机和那个相同贵。} \\
✓ \zh{这个手机和那个一样贵。} (Zhè ge shǒujī hé nà ge \textbf{yíyàng} guì.) \\
\textit{Justification :} \textbf{Piège n°2 / Choix lexical}. \zh{相同} (identique) s'emploie \textbf{sans adjectif suivant} pour dire « c'est le même objet / même prix / même aspect global ». Pour comparer le \textbf{degré} d'un adjectif (\zh{贵}), on \textbf{doit} utiliser \zh{一样} + Adj. Si on veut dire « même prix » nominale : \zh{价格相同} (jiàgé xiāngtóng).

\textbf{3.} ✗ \zh{他不跟我一样忙，他更忙。} \\
✓ \zh{他跟我不一样忙，他更忙。} (Tā gēn wǒ \textbf{bù}yíyàng máng, tā gèng máng.) \\
\textit{Justification :} \textbf{Piège n°4} (négation mal placée). \zh{不跟} = « ne pas suivre / ne pas être avec ». La négation porte sur l'égalité (\zh{不一样}), pas sur la préposition. Note : \zh{不一样忙} implique souvent une différence de degré (moins ou plus) ; la suite « \zh{他更忙} » précise qu'il est \emph{plus} occupé.

\textbf{4.} ✗ \zh{喀麦隆足球队也中国队一样强。} \\
✓ \zh{喀麦隆足球队跟中国队一样强。} (Kāmàilóng zúqiú duì \textbf{gēn} Zhōngguó duì \textbf{yíyàng} qiáng.) \\
\textit{Justification :} \textbf{Piège n°1} (absence de \zh{跟}) + \textbf{Piège n°3} (calque « aussi » $\to$ \zh{也}). \zh{也} est un adverbe d'addition (« aussi » = « de plus »), incapable de construire une comparaison d'égalité.

\textbf{5.} ✗ \zh{— 你跟老师一样高吗？ — 不，我一点儿高。} \\
✓ \zh{— 你跟老师一样高吗？ — 不，我跟老师不一样高。（或：我比老师矮一点儿。）} \\
\textit{Justification :} \textbf{Piège n°5}. La réponse négative à une question \zh{一样…吗?} exige \zh{不一样} (ou \zh{不同}) pour nier la relation d'égalité. \zh{一点儿高} signifie « un peu grand » (description absolue), ce qui ne répond pas à la question comparative.
\end{exoresolu}

\begin{savaistu}[Le comparatif au Bac : points faciles, points pièges]
À l'épreuve écrite du Baccalauréat camerounais (LV2 Chinois), les items de traduction thème (Français $\to$ Chinois) portant sur \zh{跟……一样} sont notés \textbf{au mot près}.
\begin{itemize}
    \item \textbf{1 point} pour le mot-outil correct (\zh{跟/和/与}).
    \item \textbf{1 point} pour la place correcte de \zh{一样/不一样} (après B).
    \item \textbf{1 point} pour l'adjectif en position finale.
    \item \textbf{0 point} si l'ordre est \zh{Sujet + 一样 + 跟 + B + Adj} (erreur de structure majeure).
    \item \textbf{Pénalité} si \zh{也} ou \zh{很} remplace \zh{一样}.
\end{itemize}
\textbf{Astuce de pro :} Dans la version (Chinois $\to$ Français), si vous voyez \zh{跟……一样 + Adj}, traduisez systématiquement par « aussi [Adj] que ». Si vous voyez \zh{跟……相同} sans adjectif, traduisez par « être identique à » ou « être le même que ». Ne traduisez \textbf{jamais} \zh{相同} par « aussi… que ».
\end{savaistu}

\coursec{Emplois étendus et exceptions (complément utile à l'examen)}

Dans la section précédente, nous avons confronté les structures chinoises \zh{跟/和/与……一样} et \zh{跟/和/与……不一样} au système français de la comparaison, et repéré les pièges typiques des francophones (ordre des mots, oubli de \zh{一样}, calque de « aussi … que »). Nous allons maintenant dépasser le schéma de base pour explorer des emplois plus fins, qui font souvent la différence au baccalauréat, notamment dans les épreuves de traduction (thème et version).

\courssub{Observation : un dialogue au marché de Mokolo}

\begin{textebox}[Au marché de Mokolo (Yaoundé)]
A：你看，这两件衬衫一样吗？\\
\emph{Nǐ kàn, zhè liǎng jiàn chènshān yīyàng ma?}\\
« Regarde, ces deux chemises sont-elles pareilles ? »\\[0.3em]
B：颜色相同，但是大小不一样。\\
\emph{Yánsè xiāngtóng, dànshì dàxiǎo bù yīyàng.}\\
« La couleur est identique, mais la taille n'est pas la même. »\\[0.3em]
A：我弟弟的身高几乎跟我一样，他说法语说得跟我一样好。\\
\emph{Wǒ dìdi de shēngāo jīhū gēn wǒ yīyàng, tā shuō Fǎyǔ shuō de gēn wǒ yīyàng hǎo.}\\
« La taille de mon petit frère est presque la même que la mienne, et il parle français aussi bien que moi. »\\[0.3em]
B：真的吗？我们上一样的学校，可是我们的成绩完全不同！\\
\emph{Zhēn de ma? Wǒmen shàng yīyàng de xuéxiào, kěshì wǒmen de chéngjì wánquán bùtóng!}\\
« Vraiment ? Nous allons dans la même école, pourtant nos résultats sont complètement différents ! »
\end{textebox}

\textbf{Observons.}\par
Relevez dans ce dialogue toutes les formes de comparaison. Vous remarquez que \zh{一样} et \zh{不同} apparaissent dans des positions nouvelles : devant un nom (\zh{一样的学校}), après un verbe avec \zh{得} (\zh{说得跟我一样好}), avec des adverbes (\zh{几乎一样}, \zh{完全不同}), et même sans deuxième terme de comparaison exprimé (\zh{这两件衬衫一样吗？}). C'est précisément l'objet de cette section.

\courssub{Comparaison portant sur un verbe : A + verbe + 得 + 跟/和 B + 一样 + adjectif}

Quand on compare non pas une qualité générale mais la \emph{manière} dont une action est accomplie (« il parle chinois aussi bien que moi », « elle court aussi vite que son frère »), le chinois utilise le complément de degré introduit par \zh{得}.

\begin{propriete}[Structure : comparaison d'une action]
\textbf{A + verbe + (objet) + 得 + 跟/和 + B + 一样 + adjectif}\\
Si le verbe a un objet, on répète le verbe : \zh{说汉语说得……} ou on avance l'objet : \zh{他汉语说得跟我一样好。}
\begin{itemize}
\item \zh{他说汉语说得跟我一样好。} \emph{Tā shuō Hànyǔ shuō de gēn wǒ yīyàng hǎo.} « Il parle chinois aussi bien que moi. »
\item \zh{她唱歌唱得跟老师一样好听。} \emph{Tā chànggē chàng de gēn lǎoshī yīyàng hǎotīng.} « Elle chante aussi joliment que la professeure. »
\item \zh{我弟弟足球踢得跟我不一样好。} \emph{Wǒ dìdi zúqiú tī de gēn wǒ bù yīyàng hǎo.} « Mon petit frère ne joue pas au football aussi bien que moi. »
\end{itemize}
\end{propriete}

\begin{attention}[Piège de traduction au Bac]
En version, quand vous lisez \zh{他说汉语说得跟我一样好}, ne traduisez pas « il dit le chinois dit comme moi bien » ! Repérez la chaîne \cle{verbe + 得 + 跟……一样 + adjectif} et rendez-la par « aussi + adverbe/adjectif + que ». En thème, pour traduire « elle danse aussi bien que sa sœur », pensez à répéter le verbe : \zh{她跳舞跳得跟她姐姐一样好。}
\end{attention}

\courssub{一样 et 不同 devant un nom : « le même », « différent »}

Placés devant un nom (avec \zh{的}), \zh{一样} signifie « le même, semblable » et \zh{不同} signifie « différent ». C'est un emploi \emph{adjectival}, très fréquent à l'écrit.

\begin{exemplebox}[一样 / 不同 + 的 + nom]
\begin{itemize}
\item \zh{一样的学校} \emph{yīyàng de xuéxiào} « la même école »
\item \zh{一样的想法} \emph{yīyàng de xiǎngfǎ} « la même idée »
\item \zh{不同的人} \emph{bùtóng de rén} « des personnes différentes »
\item \zh{不同的文化} \emph{bùtóng de wénhuà} « des cultures différentes »
\item \zh{中国和喀麦隆有不同的传统。} \emph{Zhōngguó hé Kāmàilóng yǒu bùtóng de chuántǒng.} « La Chine et le Cameroun ont des traditions différentes. »
\item \zh{我们在一样的学校学习。} \emph{Wǒmen zài yīyàng de xuéxiào xuéxí.} « Nous étudions dans la même école. »
\end{itemize}
\end{exemplebox}

\courssub{Emplois figés : adverbes de degré devant 一样 / 不同}

On peut nuancer la comparaison avec un adverbe de degré placé \emph{devant} \zh{一样} ou \zh{不同} :

\begin{center}
\begin{tabularx}{\linewidth}{|X|X|X|X|}
\hline
\rowcolor{popblueL} Caractères & Pinyin & Nature & Traduction \\
\hline
完全一样 & wánquán yīyàng & locution adv. & tout à fait pareil \\
\hline
几乎一样 & jīhū yīyàng & locution adv. & presque pareil \\
\hline
差不多一样 & chàbuduō yīyàng & locution adv. & à peu près pareil \\
\hline
完全不同 & wánquán bùtóng & locution adv. & complètement différent \\
\hline
很不一样 & hěn bù yīyàng & locution adv. & très différent (pas pareil du tout) \\
\hline
\end{tabularx}
\end{center}

\begin{exemplebox}[Adverbes de degré en contexte]
\begin{itemize}
\item \zh{这两个汉字完全一样。} \emph{Zhè liǎng ge Hànzì wánquán yīyàng.} « Ces deux caractères sont exactement pareils. »
\item \zh{我的身高几乎跟我爸爸一样。} \emph{Wǒ de shēngāo jīhū gēn wǒ bàba yīyàng.} « Ma taille est presque la même que celle de mon père. »
\item \zh{杜阿拉和雅温得的天气很不一样。} \emph{Dùālā hé Yǎwēndé de tiānqì hěn bù yīyàng.} « Le temps à Douala et à Yaoundé est très différent. »
\end{itemize}
\end{exemplebox}

\courssub{Restrictions et exceptions à connaître}

\begin{propriete}[相同 ne s'emploie pas devant un adjectif]
\zh{相同} exprime l'\emph{identité} d'une chose (sens, couleur, opinion) et s'emploie \textbf{sans adjectif} après lui. Pour préciser le point de comparaison (taille, prix, âge…), on utilise obligatoirement \cle{一样 + adjectif}.\\
On dit : \zh{他们一样高。} \emph{Tāmen yīyàng gāo.} « Ils sont de la même taille. »\\
On ne dit \textbf{jamais} : ✗ \zh{他们相同高。}
\end{propriete}

\begin{exemplebox}[相同 : identité d'une chose]
\begin{itemize}
\item \zh{这两个词的意思相同。} \emph{Zhè liǎng ge cí de yìsi xiāngtóng.} « Le sens de ces deux mots est identique. »
\item \zh{我们的意见相同。} \emph{Wǒmen de yìjiàn xiāngtóng.} « Nos opinions sont identiques. »
\item \zh{这两件衣服的颜色相同，大小不同。} \emph{Zhè liǎng jiàn yīfu de yánsè xiāngtóng, dàxiǎo bùtóng.} « Ces deux habits ont la même couleur, mais des tailles différentes. »
\end{itemize}
\end{exemplebox}

\begin{attention}[不一样高 ≠ « petit »]
\zh{不一样高} \emph{bù yīyàng gāo} signifie « ne pas avoir la même taille », et rien de plus. Il ne dit pas \emph{qui} est grand ou petit. En traduction, ne sur-interprétez jamais : \zh{他们俩不一样高} se traduit « ils ne font pas la même taille », et non « l'un est petit ». De même, \zh{不一样贵} = « pas au même prix », pas « bon marché ».
\end{attention}

\begin{propriete}[与 : registre soutenu, réservé à l'écrit]
\zh{与} \emph{yǔ} (« avec, et ») est la variante formelle de \zh{跟} et \zh{和}. À l'oral quotidien, on utilise presque toujours \zh{跟} ou \zh{和}. À l'examen écrit, dans un texte formel (lettre officielle, article, dissertation), \cle{与……相同/不同} est \textbf{valorisé} : \zh{中国与喀麦隆的情况不同。} \emph{Zhōngguó yǔ Kāmàilóng de qíngkuàng bùtóng.} « La situation de la Chine et celle du Cameroun sont différentes. »
\end{propriete}

\courssub{Comparaison implicite : quand B n'est pas exprimé}

Quand le deuxième terme de la comparaison est évident (deux objets déjà présents dans le contexte), on peut le sous-entendre : \zh{一样} et \zh{不同} fonctionnent alors comme de simples prédicats.

\begin{exemplebox}[Comparaison sans B exprimé]
\begin{itemize}
\item \zh{这两件衣服一样。} \emph{Zhè liǎng jiàn yīfu yīyàng.} « Ces deux habits sont pareils. »
\item \zh{这两个班的学生一样多。} \emph{Zhè liǎng ge bān de xuésheng yīyàng duō.} « Les élèves de ces deux classes sont aussi nombreux. »
\item \zh{他们的性格很不同。} \emph{Tāmen de xìnggé hěn bùtóng.} « Leurs caractères sont très différents. »
\end{itemize}
Notez que le sujet doit désigner \emph{au moins deux} entités (\zh{这两件衣服}, \zh{他们}) : la comparaison est interne au sujet.
\end{exemplebox}

\courssub{Mini-synthèse : les quatre structures de la comparaison}

\begin{center}
\begin{tabularx}{\linewidth}{|X|X|X|X|}
\hline
\rowcolor{popblueL} Structure & Schéma & Exemple chinois + pinyin & Traduction \\
\hline
一样 (égalité) & A 跟/和 B 一样 (+ adj.) & \zh{我跟他一样高。} \emph{Wǒ gēn tā yīyàng gāo.} & Je suis aussi grand que lui. \\
\hline
相同 (identité) & A 和/与 B 相同 & \zh{这两个词的意思相同。} \emph{Zhè liǎng ge cí de yìsi xiāngtóng.} & Le sens de ces deux mots est identique. \\
\hline
不一样 (différence de degré) & A 跟/和 B 不一样 (+ adj.) & \zh{杜阿拉跟雅温得不一样热。} \emph{Dùālā gēn Yǎwēndé bù yīyàng rè.} & Douala n'est pas aussi chaud que Yaoundé. \\
\hline
不同 (différence catégorique) & A 和/与 B 不同 & \zh{中国和喀麦隆的文化不同。} \emph{Zhōngguó hé Kāmàilóng de wénhuà bùtóng.} & Les cultures de la Chine et du Cameroun sont différentes. \\
\hline
\end{tabularx}
\end{center}

\begin{popfigure}
\begin{tikzpicture}[
  box/.style={draw=popblue, fill=popblueL, rounded corners=2pt, align=center, font=\small, minimum height=1.1cm},
  ex/.style={draw=popmuted, fill=popcream, rounded corners=2pt, align=center, font=\footnotesize},
  lbl/.style={font=\small\bfseries, text=popdark}
]
\node[lbl] at (0,0) {Structure};
\node[lbl] at (3.6,0) {Schéma};
\node[lbl] at (7.4,0) {Exemple};
\node[lbl] at (11.4,0) {Traduction};
\node[box, minimum width=2.6cm] at (0,-1.4) {一样};
\node[box, minimum width=4.2cm] at (3.6,-1.4) {A 跟 B 一样 + adj.};
\node[ex, minimum width=4.0cm] at (7.4,-1.4) {我跟他一样高。};
\node[ex, minimum width=4.0cm] at (11.4,-1.4) {aussi grand que lui};
\node[box, minimum width=2.6cm] at (0,-2.8) {相同};
\node[box, minimum width=4.2cm] at (3.6,-2.8) {A 与 B 相同 (sans adj.)};
\node[ex, minimum width=4.0cm] at (7.4,-2.8) {意思相同。};
\node[ex, minimum width=4.0cm] at (11.4,-2.8) {sens identique};
\node[box, minimum width=2.6cm] at (0,-4.2) {不一样};
\node[box, minimum width=4.2cm] at (3.6,-4.2) {A 跟 B 不一样 + adj.};
\node[ex, minimum width=4.0cm] at (7.4,-4.2) {今天跟昨天不一样热。};
\node[ex, minimum width=4.0cm] at (11.4,-4.2) {pas aussi chaud qu'hier};
\node[box, minimum width=2.6cm] at (0,-5.6) {不同};
\node[box, minimum width=4.2cm] at (3.6,-5.6) {A 与 B 不同};
\node[ex, minimum width=4.0cm] at (7.4,-5.6) {文化不同。};
\node[ex, minimum width=4.0cm] at (11.4,-5.6) {cultures différentes};
\end{tikzpicture}
\legende{Tableau récapitulatif des quatre structures de comparaison}
\end{popfigure}

\begin{popfigure}
\begin{tikzpicture}[
  root/.style={circle, draw=popdark, fill=poporange, text=white, font=\bfseries, minimum size=2.2cm, align=center},
  br/.style={rectangle, rounded corners=3pt, draw=popblue, fill=popblueL, align=center, font=\small, minimum width=3.4cm, minimum height=1.2cm},
  sub/.style={rectangle, rounded corners=3pt, draw=popmuted, fill=popcream, align=center, font=\footnotesize, minimum width=3.2cm}
]
\node[root, align=center] (c){COMPARER\\比较};
\node[br, above left=1.2cm and 2.2cm of c, align=center] (a){一样\\égalité + adjectif};
\node[br, below left=1.2cm and 2.2cm of c, align=center] (b){相同\\identité (sans adj.)};
\node[br, above right=1.2cm and 2.2cm of c, align=center] (d){不一样\\différence de degré};
\node[br, below right=1.2cm and 2.2cm of c, align=center] (e){不同\\différence catégorique};
\draw[-{Stealth}, thick, popdark] (c) -- (a);
\draw[-{Stealth}, thick, popdark] (c) -- (b);
\draw[-{Stealth}, thick, popdark] (c) -- (d);
\draw[-{Stealth}, thick, popdark] (c) -- (e);
\node[sub, above=0.15cm of a] {一样高 / 完全一样 / 几乎一样};
\node[sub, below=0.15cm of b] {意思相同 / 意见相同};
\node[sub, above=0.15cm of d] {不一样高 / 很不一样};
\node[sub, below=0.15cm of e] {不同的人 / 与……不同};
\end{tikzpicture}
\legende{Carte mentale : choisir la bonne structure de comparaison}
\end{popfigure}

\begin{exoresolu}[Thème : traduction guidée]
Énoncé : traduisez en chinois — « Mon frère et moi, nous sommes dans la même école, mais nos résultats sont très différents ; il parle anglais presque aussi bien que notre professeur. »
\tcblower
Solution détaillée :
\begin{enumerate}
\item « dans la même école » : \zh{一样} + \zh{的} + nom → \zh{在一样的学校} ;
\item « nos résultats sont très différents » : différence catégorique → \zh{我们的成绩很不同} ;
\item « presque aussi bien que » : adverbe \zh{几乎} + comparaison d'action avec \zh{得} → \zh{英语说得几乎跟我们老师一样好} (répétition du verbe \zh{说}).
\end{enumerate}
Traduction complète : \zh{我和我弟弟在一样的学校，可是我们的成绩很不同；他英语说得几乎跟我们老师一样好。}\\
\emph{Wǒ hé wǒ dìdi zài yīyàng de xuéxiào, kěshì wǒmen de chéngjì hěn bùtóng; tā Yīngyǔ shuō de jīhū gēn wǒmen lǎoshī yīyàng hǎo.}
\end{exoresolu}

\begin{exercice}[À vous de jouer]
\begin{enumerate}
\item Traduisez : « Ces deux phrases ont le même sens, mais les mots sont différents. »
\item Corrigez la faute : ✗ \zh{他们相同高。}
\item Traduisez : « Elle chante aussi bien que sa mère. » (utilisez \zh{得})
\item Complétez avec \zh{一样} ou \zh{相同} : \zh{这两件衣服的价钱\_\_\_\_。}
\end{enumerate}
\end{exercice}

\begin{corrige}[Exercice « À vous de jouer »]
\begin{enumerate}
\item \zh{这两个句子的意思相同，可是词不同。} \emph{Zhè liǎng ge jùzi de yìsi xiāngtóng, kěshì cí bùtóng.} (« sens » identique → \zh{相同} ; « mots » différents → \zh{不同}.)
\item \zh{他们一样高。} — \zh{相同} ne peut pas précéder un adjectif ; pour la taille, il faut \zh{一样高}.
\item \zh{她唱歌唱得跟她妈妈一样好。} \emph{Tā chànggē chàng de gēn tā māma yīyàng hǎo.} (verbe répété + \zh{得} + \zh{跟……一样好}.)
\item \zh{一样} : on compare le \emph{prix} (un degré, une mesure) → \zh{价钱一样}. \zh{相同} conviendrait pour une identité de nature (sens, couleur, opinion).
\end{enumerate}
\end{corrige}

\begin{savaistu}[Comparer, une habitude de conversation en Chine]
En Chine, comparer les enfants entre eux est un sujet de conversation quotidien et parfaitement accepté : on compare les résultats scolaires (\zh{成绩} chéngjì), la taille (\zh{身高} shēngāo), l'école fréquentée. La structure \zh{跟……一样} fait partie des tournures les plus fréquentes de la langue parlée : \zh{你跟你妈妈一样高了吗？} \emph{Nǐ gēn nǐ māma yīyàng gāo le ma?} « Tu es déjà aussi grand que ta maman ? » Au Cameroun aussi, les familles aiment comparer les enfants — un point culturel commun qui vous aidera à mémoriser ces structures !
\end{savaistu}

\coursec{Exercices d'application et de transformation}

Dans la section précédente, nous avons vu les emplois étendus de la comparaison d'égalité et de différence (\zh{跟/和/与……一样/相同} et \zh{不一样/不同}), notamment l'ajout d'un complément après \zh{一样} (\zh{跟我一样高} « aussi grand que moi ») et le registre écrit de \zh{与}. Il est temps de vérifier que ces structures sont devenues des automatismes. Cette dernière partie est entièrement consacrée à l'entraînement : reconnaissance, complétion, remise en ordre, transformation, traduction (thème et version) et production écrite. C'est exactement le type d'exercices proposés à l'examen.

\courssub{Méthode générale avant de commencer}

\begin{methode}[Réussir les exercices de traduction avec la comparaison]
Pour traduire une phrase comparative du français vers le chinois (ou l'inverse), procédez toujours en trois étapes :
\begin{enumerate}
\item \textbf{Traduire d'abord le cadre} : repérez A et B, puis posez la charpente \zh{A 跟 B 一样} (ou \zh{A 跟 B 不一样}). Ne pensez pas encore à l'adjectif.
\item \textbf{Insérer l'adjectif ou le verbe à la fin} : en chinois, le point de comparaison se place \emph{après} \zh{一样}, jamais entre A et B. « Aussi grande que » = \zh{跟……一样大}, jamais \zh{大跟……一样}.
\item \textbf{Vérifier les mots-outils} : négation \zh{不} placée devant \zh{一样}/\zh{相同} (et non devant \zh{跟}) ; adjectif seul en fin de phrase, sans \zh{是} ni \zh{很}.
\end{enumerate}
\end{methode}

\begin{attention}[L'ancre de la phrase]
Dans les exercices de remise en ordre, ne commencez jamais par le sujet. Repérez d'abord \zh{一样} ou \zh{不同} : c'est \cle{l'ancre} de la phrase. Une fois l'ancre trouvée, placez \zh{跟}/\zh{和}/\zh{与} juste devant elle, puis répartissez A et B de part et d'autre, et mettez l'adjectif tout à la fin. Cette stratégie évite 90 \% des erreurs d'ordre des mots.
\end{attention}

\begin{center}
\begin{tabularx}{\linewidth}{|X|X|X|}
\hline
\rowcolor{popblueL} Structure & Exemple (caractères + pinyin) & Traduction \\
\hline
A 跟 B 一样 + Adj. & \zh{我跟他一样高。} Wǒ gēn tā yíyàng gāo. & Je suis aussi grand que lui. \\
\hline
A 和 B 相同 & \zh{我的想法和你的相同。} Wǒ de xiǎngfǎ hé nǐ de xiāngtóng. & Mon idée est identique à la tienne. \\
\hline
A 与 B 不同 (écrit) & \zh{我的想法与你的不同。} Wǒ de xiǎngfǎ yǔ nǐ de bùtóng. & Mon idée diffère de la tienne. \\
\hline
A 跟 B 不一样 + Adj. & \zh{中国菜跟喀麦隆菜不一样辣。} Zhōngguó cài gēn Kāmàilóng cài bù yíyàng là. & La cuisine chinoise n'est pas aussi pimentée que la cuisine camerounaise. \\
\hline
Forme interrogative & \zh{你跟他一样高吗？} Nǐ gēn tā yíyàng gāo ma? & Es-tu aussi grand que lui ? \\
\hline
\end{tabularx}
\end{center}

\courssub{Exercice 1 : reconnaissance de la structure}

\begin{exercice}[Égalité ou différence ?]
Soulignez la structure de comparaison dans chaque phrase, puis indiquez si elle exprime l'\textbf{égalité} (E) ou la \textbf{différence} (D).
\begin{enumerate}
\item \zh{杜阿拉跟雅温得不一样大。} (Douala n'est pas aussi grande que Yaoundé.)
\item \zh{我的弟弟跟我一样喜欢足球。}
\item \zh{喀麦隆的时间跟中国的时间不同。}
\item \zh{这本书和那本书相同。}
\item \zh{今天的天气跟昨天的一样。}
\item \zh{他的想法与我的不一样。}
\end{enumerate}
\end{exercice}

\begin{corrige}[Exercice 1]
\begin{enumerate}
\item \zh{跟……不一样} → \textbf{D} (négation \zh{不} devant \zh{一样}).
\item \zh{跟……一样} + verbe \zh{喜欢} → \textbf{E}.
\item \zh{跟……不同} → \textbf{D} (\zh{不同} est la forme négative lexicale de \zh{相同}).
\item \zh{相同} seul → \textbf{E} (structure raccourcie : A 和 B 相同).
\item \zh{跟……一样} → \textbf{E} (B est remplacé par \zh{昨天的}, ellipse fréquente).
\item \zh{与……不一样} → \textbf{D} (\zh{与} marque ici le registre écrit).
\end{enumerate}
\end{corrige}

\courssub{Exercice 2 : texte à trous}

\begin{exercice}[Complétez avec les mots-outils]
Complétez avec : \zh{跟}, \zh{和}, \zh{与}, \zh{一样}, \zh{相同}, \zh{不同}, \zh{不一样}.
\begin{enumerate}
\item \zh{我的学校\_\_\_\_\_\_你的学校\_\_\_\_\_\_大。} (Mon école est aussi grande que la tienne.)
\item \zh{法语\_\_\_\_\_\_汉语不\_\_\_\_\_\_。} (Le français n'est pas identique au chinois.)
\item \zh{他的成绩\_\_\_\_\_\_我的成绩\_\_\_\_\_\_。} (Ses résultats sont identiques aux miens.)
\item \zh{这个城市\_\_\_\_\_\_那个城市\_\_\_\_\_\_热闹。} (Cette ville n'est pas aussi animée que l'autre.)
\end{enumerate}
\end{exercice}

\begin{corrige}[Exercice 2]
\begin{enumerate}
\item \zh{我的学校\textbf{跟}你的学校\textbf{一样}大。} — cadre \zh{A 跟 B 一样} + adjectif final \zh{大}.
\item \zh{法语\textbf{跟}汉语\textbf{不同}。} — \zh{不同} (mot unique) est préférable à \zh{不相同} pour exprimer la différence ; registre neutre.
\item \zh{他的成绩\textbf{和}我的成绩\textbf{相同}。} — \zh{和} possible car phrase déclarative orale ; \zh{与} serait aussi correct à l'écrit.
\item \zh{这个城市\textbf{跟}那个城市\textbf{不一样}热闹。} — négation intégrée dans \zh{不一样}, adjectif \zh{热闹} en finale.
\end{enumerate}
\end{corrige}

\courssub{Exercice 3 : remise en ordre}

\begin{exoresolu}[Remettre les mots dans l'ordre]
Énoncé : remettez dans l'ordre les mots suivants : \zh{高 / 我 / 一样 / 跟 / 他}.
\tcblower
Solution détaillée :
\begin{enumerate}
\item \textbf{Repérer l'ancre} : \zh{一样} est présent → phrase d'égalité.
\item \textbf{Poser le cadre} : \zh{我 跟 他 一样} (A + \zh{跟} + B + \zh{一样}).
\item \textbf{Placer l'adjectif en finale} : \zh{我跟他一样高。} Wǒ gēn tā yíyàng gāo. « Je suis aussi grand que lui. »
\end{enumerate}
Piège évité : ne jamais écrire \zh{我高跟他一样} (calque du français « je suis grand comme lui »).
\end{exoresolu}

\begin{exercice}[À vous]
Remettez dans l'ordre :
\begin{enumerate}
\item \zh{辣 / 不一样 / 这个菜 / 跟 / 那个菜}
\item \zh{相同 / 我的答案 / 与 / 你的答案}
\end{enumerate}
\end{exercice}

\begin{corrige}[Exercice 3 bis]
1. \zh{这个菜跟那个菜不一样辣。} « Ce plat n'est pas aussi pimenté que l'autre. »\\
2. \zh{我的答案与你的答案相同。} « Ma réponse est identique à la tienne. »
\end{corrige}

\courssub{Exercice 4 : transformation de phrases}

\begin{exercice}[Affirmatif, négatif, interrogatif]
Transformez les phrases selon la consigne : (a) passez au négatif ou à l'affirmatif selon le cas ; (b) mettez la phrase obtenue à la forme interrogative avec \zh{吗}.
\begin{enumerate}
\item \zh{雅温得跟杜阿拉一样热。} → négatif → interrogatif
\item \zh{我的学校跟你的学校不一样大。} → affirmatif → interrogatif
\item \zh{这两件衣服相同。} → négatif → interrogatif
\item \zh{他的汉语水平与我的不同。} → affirmatif → interrogatif
\end{enumerate}
\end{exercice}

\begin{corrige}[Exercice 4]
\begin{enumerate}
\item \zh{雅温得跟杜阿拉\textbf{不}一样热。} → \zh{雅温得跟杜阿拉不一样热\textbf{吗}？} (Yaoundé n'est-elle pas aussi chaude que Douala ?)
\item \zh{我的学校跟你的学校一样大。} → \zh{我的学校跟你的学校一样大吗？}
\item \zh{这两件衣服\textbf{不同}。} → \zh{这两件衣服不同吗？} (\zh{不相同} est possible mais moins idiomatique que \zh{不同}).
\item \zh{他的汉语水平与我的相同。} → \zh{他的汉语水平与我的相同吗？}
\end{enumerate}
Règle appliquée : la négation porte sur \zh{一样}/\zh{相同} (\zh{不一样}, \zh{不相同} ou mieux \zh{不同}), jamais sur \zh{跟}. La particule \zh{吗} se place simplement en fin de phrase, sans changer l'ordre des mots.
\end{corrige}

\courssub{Exercices 5 et 6 : thème et version}

\begin{exercice}[Thème : traduisez en chinois]
\begin{enumerate}
\item « Ma sœur est aussi intelligente que moi. »
\item « Le Nigeria est différent du Cameroun. »
\item « Ces deux livres sont identiques. »
\end{enumerate}
\end{exercice}

\begin{corrige}[Exercice 5]
\begin{enumerate}
\item \zh{我的姐姐跟我一样聪明。} Wǒ de jiějie gēn wǒ yíyàng cōngming. — cadre \zh{跟……一样} + adjectif \zh{聪明} en finale, sans \zh{是}.
\item \zh{尼日利亚跟喀麦隆不同。} Nírìlìyà gēn Kāmàilóng bùtóng. — « différent de » = \zh{跟……不同} ; pas d'adjectif final car \zh{不同} suffit.
\item \zh{这两本书相同。} Zhè liǎng běn shū xiāngtóng. — attention au classificateur : \zh{两本书} ; \zh{相同} fonctionne comme prédicat, sans \zh{是}.
\end{enumerate}
\end{corrige}

\begin{exemplebox}[Version guidée]
\zh{中国菜跟喀麦隆菜不一样辣。}\\
\emph{Zhōngguó cài gēn Kāmàilóng cài bù yíyàng là.}\\
« La cuisine chinoise n'est pas aussi pimentée que la cuisine camerounaise. »\\
Analyse : A = \zh{中国菜}, B = \zh{喀麦隆菜}, négation \zh{不一样}, adjectif final \zh{辣} « pimenté ». Le français rend \zh{不一样辣} par « pas aussi pimenté que ».
\end{exemplebox}

\begin{textebox}[Version : phrase à traduire et à commenter]
\zh{我的想法与你的相同。}\\
\emph{Wǒ de xiǎngfǎ yǔ nǐ de xiāngtóng.}\\
Traduction : « Mon idée est identique à la tienne. »\\
Justification : \zh{与} est la conjonction de comparaison du \textbf{registre écrit} (dissertation, lettre formelle, presse). À l'oral, on préfère \zh{跟} ou \zh{和}. À l'examen, employer \zh{与} dans une rédaction valorise la copie.
\end{textebox}

\courssub{Exercice 7 : production écrite guidée}

\begin{exercice}[Comparez deux villes du Cameroun]
Rédigez cinq phrases comparant votre ville à une autre ville du Cameroun, en utilisant au moins une fois chaque structure : \zh{跟……一样} + adjectif, \zh{跟……不一样} + adjectif, \zh{不同}, \zh{相同}. Thèmes suggérés : la taille (\zh{大/小}), le climat (\zh{热/凉快}), les langues (\zh{语言}), la nourriture (\zh{菜}), l'école (\zh{学校}).\\
Modèle : \zh{布埃亚跟雅温得不一样热。} Bù'āiyà gēn Yǎwēndé bù yíyàng rè. « Buea n'est pas aussi chaude que Yaoundé. »
\end{exercice}

\begin{corrige}[Exercice 7 : exemple de production attendue]
\zh{我的家乡跟杜阿拉不一样大，它比较小。} « Ma ville natale n'est pas aussi grande que Douala, elle est plus petite. »\\
\zh{这里的天气跟布埃亚的一样凉快。} « Le climat d'ici est aussi frais qu'à Buea. »\\
\zh{我们说的语言跟北方说的语言不同。} « La langue que nous parlons diffère de celle du Nord. »\\
\zh{我们学校的规则跟你们的相同。} « Le règlement de notre école est identique au vôtre. »\\
\zh{这里的菜跟首都的菜一样辣。} « Les plats d'ici sont aussi pimentés que ceux de la capitale. »
\end{corrige}

\courssub{Correction collective et auto-évaluation}

La correction se fait au tableau : chaque élève justifie son choix en citant le schéma de structure (\zh{A 跟 B 一样} + adjectif final). L'enseignant fait vérifier les quatre points critiques de la grille ci-dessous.

\begin{popfigure}
\begin{tikzpicture}[font=\small]
\draw[fill=popcream, rounded corners=4pt] (0,0) rectangle (11.5,4.6);
\node[font=\bfseries, text=popdark] at (5.75,4.2) {Grille d'auto-évaluation de ma phrase comparative};
\node[draw=popblue, minimum size=0.35cm] at (0.6,3.3) {};
\node[anchor=west] at (1.1,3.3) {Le mot-outil de comparaison est présent (跟 / 和 / 与)};
\node[draw=popblue, minimum size=0.35cm] at (0.6,2.4) {};
\node[anchor=west] at (1.1,2.4) {一样 / 相同 est bien placé juste après B};
\node[draw=popblue, minimum size=0.35cm] at (0.6,1.5) {};
\node[anchor=west] at (1.1,1.5) {La négation 不 est devant 一样 / 相同 (pas devant 跟)};
\node[draw=popblue, minimum size=0.35cm] at (0.6,0.6) {};
\node[anchor=west] at (1.1,0.6) {L'adjectif est en fin de phrase, sans 是 ni 很};
\end{tikzpicture}
\legende{Grille d'auto-évaluation : quatre cases à cocher avant de rendre sa copie.}
\end{popfigure}

\begin{popfigure}
\begin{tikzpicture}[font=\small]
\draw[fill=popgreenL, rounded corners=4pt] (0,0) rectangle (11.5,3.4);
\node[font=\bfseries, text=popdark] at (5.75,3.0) {Barème visuel type Bac : phrase de traduction (4 points)};
\foreach \x/\t in {0.3/{1 pt : A + 跟 + B}, 3.2/{1 pt : 一样 / 相同}, 6.1/{1 pt : négation 不 correcte}, 9.0/{1 pt : adjectif final}} {
  \draw[fill=white, draw=popgreen, rounded corners=2pt] (\x,0.5) rectangle (\x+2.3,2.2);
  \node[align=center, text width=2.1cm] at (\x+1.15,1.35) {\t};
}
\end{tikzpicture}
\legende{Un point par élément correctement placé : cadre, ancre, négation, adjectif.}
\end{popfigure}

\begin{aretenir}[Bilan de la leçon]
Pour toute phrase comparative : (1) poser le cadre \zh{A 跟/和/与 B 一样/相同} ; (2) nier avec \zh{不一样} ou \zh{不同} ; (3) placer l'adjectif en toute fin ; (4) réserver \zh{与} à l'écrit. Avec la méthode de l'ancre et la grille d'auto-évaluation, les exercices de transformation et de traduction deviennent des points faciles à l'examen.
\end{aretenir}

\newpage

\coursec{Activité d'intégration}

\courssub{Objectifs de l'activité}

À l'issue de cette activité d'intégration, tu seras capable de :
\begin{itemize}
\item utiliser correctement la structure d'égalité \cle{A 跟/和/与 B 一样/相同 (+ adjectif)} pour dire que deux réalités se ressemblent ;
\item utiliser correctement la structure de différence \cle{A 跟/和/与 B 不一样/不同} pour opposer deux réalités ;
\item employer le lexique de la ville, de la vie quotidienne et de la comparaison (météo, transports, nourriture, population, rythme de vie) ;
\item produire par écrit une fiche comparative de six phrases grammaticalement exactes, en caractères, avec pinyin et traduction ;
\item présenter oralement, en deux minutes, une comparaison claire et structurée devant la classe ;
\item évaluer la production de tes camarades à l'aide d'une grille d'auto-évaluation portant sur la correction des structures.
\end{itemize}

\courssub{Situation}

Dans le cadre de la semaine de la francophonie et des échanges culturels, ton lycée de Yaoundé prépare une exposition intitulée « Deux villes, deux mondes ». Des élèves d'une école partenaire de Pékin participeront en visioconférence. Le professeur de chinois demande à chaque binôme de préparer une fiche comparative entre deux villes : par exemple Douala et Yaoundé, ou ta ville et une ville chinoise. La fiche sera affichée au mur de la classe, et les meilleures seront envoyées aux partenaires chinois. Pour t'inspirer, le professeur vous distribue le petit texte suivant, écrit par une ancienne élève du lycée qui a visité les deux capitales :

\begin{textebox}[Texte support : Deux capitales]
雅温得和北京都是首都，但是这两个城市很不一样。\\
Yǎwēndé hé Běijīng dōu shì shǒudū, dànshì zhè liǎng ge chéngshì hěn bù yíyàng.\\
« Yaoundé et Pékin sont toutes deux des capitales, mais ces deux villes sont très différentes. »\\[4pt]
北京的人口跟雅温得的不一样：北京比雅温得大得多，人也多得多。\\
Běijīng de rénkǒu gēn Yǎwēndé de bù yíyàng : Běijīng bǐ Yǎwēndé dà de duō, rén yě duō de duō.\\
« La population de Pékin n'est pas la même que celle de Yaoundé : Pékin est beaucoup plus grande que Yaoundé, et il y a aussi beaucoup plus de monde. »\\[4pt]
北京的冬天很冷，雅温得的天气跟北京不一样，一年都很热。\\
Běijīng de dōngtiān hěn lěng, Yǎwēndé de tiānqì gēn Běijīng bù yíyàng, yì nián dōu hěn rè.\\
« L'hiver à Pékin est très froid ; le temps à Yaoundé n'est pas le même qu'à Pékin, il fait chaud toute l'année. »\\[4pt]
不过，两个城市也有相同的地方：北京的交通跟雅温得的一样忙，早上和晚上车都很多。\\
Bùguò, liǎng ge chéngshì yě yǒu xiāngtóng de dìfang : Běijīng de jiāotōng gēn Yǎwēndé de yíyàng máng, zǎoshang hé wǎnshang chē dōu hěn duō.\\
« Cependant, les deux villes ont aussi des points communs : la circulation à Pékin est aussi chargée qu'à Yaoundé, il y a beaucoup de voitures le matin et le soir. »\\[4pt]
两个城市的人都一样友好，我很喜欢这两个地方。\\
Liǎng ge chéngshì de rén dōu yíyàng yǒuhǎo, wǒ hěn xǐhuan zhè liǎng ge dìfang.\\
« Les habitants des deux villes sont tout aussi accueillants, j'aime beaucoup ces deux endroits. »
\end{textebox}

\begin{propriete}[Rappel des deux structures de la leçon]
\textbf{1. Égalité} : \cle{A + 跟/和/与 + B + 一样/相同 (+ adjectif)}.\\
L'adjectif, s'il est présent, se place \emph{après} \zh{一样}. La structure ne contient jamais le verbe \zh{是}.\\
\zh{杜阿拉跟雅温得一样大。} (Dùālā gēn Yǎwēndé yíyàng dà.) « Douala est aussi grande que Yaoundé. »\\[4pt]
\textbf{2. Différence} : \cle{A + 跟/和/与 + B + 不一样/不同 (+ adjectif)}.\\
La négation \zh{不} se place \emph{avant} \zh{一样} (ou devant \zh{同} : \zh{不同}).\\
\zh{雅温得的天气跟北京不一样。} (Yǎwēndé de tiānqì gēn Běijīng bù yíyàng.) « Le temps à Yaoundé n'est pas le même qu'à Pékin. »\\[4pt]
\textbf{3. Variante développée} : \cle{A 跟 B 不一样，A 很 X，B 很 Y} — après la constatation de la différence, une proposition précise le contraste.\\
\zh{杜阿拉跟雅温得不一样：杜阿拉很热，雅温得比较凉快。} « Douala n'est pas comme Yaoundé : il fait très chaud à Douala, alors qu'il fait plutôt frais à Yaoundé. »
\end{propriete}

\courssub{Consignes}

Travail en binômes. Suis les étapes dans l'ordre.

\begin{enumerate}
\item \textbf{Lecture et repérage (10 min).} Lis le texte support à voix haute avec ton camarade. Souligne en bleu toutes les structures d'égalité (\zh{跟……一样} ou \zh{相同}) et en rouge toutes les structures de différence (\zh{跟……不一样} ou \zh{不同}). Combien y en a-t-il de chaque type ? Recopie-les dans ton cahier avec leur traduction.
\item \textbf{Choix des villes et collecte du lexique (10 min).} Choisissez deux villes que vous connaissez bien : Douala et Yaoundé, Bafoussam et Buea, ou votre ville et Pékin (ou Shanghai). Complétez ensemble une liste de huit mots utiles à la comparaison (météo, population, nourriture, transports, prix, rythme de vie, langues, habitants). Vérifiez les caractères et le pinyin dans le tableau de vocabulaire de la leçon.
\item \textbf{Rédaction de la fiche comparative (20 min).} Rédigez six phrases complètes :
\begin{itemize}
\item trois phrases d'égalité avec \zh{跟/和/与……一样} + adjectif (ou \zh{相同}) ;
\item trois phrases de différence avec \zh{跟/和/与……不一样} (ou \zh{不同}), dont au moins une développée avec un adjectif précis : \zh{A 跟 B 不一样，A 很……，B 很……}.
\end{itemize}
Chaque phrase doit être écrite en caractères, avec le pinyin au-dessus ou en dessous, et la traduction française.
\item \textbf{Relecture croisée (10 min).} Échangez votre fiche avec le binôme voisin. Vérifiez chaque phrase avec la grille d'auto-évaluation : (a) l'ordre A + \zh{跟} + B est-il respecté ? (b) la négation \zh{不} est-elle bien placée \emph{avant} \zh{一样} ? (c) l'adjectif suit-il bien \zh{一样} ? (d) y a-t-il un \zh{是} interdit avant l'adjectif ? Corrigez ensemble les erreurs trouvées.
\item \textbf{Présentation orale (2 min par binôme).} Présentez votre comparaison à la classe. Pendant les présentations des autres, notez sur ton cahier une phrase d'égalité et une phrase de différence entendues, puis vérifie leur correction avec la grille.
\item \textbf{Affichage.} La classe vote pour la meilleure fiche selon deux critères : exactitude grammaticale et richesse du vocabulaire. La fiche gagnante est affichée au mur et envoyée aux partenaires de Pékin.
\end{enumerate}

\begin{methode}[Réussir sa présentation orale en quatre étapes]
\begin{enumerate}
\item \textbf{Annoncer} les deux villes comparées : \zh{我们要比较杜阿拉和雅温得。} (Wǒmen yào bǐjiào Dùālā hé Yǎwēndé.) « Nous allons comparer Douala et Yaoundé. »
\item \textbf{Présenter d'abord les ressemblances} avec \zh{一样} / \zh{相同}, en articulant lentement les tons.
\item \textbf{Présenter ensuite les différences} avec \zh{不一样} / \zh{不同}, en développant au moins une avec deux propositions (\zh{A 很 X，B 很 Y}).
\item \textbf{Conclure} par une phrase d'opinion : \zh{我很喜欢这两个城市。} (Wǒ hěn xǐhuan zhè liǎng ge chéngshì.) « J'aime beaucoup ces deux villes. »
\end{enumerate}
\end{methode}

\begin{attention}[Pièges à éviter dans ta fiche]
\begin{itemize}
\item Ne dis jamais \zh{雅温得是跟杜阿拉一样热} : pas de \zh{是} dans la structure d'égalité. Dis \zh{雅温得跟杜阿拉一样热}。Le français « est aussi chaude que » pousse à traduire « est » par \zh{是} : c'est un faux ami de structure, car en chinois l'adjectif seul joue le rôle de prédicat.
\item La négation porte sur \zh{一样}, pas sur l'adjectif : \zh{跟……不一样大} signifie « pas aussi grand que », alors que \zh{跟……一样不大} est une faute de sens.
\item L'ordre est figé : A + \zh{跟} + B + \zh{一样} + adjectif. Ne mets jamais l'adjectif avant \zh{跟} comme en français (« aussi grand que ») : le français place l'adjectif entre les deux termes, le chinois le place à la fin.
\item \zh{与} est plus écrit et plus formel que \zh{跟} et \zh{和} : réserve-le aux phrases de la fiche écrite, pas nécessairement à l'oral.
\item Ne confonds pas \zh{一样} (comparaison d'égalité) avec \zh{比} (comparaison de supériorité) : \zh{北京比雅温得大} « Pékin est plus grande que Yaoundé » ; \zh{北京跟雅温得一样大} « Pékin est aussi grande que Yaoundé ». Les deux structures ne se mélangent pas : on ne dit pas \zh{北京跟雅温得比一样大}。
\end{itemize}
\end{attention}

\courssub{Ressources à mobiliser}

\begin{aretenir}[Structures à réutiliser]
\begin{itemize}
\item Égalité : \cle{A + 跟/和/与 + B + 一样/相同 (+ adjectif)}. Exemple : \zh{杜阿拉的天气跟雅温得的一样热。} (Dùālā de tiānqì gēn Yǎwēndé de yíyàng rè.) « Il fait aussi chaud à Douala qu'à Yaoundé. »
\item Différence : \cle{A + 跟/和/与 + B + 不一样/不同}. Exemple : \zh{雅温得的天气跟北京不一样。} (Yǎwēndé de tiānqì gēn Běijīng bù yíyàng.) « Le temps à Yaoundé n'est pas le même qu'à Pékin. »
\item Différence développée : \cle{A 跟 B 不一样，A 很 X，B 很 Y}.
\end{itemize}
\end{aretenir}

\begin{center}
\begin{tabularx}{\linewidth}{|X|X|X|}
\hline
\rowcolor{popblueL} Structure & Exemple (caractères + pinyin) & Traduction \\
\hline
A 跟 B 一样 + adj. & \zh{杜阿拉跟雅温得一样大。} Dùālā gēn Yǎwēndé yíyàng dà. & Douala est aussi grande que Yaoundé. \\
\hline
A 和 B 相同 & \zh{两个城市的交通相同。} Liǎng ge chéngshì de jiāotōng xiāngtóng. & La circulation est identique dans les deux villes. \\
\hline
A 跟 B 不一样 & \zh{雅温得的天气跟北京不一样。} Yǎwēndé de tiānqì gēn Běijīng bù yíyàng. & Le temps à Yaoundé n'est pas le même qu'à Pékin. \\
\hline
A 与 B 不同 & \zh{杜阿拉与雅温得的天气不同。} Dùālā yǔ Yǎwēndé de tiānqì bùtóng. & Le temps à Douala est différent de celui de Yaoundé. \\
\hline
A 跟 B 不一样 + adj. & \zh{杜阿拉的东西跟雅温得的不一样贵。} Dùālā de dōngxi gēn Yǎwēndé de bù yíyàng guì. & Les choses à Douala ne sont pas aussi chères qu'à Yaoundé. \\
\hline
\end{tabularx}
\end{center}

\begin{center}
\begin{tabularx}{\linewidth}{|X|X|X|X|}
\hline
\rowcolor{popblueL} Caractères & Pinyin & Nature & Traduction \\
\hline
城市 & chéngshì & nom & ville \\
\hline
首都 & shǒudū & nom & capitale \\
\hline
人口 & rénkǒu & nom & population \\
\hline
天气 & tiānqì & nom & temps (météo) \\
\hline
交通 & jiāotōng & nom & circulation, transports \\
\hline
生活 & shēnghuó & nom & vie (quotidienne) \\
\hline
东西 & dōngxi & nom & chose, affaire \\
\hline
一样 & yíyàng & adj./adv. & pareil, identique \\
\hline
相同 & xiāngtóng & adj. & identique, semblable \\
\hline
不同 & bùtóng & adj. & différent \\
\hline
比较 & bǐjiào & adv./v. & plutôt, relativement ; comparer \\
\hline
友好 & yǒuhǎo & adj. & amical, accueillant \\
\hline
忙 & máng & adj. & chargé, animé, occupé \\
\hline
贵 & guì & adj. & cher \\
\hline
凉快 & liángkuai & adj. & frais (agréablement) \\
\hline
安静 & ānjìng & adj. & calme, tranquille \\
\hline
干净 & gānjìng & adj. & propre \\
\hline
\end{tabularx}
\end{center}

\begin{savaistu}[跟, 和, 与 : trois mots, trois registres]
Les trois mots \zh{跟}, \zh{和} et \zh{与} signifient tous « et / avec » et peuvent introduire le second terme de la comparaison. Mais ils ne s'emploient pas dans les mêmes contextes : \zh{跟} est le plus courant à l'oral dans le nord de la Chine (Pékin), \zh{和} est neutre et partout accepté, \zh{与} appartient à la langue écrite et formelle (articles, discours, documents officiels). À l'examen, dans une traduction vers le chinois, \zh{跟} et \zh{和} sont toujours corrects ; dans un texte écrit soutenu, \zh{与} donne un ton plus élégant. Retiens aussi que \zh{跟} et \zh{和} peuvent être conjonctions (« et ») : \zh{雅温得和北京都是首都} « Yaoundé et Pékin sont toutes deux des capitales ».
\end{savaistu}

\courssub{Proposition de production et corrigé commenté}

\begin{exoresolu}[Fiche comparative modèle : Douala et Yaoundé]
Énoncé : rédige une fiche de six phrases (trois d'égalité, trois de différence) comparant Douala et Yaoundé, avec pinyin et traduction, puis présente-la oralement en deux minutes.
\tcblower
\textbf{Production modèle.}

Égalité :
\begin{itemize}
\item \zh{杜阿拉跟雅温得一样大。} (Dùālā gēn Yǎwēndé yíyàng dà.) « Douala est aussi grande que Yaoundé. »
\item \zh{杜阿拉的人跟雅温得的人一样友好。} (Dùālā de rén gēn Yǎwēndé de rén yíyàng yǒuhǎo.) « Les gens de Douala sont aussi accueillants que ceux de Yaoundé. »
\item \zh{两个城市的交通相同：早上车都很多。} (Liǎng ge chéngshì de jiāotōng xiāngtóng : zǎoshang chē dōu hěn duō.) « La circulation est identique dans les deux villes : le matin, il y a beaucoup de voitures. »
\end{itemize}

Différence :
\begin{itemize}
\item \zh{杜阿拉的天气跟雅温得不一样：杜阿拉很热，雅温得比较凉快。} (Dùālā de tiānqì gēn Yǎwēndé bù yíyàng : Dùālā hěn rè, Yǎwēndé bǐjiào liángkuai.) « Le temps à Douala n'est pas le même qu'à Yaoundé : il fait très chaud à Douala, alors qu'il fait plutôt frais à Yaoundé. »
\item \zh{雅温得的生活跟杜阿拉不同：雅温得比较安静。} (Yǎwēndé de shēnghuó gēn Dùālā bùtóng : Yǎwēndé bǐjiào ānjìng.) « La vie à Yaoundé est différente de celle de Douala : Yaoundé est plus calme. »
\item \zh{杜阿拉的东西跟雅温得的不一样贵。} (Dùālā de dōngxi gēn Yǎwēndé de bù yíyàng guì.) « Les choses à Douala ne sont pas aussi chères qu'à Yaoundé. »
\end{itemize}

\textbf{Commentaire des choix de langue.}
\begin{enumerate}
\item Dans les phrases 1 et 2, l'adjectif (\zh{大}, \zh{友好}) est placé \emph{après} \zh{一样}, conformément au schéma A + \zh{跟} + B + \zh{一样} + adjectif ; aucun \zh{是} n'est utilisé.
\item La phrase 3 montre l'emploi de \zh{相同} sans adjectif après : la comparaison porte sur le nom \zh{交通} ; la précision est donnée dans une seconde proposition introduite par les deux-points chinois.
\item Les phrases 4 et 5 illustrent la différence développée : après \zh{不一样} / \zh{不同}, une proposition explicative précise le contraste avec \zh{很} ou \zh{比较} + adjectif, ce qui enrichit la fiche sans sortir du niveau HSK 3-4.
\item La phrase 6 montre la négation correctement placée : \zh{不一样贵} = « pas aussi cher que ». Attention à ne pas écrire \zh{一样不贵}, qui changerait le sens.
\item À l'oral, on préfère \zh{跟} ou \zh{和} ; à l'écrit, on peut varier avec \zh{与} : \zh{杜阿拉与雅温得的天气不同}。
\end{enumerate}
\end{exoresolu}

\begin{exercice}[Entraînement à l'examen : traduction et transformation]
\begin{enumerate}
\item Traduis en chinois (caractères + pinyin) : « Le marché de Bafoussam est aussi animé que le marché de Douala. »
\item Traduis en chinois : « La nourriture chinoise n'est pas la même que la nourriture camerounaise. »
\item Transforme en phrase de différence développée : \zh{北京的冬天跟雅温得的冬天不一样。} (précise : froid / chaud).
\item Corrige la phrase fautive : \zh{雅温得是跟杜阿拉一样大。}
\item Traduis en français : \zh{我的学校跟你的学校不一样大。} (Wǒ de xuéxiào gēn nǐ de xuéxiào bù yíyàng dà.)
\end{enumerate}
\end{exercice}

\begin{corrige}[Entraînement à l'examen]
\begin{enumerate}
\item \zh{巴富萨姆的市场跟杜阿拉的市场一样热闹。} (Bāfùsàmǔ de shìchǎng gēn Dùālā de shìchǎng yíyàng rènao.) Structure : A + \zh{跟} + B + \zh{一样} + adjectif \zh{热闹} « animé » ; pas de \zh{是}.
\item \zh{中国菜跟喀麦隆菜不一样。} (Zhōngguó cài gēn Kāmàilóng cài bù yíyàng.) On peut aussi dire \zh{中国菜与喀麦隆菜不同} en style écrit.
\item \zh{北京的冬天跟雅温得的冬天不一样：北京很冷，雅温得很热。} (Běijīng de dōngtiān gēn Yǎwēndé de dōngtiān bù yíyàng : Běijīng hěn lěng, Yǎwēndé hěn rè.) La deuxième proposition précise le contraste avec \zh{很} + adjectif.
\item \zh{雅温得跟杜阿拉一样大。} (Yǎwēndé gēn Dùālā yíyàng dà.) Il faut supprimer le \zh{是}, interdit dans la structure d'égalité : l'adjectif \zh{大} suffit comme prédicat.
\item « Mon école n'est pas aussi grande que ton école. » La négation \zh{不} devant \zh{一样} porte sur toute la comparaison : \zh{不一样大} = « pas de la même taille ».
\end{enumerate}
\end{corrige}

\courssub{Bilan de l'activité}

Cette activité t'a permis de mobiliser en situation réelle les deux grandes structures de la leçon : l'égalité \zh{跟/和/与……一样/相同} et la différence \zh{跟/和/与……不一样/不同}. Tu as d'abord \emph{repéré} ces structures dans un texte authentique (compréhension de l'écrit), puis tu les as \emph{réutilisées} dans une fiche écrite (expression écrite, caractères et pinyin), avant de les \emph{présenter oralement} (expression orale) et de les \emph{vérifier} chez tes camarades grâce à la grille d'auto-évaluation (grammaire réflexive). Retiens les trois points critiques : l'ordre figé A + \zh{跟} + B + \zh{一样} + adjectif, l'absence de \zh{是} dans la structure, et la place de la négation \zh{不} devant \zh{一样}. Ces structures, très fréquentes à l'examen en traduction (thème et version) comme en production écrite, te serviront aussi dans toutes les comparaisons Cameroun–Chine que tu seras amené à formuler : villes, écoles, cuisines, fêtes, rythmes de vie.

\newpage

\coursec{Méthodes et exercices résolus}

\coursec{Leçon 2 — Grammaire : Le comparatif d'égalité et de différence}

\courssub{Observation guidée : comparer deux réalités camerounaises}

\begin{textebox}[Dialogue : Deux lycéens comparent Yaoundé et Douala]
\zh{— 小张，你觉得雅温得跟杜阿拉一样吗？} \\
\emph{— Xiǎo Zhāng, nǐ juéde Yǎwēndé gēn Dù'ālā yíyàng ma?} \\
\zh{— 不完全一样。杜阿拉跟雅温得不一样热，杜阿拉更热。可是，两个城市的人都一样喜欢足球。} \\
\emph{— Bù wánquán yíyàng. Dù'ālā gēn Yǎwēndé bù yíyàng rè, Dù'ālā gèng rè. Kěshì, liǎng ge chéngshì de rén dōu yíyàng xǐhuan zúqiú.} \\
\zh{— 真的？我觉得交通差不多一样乱。} \\
\emph{— Zhēnde? Wǒ juéde jiāotōng chàbuduō yíyàng luàn.}
\tcblower
\emph{— « Petit Zhang, tu trouves que Yaoundé et Douala se ressemblent ? »} \\
\emph{— « Pas tout à fait. Douala n'a pas la même chaleur que Yaoundé, Douala est plus chaud. Mais les gens des deux villes aiment le football de la même façon. »} \\
\emph{— « Vraiment ? Moi je trouve la circulation à peu près aussi chaotique. »}
\end{textebox}

\begin{propriete}[Vocabulaire clé de la leçon]
\begin{center}
\begin{tabularx}{\linewidth}{|X|X|X|X|}
\hline
\rowcolor{popblueL} \textbf{Caractères} & \textbf{Pinyin} & \textbf{Nature} & \textbf{Traduction} \\ \hline
比较 & bǐjiào & adv. & relativement, assez \\ \hline
跟 & gēn & prep. & avec (oral), par rapport à \\ \hline
和 & hé & prep. & et, avec (neutre) \\ \hline
与 & yǔ & prep. & avec, et (écrit, formel) \\ \hline
一样 & yíyàng & adj. & pareil, semblable, identique \\ \hline
相同 & xiāngtóng & adj. & identique (formel, abstrait) \\ \hline
不一样 & bù yíyàng & adj. & pas pareil, différent \\ \hline
不同 & bùtóng & adj. & différent (formel) \\ \hline
完全 & wánquán & adv. & complètement, tout à fait \\ \hline
差不多 & chàbuduō & adv. & à peu près, presque \\ \hline
不太 & bù tài & adv. & pas très, pas tout à fait \\ \hline
气候 & qìhòu & n. & climat \\ \hline
交通 & jiāotōng & n. & circulation, transports \\ \hline
足球 & zúqiú & n. & football \\ \hline
乱 & luàn & adj. & en désordre, chaotique \\ \hline
觉得 & juéde & v. & trouver, penser que \\ \hline
\end{tabularx}
\end{center}
\end{propriete}

\courssub{Structure d'égalité : A 跟/和/与 B 一样/相同}

\begin{propriete}[Schéma fondamental de l'égalité]
\textbf{Structure de base :} \zh{A + 跟/和/与 + B + 一样/相同} \\
\textbf{Structure étendue (point de comparaison) :} \zh{A + 跟/和/与 + B + 一样/相同 + Adjectif / Verbe} \\
\begin{itemize}
\item \zh{跟} (gēn) : registre oral, familier, le plus fréquent à l'oral.
\item \zh{和} (hé) : registre neutre, passe-partout, sûr à l'écrit comme à l'oral.
\item \zh{与} (yǔ) : registre soutenu, écrit formel (presse, administration, dissertation).
\item \zh{一样} (yíyàng) : « pareil, semblable » ; usage général, accepte un adjectif/verbe derrière.
\item \zh{相同} (xiāngtóng) : « identique » ; registre soutenu, sens plus abstrait, rarement suivi d'adjectif.
\end{itemize}
\begin{popfigure}
\begin{tikzpicture}[node distance=0.8cm and 1.2cm, every node/.style={font=\small}]
\node (A) [draw, rounded corners, fill=popblueL] {A (Sujet)};
\node (prep) [right=of A, draw, diamond, aspect=2, fill=poporangeL] {跟 / 和 / 与};
\node (B) [right=of prep, draw, rounded corners, fill=popblueL] {B (Terme de comparaison)};
\node (pred) [right=of B, draw, rounded corners, fill=popgreenL] {一样 / 相同};
\node (adj) [right=of pred, draw, rounded corners, fill=poppinkL] {(Adj. / Verbe)};
\draw[-Stealth, thick, popdark] (A) -- (prep);
\draw[-Stealth, thick, popdark] (prep) -- (B);
\draw[-Stealth, thick, popdark] (B) -- (pred);
\draw[-Stealth, thick, popdark] (pred) -- (adj);
\node[below=0.3cm of prep, align=center, text width=2cm] {\textbf{Mot-outil}\\(choix registral)};
\node[below=0.3cm of pred, align=center, text width=2cm] {\textbf{Prédicat}\\(pas de \zh{是}, pas de \zh{很})};
\end{tikzpicture}
\legende{Schéma de l'ordre des mots : le bloc prépositionnel précède obligatoirement le prédicat.}
\end{popfigure}
\end{propriete}

\begin{methode}[Construire une phrase d'égalité en 5 étapes]
\begin{enumerate}
\item \cle{Identifier} les deux termes comparés : A (sujet) et B (terme de référence). Ce sont des noms, pronoms ou groupes nominaux.
\item \cle{Choisir} le mot-outil de liaison selon le registre : \zh{跟} (oral), \zh{和} (neutre), \zh{与} (écrit formel).
\item \cle{Placer} le bloc \zh{跟/和/与 + B} \textbf{immédiatement après le sujet A}, avant le prédicat. C'est un complément prépositionnel, jamais un objet direct.
\item \cle{Conclure} par le prédicat : \zh{一样} (général) ou \zh{相同} (formel/abstrait). \textbf{Interdiction absolue} : pas de \zh{是} (shì), pas de \zh{很} (hěn) devant \zh{一样}.
\item \cle{Préciser} (optionnel) le point de comparaison en ajoutant un adjectif (\zh{大}, \zh{高}, \zh{贵}) ou un verbe (\zh{喜欢}, \zh{会}) après \zh{一样}.
\end{enumerate}
\cle{Règle d'or :} En français « A est aussi adj. que B » $\rightarrow$ En chinois \zh{A 跟 B 一样 + Adj.} (pas de mot pour « aussi »).
\end{methode}

\begin{exemplebox}[Exemples d'égalité simples et étendues]
\begin{itemize}
\item \zh{我的手机跟你的一样。} \emph{Wǒ de shǒujī gēn nǐ de yíyàng.} (Mon téléphone est pareil au tien.) — \textit{Prédicat seul.}
\item \zh{这件衣服和那件一样贵。} \emph{Zhè jiàn yīfu hé nà jiàn yíyàng guì.} (Ce vêtement est aussi cher que celui-là.) — \textit{Avec adjectif.}
\item \zh{哥哥与弟弟相同。} \emph{Gēge yǔ dìdi xiāngtóng.} (Le grand frère et le petit frère sont identiques.) — \textit{Registre formel, \zh{相同} sans adjectif.}
\item \zh{喀麦隆学生跟中国学生一样努力。} \emph{Kāmàilóng xuésheng gēn Zhōngguó xuésheng yíyàng nǔlì.} (Les élèves camerounais sont aussi travailleurs que les élèves chinois.) — \textit{Contexte socioculturel.}
\end{itemize}
\end{exemplebox}

\begin{exoresolu}[Transformation : fusionner deux phrases en une égalité]
Énoncé — Transforme en une phrase avec \zh{一样} : \\
\zh{我妹妹十八岁。我十八岁。} (Wǒ mèimei shíbā suì. Wǒ shíbā suì.)
\tcblower
Solution détaillée :
\begin{enumerate}
\item Termes : A = \zh{我妹妹}, B = \zh{我}. Point de comparaison : l'âge $\rightarrow$ adjectif \zh{大} (grand/âgé).
\item Mot-outil : \zh{跟} (contexte énonciatif simple).
\item Assemblage : Sujet A + \zh{跟} + B + \zh{一样} + Adj.
\end{enumerate}
Réponse : \zh{我妹妹跟我一样大。} \emph{Wǒ mèimei gēn wǒ yíyàng dà.} \\
\emph{Justification :} \zh{一样大} = « même âge ». Pas de \zh{是}, pas de \zh{很}. L'ordre \zh{跟我} avant \zh{一样大} est immuable.
\end{exoresolu}

\courssub{Structure de différence : A 跟/和/与 B 不一样/不同}

\begin{propriete}[Schéma fondamental de la différence]
\textbf{Structure de base :} \zh{A + 跟/和/与 + B + 不一样/不同} \\
\textbf{Structure étendue :} \zh{A + 跟/和/与 + B + 不一样/不同 + Adjectif} \\
\begin{itemize}
\item Négation : \zh{不} (bù) se place \textbf{uniquement} devant \zh{一样} $\rightarrow$ \zh{不一样} (bù yíyàng).
\item \cle{Sandhi tonal :} \zh{不} (4e ton) $\rightarrow$ \zh{不} (2e ton, \emph{bú}) devant un 4e ton (\zh{一} yì $\rightarrow$ yí). Prononciation réelle : \emph{bú yíyàng}. À l'écrit, on note \zh{不一样}.
\item \zh{不同} (bùtóng) : forme contractée formelle de \zh{不相同}, s'utilise comme \zh{不一样} mais sans adjectif suivant généralement.
\item \textbf{Erreur interdite :} \zh{不跟……一样} (négation du mot-outil). On nie la ressemblance, pas la relation.
\end{itemize}
\end{propriete}

\begin{methode}[Construire une phrase de différence sans faute]
\begin{enumerate}
\item \cle{Repérer} A, B et le point de comparaison (prix, taille, climat, opinion…).
\item \cle{Former} le prédicat négatif : \zh{不一样} (oral/neutre) ou \zh{不同} (écrit).
\item \cle{Assembler} : \zh{A + 跟/和/与 + B + 不一样/不同 (+ Adj.)}.
\item \cle{Gérer l'ellipse} : si B est un nom « possesseur » de la qualité comparée, garder le \zh{的} de reprise. Ex: \zh{北京的天气跟上海的不一样。} (Celui de Shanghai).
\end{enumerate}
\cle{Test :} Si tu peux dire en français « A n'est pas le même que B » ou « A diffère de B », utilise cette structure.
\end{methode}

\begin{exemplebox}[Exemples de différence avec ellipse du \zh{的}]
\begin{itemize}
\item \zh{这道菜跟那道菜不一样辣。} \emph{Zhè dào cài gēn nà dào cài bù yíyàng là.} (Ce plat n'est pas aussi épicé que celui-là.)
\item \zh{雅温得的气候与杜阿拉的不同。} \emph{Yǎwēndé de qìhòu yǔ Dù'ālā de bùtóng.} (Le climat de Yaoundé diffère de celui de Douala.) — \textit{\zh{的} remplace \zh{气候}, registre formel \zh{与}/\zh{不同}.}
\item \zh{我的想法跟老师的不太一样。} \emph{Wǒ de xiǎngfa gēn lǎoshī de bù tài yíyàng.} (Mon avis n'est pas tout à fait le même que celui du prof.) — \textit{Nuance \zh{不太}.}
\end{itemize}
\end{exemplebox}

\begin{exoresolu}[Thème : traduire une phrase de différence contextuelle]
Énoncé — « Le prix du ndolé au marché n'est pas le même qu'au supermarché. » (ndolé \zh{恩多莱} ēnduōlái ; marché \zh{市场} shìchǎng ; supermarché \zh{超市} chāoshì ; prix \zh{价格} jiàgé).
\tcblower
Solution détaillée :
\begin{enumerate}
\item A = \zh{市场的恩多莱价格} (ou \zh{市场买恩多莱的价格}), B = \zh{超市}.
\item Structure différence + ellipse : \zh{市场的恩多莱价格跟超市的不一样。}
\item Alternative avec \zh{不同} (plus formel) : \zh{市场的恩多莱价格与超市的不同。}
\end{enumerate}
Réponse rédigée : \zh{市场的恩多莱价格跟超市的不一样。} \\
\emph{Shìchǎng de ēnduōlái jiàgé gēn chāoshì de bù yíyàng.} \\
\emph{Point de vigilance :} Le \zh{的} final après \zh{超市} est obligatoire pour nominaliser « celui du supermarché ». Sans lui, on compare un prix à un lieu.
\end{exoresolu}

\courssub{Registres, nuances et emplois étendus (niveau Bac)}

\begin{aretenir}[Tableau récapitulatif : Choix du mot-outil et du prédicat]
\begin{center}
\begin{tabularx}{\linewidth}{|X|X|X|X|}
\hline
\rowcolor{popblueL} \textbf{Contexte} & \textbf{Mot-outil} & \textbf{Prédicat égalité} & \textbf{Prédicat différence} \\ \hline
Dialogue, message, oral & \zh{跟} (gēn) & \zh{一样} (yíyàng) & \zh{不一样} (bù yíyàng) \\ \hline
Rédaction courante, neutre & \zh{和} (hé) & \zh{一样} / \zh{相同} & \zh{不一样} / \zh{不同} \\ \hline
Dissertation, article, officiel & \zh{与} (yǔ) & \zh{相同} (xiāngtóng) & \zh{不同} (bùtóng) \\ \hline
\end{tabularx}
\end{center}
\cle{Exception importante :} \zh{相同} ne se prolonge quasi jamais par un adjectif (*\zh{相同大} faux). Pour « identiquement grand », on dit \zh{一样大}.
\end{aretenir}

\begin{methode}[Nuancer la comparaison : adverbes de degré]
Les adverbes se placent \textbf{immédiatement avant} \zh{一样} ou \zh{不一样} (jamais avant le mot-outil \zh{跟}).
\begin{center}
\begin{tabularx}{\linewidth}{|X|X|X|}
\hline
\rowcolor{popblueL} \textbf{Adverbe} & \textbf{Structure} & \textbf{Sens / Nuance} \\ \hline
\zh{完全} wánquán & A 跟 B \zh{完全一样} & Tout à fait pareil (renforcement) \\ \hline
\zh{真的} zhēnde & A 跟 B \zh{真的一样} & Vraiment pareil (insistance) \\ \hline
\zh{差不多} chàbuduō & A 跟 B \zh{差不多一样} & À peu près pareil (approximation) \\ \hline
\zh{不太} bù tài & A 跟 B \zh{不太一样} & Pas tout à fait pareil (atténuation polie) \\ \hline
\zh{完全不} wánquán bù & A 跟 B \zh{完全不一样} & Complètement différent (renforcement nég.) \\ \hline
\end{tabularx}
\end{center}
\cle{Interrogation :} \zh{A 跟 B 一样吗？} $\rightarrow$ Réponse : \zh{一样。} / \zh{不一样。} / \zh{差不多。}
\cle{Prédicat verbal :} \zh{A 跟 B 一样 + Verbe.} Ex: \zh{他跟我一样喜欢看电影。} (Il aime regarder des films comme moi.)
\end{methode}

\begin{exoresolu}[Composition guidée : paragraphe comparatif varié]
Énoncé — Rédige 5 phrases comparant ton lycée et un lycée chinois imaginaire, en utilisant : 1× \zh{一样}+Adj, 1× \zh{不一样}+Adj, 1× \zh{差不多一样}, 1× \zh{与……不同}, 1× \zh{完全一样}.
\tcblower
Modèle de réponse :
\begin{enumerate}
\item \zh{我们学校跟北京那所中学一样大。} (Égalité + Adj)
\item \zh{我们的食堂饭菜跟他们的不一样好吃，我们的更好吃！} (Différence + Adj + comparatif supériorité)
\item \zh{两个学校的上课时间差不多一样早。} (Nuance \zh{差不多})
\item \zh{我们的校服与他们的校服不同。} (Registre formel \zh{与}/\zh{不同})
\item \zh{但是，学生对高考的压力完全一样。} (Renforcement \zh{完全})
\end{enumerate}
\emph{Justification :} Démonstration de la maîtrise des registres (\zh{跟} vs \zh{与}), des prédicats (\zh{一样} vs \zh{不同}), des nuances (\zh{差不多}, \zh{完全}) et de la syntaxe stricte.
\end{exoresolu}

\courssub{Pièges francophones, erreurs types et stratégies de correction}

\begin{attention}[Les 6 erreurs fatales au Bac — Diagnostic et Correction]
\begin{enumerate}
\item \cle{Calque syntaxique français (ordre des mots)} \\
Faux : \zh{*我一样高跟他。} / \zh{*一样我跟他高。} \\
Vrai : \zh{我跟他一样高。} \\
\cle{Règle :} Le complément de comparaison (\zh{跟}+B) est un adverbial ; il précède le prédicat. \textbf{Sujet + 跟B + 一样 + Adj.}
\item \cle{Négation déplacée : \zh{*不跟……一样}} \\
Faux : \zh{*这本书不跟那本书一样。} \\
Vrai : \zh{这本书跟那本书不一样。} \\
\cle{Règle :} \zh{不} nie \zh{一样} (la ressemblance), pas la préposition \zh{跟}.
\item \cle{Verbe copule \zh{是} ou adverbe \zh{很} parasites} \\
Faux : \zh{*他是跟我一样大。} / \zh{*他很一样高。} \\
Vrai : \zh{他跟我一样大。} \\
\cle{Règle :} \zh{一样} est un adjectif prédicatif autonome. Seuls les adverbes de degré (\zh{完全}, \zh{差不多}, \zh{不太}) le précèdent.
\item \cle{Confusion graphique / phonétique} \\
\zh{跟} (gēn, pied 足) $\neq$ \zh{很} (hěn, très) ; \zh{与} (yǔ, 3 traits) $\neq$ \zh{兴} (xīng/xìng) ; \zh{相同} (xiāngtóng, 1er ton) $\neq$ \zh{想同} (xiǎng tóng). \zh{一} (yī) change de ton : yí (devant 4e), yì (devant 1/2/3), yi (neutre).
\item \cle{Oubli du \zh{的} de reprise (ellipse mal faite)} \\
Faux : \zh{*雅温得的气候跟杜阿拉不一样。} (Compare climat vs ville). \\
Vrai : \zh{雅温得的气候跟杜阿拉的不一样。} (\zh{的} = « celui de »).
\item \cle{Adjectif après \zh{相同}} \\
Faux : \zh{*这两本书相同厚。} \\
Vrai : \zh{这两本书一样厚。} / \zh{这两本书的厚度相同。} (Nominalisation).
\end{enumerate}
\end{attention}

\begin{methode}[Protocole de correction d'exercices type Bac (Transformation / Correction)]
\begin{enumerate}
\item \cle{Analyser} la phrase source : identifier A, B, le mot-outil, le prédicat (\zh{一样}/\zh{不一样}), l'adjectif éventuel.
\item \cle{Cibler} la consigne : Affirmatif $\leftrightarrow$ Négatif (ajouter/retirer \zh{不}) ; Fusion (2 phrases $\rightarrow$ 1) ; Correction (repérer l'erreur parmi les 6 ci-dessus).
\item \cle{Appliquer} la transformation minimale : ne toucher qu'à l'élément demandé (ex: ajouter \zh{不} devant \zh{一样}), sans réordonnancer la phrase.
\item \cle{Vérifier} la checklist : 1. Ordre A 跟 B ? 2. \zh{不} devant \zh{一样} ? 3. Pas de \zh{是}/\zh{很} ? 4. \zh{的} de reprise si ellipse ? 5. Registre cohérent (\zh{与} $\leftrightarrow$ \zh{不同}) ?
\item \cle{Relire} à voix haute (mentalement) en contrôlant les tons de \zh{一} et \zh{不}.
\end{enumerate}
\end{methode}

\begin{exoresolu}[Correction de phrases fautives (Type Bac)]
Énoncé — Corrige et justifie chaque phrase.
\begin{enumerate}
\item \zh{北京的冬天跟哈尔滨的一样不冷。}
\item \zh{我和我哥哥相同高。}
\item \zh{这双鞋不跟那双鞋一样贵。}
\item \zh{我的中文老师跟你的完全相同好。}
\end{enumerate}
\tcblower
Solution détaillée :
\begin{enumerate}
\item \textbf{Erreur :} Place de la négation / Sens. \zh{一样不冷} = « aussi pas froid » (égalité dans le non-froid). Si le sens est « pas aussi froid », structure différence. \\
\textbf{Correction :} \zh{北京的冬天跟哈尔滨的不一样冷。} (Pékin n'est pas aussi froid que Harbin.) Ou : \zh{北京的冬天不跟哈尔滨的一样冷} (FAUX, voir erreur 2). \textbf{Bonne correction :} \zh{北京的冬天跟哈尔滨的不太一样冷。} (Pas tout aussi froid).
\item \textbf{Erreur :} \zh{相同} + Adj interdit. \zh{和} neutre mais \zh{相同} formel $\rightarrow$ incohérence registrale. \\
\textbf{Correction :} \zh{我和我哥哥一样高。} (Ou \zh{我跟我哥哥一样高。})
\item \textbf{Erreur :} \zh{不跟} (négation mot-outil). \\
\textbf{Correction :} \zh{这双鞋跟那双鞋不一样贵。}
\item \textbf{Erreur :} \zh{相同} + Adj (\zh{好}) ; \zh{完全} modifie \zh{相同} mais structure fausse. \\
\textbf{Correction :} \zh{我的中文老师跟你的完全一样好。} (Registre oral/neutre \zh{跟} + \zh{一样} + Adj).
\end{enumerate}
\end{exoresolu}

\courssub{Entraînement : Exercices d'application}

\begin{exercice}[Exercices de la leçon — À rendre / Travailler en classe]
\textbf{Exercice 1 : Choix multiple (Registre).} Complète avec \zh{跟}, \zh{和} ou \zh{与} et justifie le registre.
\begin{enumerate}
\item (SMS à un ami) \zh{你的新书\_\_\_\_我的\_\_\_\_。} (一样)
\item (Rapport officiel) \zh{本年度预算\_\_\_\_去年\_\_\_\_。} (相同)
\item (Conversation) \zh{今天的天气\_\_\_\_昨天\_\_\_\_。} (不一样)
\end{enumerate}

\textbf{Exercice 2 : Transformation (Affirmatif $\rightarrow$ Négatif / Négatif $\rightarrow$ Affirmatif).}
\begin{enumerate}
\item \zh{这两个汉字写法一样。} $\rightarrow$ (Différence)
\item \zh{我的观点跟你的不一样。} $\rightarrow$ (Égalité)
\item \zh{杜阿拉的湿度跟雅温得的完全一样。} $\rightarrow$ (Différence nuancée : \zh{不太})
\end{enumerate}

\textbf{Exercice 3 : Correction d'erreurs (Une erreur par phrase).}
\begin{enumerate}
\item \zh{我跟他不一样高。} (Sens visé : « Je ne suis pas aussi grand que lui »)
\item \zh{这件衣服跟那件很一样贵。}
\item \zh{喀麦隆的足球与中国的足球不一样强。} (Registre)
\item \zh{妹妹跟姐姐一样大岁。} (Classificateur / Structure âge)
\end{enumerate}

\textbf{Exercice 4 : Thème (Français $\rightarrow$ Chinois).}
\begin{enumerate}
\item Ma ville natale a le même climat que Yaoundé.
\item Le prix du riz au supermarché n'est pas le même qu'au marché.
\item Les élèves camerounais et les élèves chinois aiment tous le football. (Utiliser \zh{一样喜欢})
\item Notre professeur de chinois est aussi strict que le vôtre. (strict \zh{严厉} yánlì)
\item Les coutumes de ce village sont à peu près identiques à celles de la ville. (coutume \zh{风俗} fēngsú ; village \zh{村庄} cūnzhuāng ; ville \zh{城市} chéngshì)
\end{enumerate}

\textbf{Exercice 5 : Expression écrite (Production).}
Rédige un court paragraphe (60-80 caractères) comparant \textbf{la vie étudiante à Yaoundé et à Pékin}. Utilise obligatoirement : une structure d'égalité avec adjectif, une structure de différence avec adjectif, un adverbe de nuance (\zh{差不多} ou \zh{不太} ou \zh{完全}), et le mot-outil \zh{与} au moins une fois. Soigne la ponctuation chinoise.
\end{exercice}

\begin{corrige}[Exercices de la leçon]
\textbf{Exercice 1 :}
1. \zh{跟} / \zh{一样} (Oral, SMS). 2. \zh{与} / \zh{相同} (Écrit, officiel, abstrait). 3. \zh{跟} / \zh{不一样} (Oral, conversation).

\textbf{Exercice 2 :}
1. \zh{这两个汉字写法不一样。} / \zh{这两个汉字写法不同。}
2. \zh{我的观点跟你的一样。}
3. \zh{杜阿拉的湿度跟雅温得的不太一样。} (ou \zh{差不多一样} pour « à peu près pareil » mais consigne = différence nuancée).

\textbf{Exercice 3 :}
1. \textbf{Erreur :} \zh{不一样高} = « pas de la même taille » (différence absolue). Pour « pas aussi grand » (infériorité), on utilise le comparatif \zh{比} : \zh{我比他不高} / \zh{我没有他高}. \textit{Note : Si le sens est bien « tailles différentes », \zh{不一样高} est correct. Ici, correction attendue pour « pas aussi grand » :} \zh{我没有他高。} (HSK 3/4).
2. \textbf{Erreur :} \zh{很一样}. \textbf{Correction :} \zh{这件衣服跟那件一样贵。}
3. \textbf{Erreur :} Registre mixte \zh{与} (formel) + \zh{不一样} (oral/neutre) + \zh{强} (adj après \zh{不一样} possible mais \zh{与} appelle \zh{不同}). \textbf{Correction :} \zh{喀麦隆的足球与中国的足球不同。} (Ou \zh{喀麦隆足球跟中国足球不一样强。} si on garde l'adj).
4. \textbf{Erreur :} \zh{大岁} n'existe pas. Âge = \zh{岁} seul ou \zh{年纪} niánjì. \textbf{Correction :} \zh{妹妹跟姐姐一样大。} / \zh{妹妹跟姐姐年纪一样。}

\textbf{Exercice 4 :}
1. \zh{我的家乡跟雅温得气候一样。} (ou \zh{一样热/冷} selon contexte).
2. \zh{超市的米价跟市场的不一样。} (\zh{米价} mǐjià prix du riz ; \zh{的} de reprise obligatoire).
3. \zh{喀麦隆学生跟中国学生一样喜欢足球。}
4. \zh{我们的中文老师跟你们的一样严厉。}
5. \zh{这个村庄的风俗与城市的差不多一样。} (ou \zh{与城市的风俗差不多相同}).

\textbf{Exercice 5 :} (Exemple de production attendue)
\zh{雅温得的学生生活与北京的不太一样：北京的作业完全一样多，可是雅温得的课外活动跟北京的一样丰富。两地学生差不多一样努力，因为高考压力相同。} \\
\emph{Yǎwēndé de xuésheng shēnghuó yǔ Běijīng de bù tài yíyàng: Běijīng de zuòyè wánquán yíyàng duō, kěshì Yǎwēndé de kèwài huódòng gēn Běijīng de yíyàng fēngfù. Liǎng dì xuésheng chàbuduō yíyàng nǔlì, yīnwèi gāokǎo yālì xiāngtóng.}
\end{corrige}

\begin{savaistu}[Le comparatif dans la culture chinoise : \zh{比} vs \zh{跟……一样}]
En chinois classique, la comparaison d'égalité s'exprimait souvent par \zh{若} (ruò, « comme, semblable à ») ou \zh{如} (rú). La structure moderne \zh{跟……一样} s'est imposée sous l'influence de la grammaire vernaculaire (báihuà) au XXe siècle.
\cle{Note culturelle :} Dans la politesse chinoise, on évite souvent la différence frontale (\zh{不一样}). On préfère \zh{不太一样} (pas tout à fait pareil) ou \zh{各有特色} (gè yǒu tèsè, « chacun a ses particularités ») pour sauvegarder la face (\zh{面子} miànzi) de l'interlocuteur. Au Cameroun, la franchise est souvent valorisée ; en Chine, l'harmonie (\zh{和谐} héxié) guide le choix linguistique.
\end{savaistu}

\newpage

\coursec{Exercices}

\courssub{Je vérifie mes connaissances}

\begin{exercice}[QCM : la bonne particule]
Choisis la bonne réponse pour compléter chaque phrase.
\begin{enumerate}
\item 这本书\_\_\_那本书相同。(\emph{Zhè běn shū \_\_\_ nà běn shū xiāngtóng.})\\
A. 很 \hspace{1cm} B. 与 \hspace{1cm} C. 也 \hspace{1cm} D. 得
\item 我\_\_\_我哥哥一样高。(\emph{Wǒ \_\_\_ wǒ gēge yīyàng gāo.})\\
A. 是 \hspace{1cm} B. 跟 \hspace{1cm} C. 很 \hspace{1cm} D. 比
\item 汉语\_\_\_法语不一样难。(\emph{Hànyǔ \_\_\_ Fǎyǔ bù yīyàng nán.})\\
A. 和 \hspace{1cm} B. 的 \hspace{1cm} C. 了 \hspace{1cm} D. 得
\item 杜阿拉的天气\_\_\_雅温得的天气不同。(\emph{Dù'ālā de tiānqì \_\_\_ Yǎwēndé de tiānqì bùtóng.})\\
A. 很 \hspace{1cm} B. 太 \hspace{1cm} C. 与 \hspace{1cm} D. 不
\end{enumerate}
\end{exercice}

\begin{exercice}[Vrai ou faux ? Justifie ta réponse]
Dis si chaque phrase est correcte (V) ou incorrecte (F), puis justifie en citant la règle de grammaire concernée.
\begin{enumerate}
\item \zh{我跟他一样高。} (\emph{Wǒ gēn tā yīyàng gāo.}) « Je suis aussi grand que lui. »
\item \zh{我一样跟他喜欢足球。} (\emph{Wǒ yīyàng gēn tā xǐhuan zúqiú.}) « J'aime le football comme lui. »
\item \zh{我的学校与他的学校相同。} (\emph{Wǒ de xuéxiào yǔ tā de xuéxiào xiāngtóng.}) « Mon école est identique à la sienne. »
\item \zh{她不跟我一样喜欢音乐。} (\emph{Tā bù gēn wǒ yīyàng xǐhuan yīnyuè.}) « Elle n'aime pas la musique autant que moi. »
\end{enumerate}
\end{exercice}

\begin{exercice}[Vocabulaire : complète le tableau]
Recopie et complète le tableau suivant en donnant le pinyin et la traduction française de chaque mot.
\begin{center}
\begin{tabularx}{\linewidth}{|X|X|X|X|}
\hline
\rowcolor{popblueL} Caractères & Pinyin & Nature & Traduction \\
\hline
一样 & \ldots\ldots & \ldots\ldots & \ldots\ldots \\
\hline
相同 & \ldots\ldots & \ldots\ldots & \ldots\ldots \\
\hline
不同 & \ldots\ldots & \ldots\ldots & \ldots\ldots \\
\hline
跟 & \ldots\ldots & \ldots\ldots & \ldots\ldots \\
\hline
与 & \ldots\ldots & \ldots\ldots & \ldots\ldots \\
\hline
比较 & \ldots\ldots & \ldots\ldots & \ldots\ldots \\
\hline
\end{tabularx}
\end{center}
\end{exercice}

\begin{exercice}[Associe le pinyin aux caractères]
Relie chaque expression en caractères à son pinyin.
\begin{enumerate}
\item \zh{跟……一样} \hspace{2cm} a. \emph{yǔ \ldots\ldots bùtóng}
\item \zh{和……不同} \hspace{2cm} b. \emph{gēn \ldots\ldots yīyàng}
\item \zh{与……相同} \hspace{2cm} c. \emph{hé \ldots\ldots bùtóng}
\item \zh{跟……不一样} \hspace{2cm} d. \emph{yǔ \ldots\ldots xiāngtóng}
\item \hspace{3cm} e. \emph{gēn \ldots\ldots bù yīyàng}
\end{enumerate}
\end{exercice}

\courssub{Je m'entraîne}

\begin{exercice}[Traduis en français]
Traduis les phrases suivantes en français.
\begin{enumerate}
\item \zh{喀麦隆的足球跟巴西的足球一样有名。} (\emph{Kāmàilóng de zúqiú gēn Bāxī de zúqiú yīyàng yǒumíng.})
\item \zh{我的想法与你的相同。} (\emph{Wǒ de xiǎngfǎ yǔ nǐ de xiāngtóng.})
\item \zh{今天的天气和昨天不一样热。} (\emph{Jīntiān de tiānqì hé zuótiān bù yīyàng rè.})
\item \zh{我妹妹跟我不一样，她不喜欢运动。} (\emph{Wǒ mèimei gēn wǒ bù yīyàng, tā bù xǐhuan yùndòng.})
\end{enumerate}
\end{exercice}

\begin{exercice}[Transforme les phrases]
Transforme chaque phrase affirmative en phrase négative avec \zh{不一样} ou \zh{不同}, puis traduis les deux versions.

Exemple : \zh{我跟他一样高。} → \zh{我跟他不一样高。} « Je ne suis pas aussi grand que lui. »
\begin{enumerate}
\item \zh{这本书跟那本书一样有意思。} (\emph{Zhè běn shū gēn nà běn shū yīyàng yǒu yìsi.})
\item \zh{我的房间与弟弟的房间相同。} (\emph{Wǒ de fángjiān yǔ dìdi de fángjiān xiāngtóng.})
\item \zh{汉语和法语一样难。} (\emph{Hànyǔ hé Fǎyǔ yīyàng nán.})
\end{enumerate}
\end{exercice}

\begin{exercice}[Texte à trous : les mots-outils]
Complète le texte suivant avec les mots de la liste : \zh{跟} — \zh{与} — \zh{一样} — \zh{不同} — \zh{和} — \zh{相同}.

\begin{textebox}[Mon ami chinois]
我有一个中国朋友，他叫王明。他\_\_\_我\_\_\_大，我们都是十八岁。\\
(\emph{Wǒ yǒu yí ge Zhōngguó péngyou, tā jiào Wáng Míng. Tā \_\_\_ wǒ \_\_\_ dà, wǒmen dōu shì shíbā suì.})\\
他的爱好\_\_\_我的爱好\_\_\_：我们都喜欢足球。\\
(\emph{Tā de àihào \_\_\_ wǒ de àihào \_\_\_ : wǒmen dōu xǐhuan zúqiú.})\\
可是我们的学校\_\_\_：他的学校在北京，我的学校在雅温得。\\
(\emph{Kěshì wǒmen de xuéxiào \_\_\_ : tā de xuéxiào zài Běijīng, wǒ de xuéxiào zài Yǎwēndé.})\\
北京\_\_\_雅温得的天气不\_\_\_：北京冬天很冷，雅温得不冷。\\
(\emph{Běijīng \_\_\_ Yǎwēndé de tiānqì bù \_\_\_ : Běijīng dōngtiān hěn lěng, Yǎwēndé bù lěng.})
\end{textebox}
\end{exercice}

\begin{exercice}[Remets les mots dans l'ordre]
Remets les mots dans le bon ordre pour former une phrase correcte, puis traduis-la.
\begin{enumerate}
\item 一样 / 汉语 / 跟 / 英语 / 难 / 不
\item 我的 / 与 / 相同 / 意见 / 老师 / 的
\item 他 / 和 / 一样 / 我 / 高 / 不
\item 跟 / 杜阿拉 / 一样 / 雅温得 / 大 / 吗 / ？
\end{enumerate}
\end{exercice}

\begin{exercice}[Corrige les phrases fautives]
Chaque phrase contient une erreur typique des francophones. Identifie l'erreur, explique-la en français, puis réécris la phrase correctement avec le pinyin.
\begin{enumerate}
\item \zh{我不跟他一样高。} (sens voulu : « Je ne suis pas aussi grand que lui. »)
\item \zh{我一样跟他喜欢足球。} (sens voulu : « J'aime le football comme lui. »)
\item \zh{她跟我一样是很高。} (sens voulu : « Elle est aussi grande que moi. »)
\end{enumerate}
\end{exercice}

\courssub{Je me prépare au Bac}

\begin{exercice}[Compréhension de texte et questions de langue]
Lis le texte suivant, puis réponds aux questions.

\begin{textebox}[Deux villes, deux vies]
我叫恩东，我住在杜阿拉。我的表哥住在雅温得。\\
(\emph{Wǒ jiào Ēndōng, wǒ zhù zài Dù'ālā. Wǒ de biǎogē zhù zài Yǎwēndé.})\\
杜阿拉跟雅温得不一样：杜阿拉很热，雅温得的天气与杜阿拉不同，那里不太热。\\
(\emph{Dù'ālā gēn Yǎwēndé bù yīyàng : Dù'ālā hěn rè, Yǎwēndé de tiānqì yǔ Dù'ālā bùtóng, nàli bú tài rè.})\\
可是我们的生活相同：我们都早上七点去学校，下午五点回家。\\
(\emph{Kěshì wǒmen de shēnghuó xiāngtóng : wǒmen dōu zǎoshang qī diǎn qù xuéxiào, xiàwǔ wǔ diǎn huí jiā.})\\
我的表哥跟我一样喜欢足球，我们也一样喜欢恩多雷。\\
(\emph{Wǒ de biǎogē gēn wǒ yīyàng xǐhuan zúqiú, wǒmen yě yīyàng xǐhuan ndolé.})
\end{textebox}

\textbf{A. Questions de compréhension} (réponds en français)
\begin{enumerate}
\item Où habitent le narrateur et son cousin ?
\item Quelle différence y a-t-il entre les deux villes ?
\item Qu'est-ce que les deux cousins ont en commun ?
\end{enumerate}

\textbf{B. Questions de langue}
\begin{enumerate}
\item Relève dans le texte deux structures de comparaison différentes et donne leur schéma (par exemple : A + 跟 + B + 一样).
\item Dans la phrase \zh{杜阿拉跟雅温得不一样}, pourquoi \zh{不} est-il placé avant \zh{一样} et non avant \zh{跟} ?
\item Transforme en phrase négative : \zh{我们的生活相同。}
\end{enumerate}
\end{exercice}

\begin{exercice}[Thème et version]
\textbf{A. Thème.} Traduis en chinois (caractères et pinyin) les phrases suivantes :
\begin{enumerate}
\item Mon frère est aussi grand que mon père, mais il n'est pas aussi fort que lui.
\item Le chinois n'est pas aussi difficile que le croient les élèves.
\item Mon école est identique à celle de mon correspondant chinois.
\end{enumerate}

\textbf{B. Version.} Traduis en français le texte suivant :

\begin{textebox}[Version]
我的学校跟他的学校不一样大：我的学校很大，他的学校比较小。\\
(\emph{Wǒ de xuéxiào gēn tā de xuéxiào bù yīyàng dà : wǒ de xuéxiào hěn dà, tā de xuéxiào bǐjiào xiǎo.})\\
可是我们的老师与他们的老师相同：都很认真。\\
(\emph{Kěshì wǒmen de lǎoshī yǔ tāmen de lǎoshī xiāngtóng : dōu hěn rènzhēn.})\\
他和我不一样，他喜欢数学，我喜欢汉语。\\
(\emph{Tā hé wǒ bù yīyàng, tā xǐhuan shùxué, wǒ xǐhuan Hànyǔ.})\\
我们跟中国朋友一样努力学习。\\
(\emph{Wǒmen gēn Zhōngguó péngyou yīyàng nǔlì xuéxí.})
\end{textebox}
\end{exercice}

\begin{exercice}[Expression écrite guidée]
\textbf{Sujet :} Tu as un correspondant chinois qui s'appelle 李明 (\emph{Lǐ Míng}). Il te demande de lui décrire ton école. Rédige un paragraphe de \textbf{6 à 8 phrases} (environ \textbf{80 à 120 caractères}) comparant ton école et son école.

\textbf{Contraintes obligatoires :}
\begin{itemize}
\item utiliser au moins \textbf{deux structures différentes} de la leçon parmi : \zh{A 跟 B 一样 + adj.}, \zh{A 与 B 相同}, \zh{A 和 B 不一样 + adj.}, \zh{A 与 B 不同} ;
\item employer le vocabulaire de la vie scolaire : \zh{学校、老师、学生、教室、课、运动} ;
\item écrire en caractères simplifiés, avec le pinyin sous chaque phrase.
\end{itemize}

\textbf{Aide au démarrage :} \zh{我的学校跟你的学校……} (\emph{Wǒ de xuéxiào gēn nǐ de xuéxiào \ldots\ldots})
\end{exercice}

\newpage

\coursec{Corrigés des exercices}

\begin{corrige}[Exercice 1 — QCM : la bonne particule]
\begin{enumerate}
\item \textbf{B. 与} — \zh{这本书与那本书相同。} (\emph{Zhè běn shū yǔ nà běn shū xiāngtóng.}) « Ce livre est identique à celui-là. » La structure \zh{A 与 B 相同} exige la préposition de comparaison \zh{与} ; \zh{很} et \zh{也} sont des adverbes, \zh{得} est un mot-outil de complément.
\item \textbf{B. 跟} — \zh{我跟我哥哥一样高。} (\emph{Wǒ gēn wǒ gēge yīyàng gāo.}) « Je suis aussi grand que mon grand frère. » Schéma : A + 跟 + B + 一样 + adjectif.
\item \textbf{A. 和} — \zh{汉语和法语不一样难。} (\emph{Hànyǔ hé Fǎyǔ bù yīyàng nán.}) « Le chinois n'est pas aussi difficile que le français. » Négation : 不 se place avant \zh{一样}.
\item \textbf{C. 与} — \zh{杜阿拉的天气与雅温得的天气不同。} (\emph{Dù'ālā de tiānqì yǔ Yǎwēndé de tiānqì bùtóng.}) « Le temps de Douala est différent de celui de Yaoundé. » \zh{不同} est invariable : on ne met jamais \zh{不} devant (il contient déjà la négation).
\end{enumerate}
\end{corrige}

\begin{corrige}[Exercice 2 — Vrai ou faux ?]
\begin{enumerate}
\item \textbf{V.} \zh{我跟他一样高。} respecte le schéma A + 跟 + B + 一样 + adjectif. Traduction correcte.
\item \textbf{F.} Ordre des mots fautif : la préposition \zh{跟} et son complément doivent précéder \zh{一样}. Correction : \zh{我跟他一样喜欢足球。} (\emph{Wǒ gēn tā yīyàng xǐhuan zúqiú.})
\item \textbf{V.} \zh{我的学校与他的学校相同。} suit le schéma A + 与 + B + 相同, registre soutenu, correct.
\item \textbf{V.} \zh{她不跟我一样喜欢音乐。} : la négation \zh{不} est bien placée avant \zh{跟} quand on nie le verbe qui suit (elle n'aime pas autant que moi). Attention : ce sens diffère de \zh{她跟我一样不喜欢音乐} (« comme moi, elle n'aime pas la musique »).
\end{enumerate}
\end{corrige}

\begin{corrige}[Exercice 3 — Vocabulaire : complète le tableau]
\begin{center}
\begin{tabularx}{\linewidth}{|X|X|X|X|}
\hline
\rowcolor{popblueL} Caractères & Pinyin & Nature & Traduction \\
\hline
一样 & yīyàng & adjectif / adverbe & pareil, identique \\
\hline
相同 & xiāngtóng & adjectif & identique, semblable \\
\hline
不同 & bùtóng & adjectif & différent \\
\hline
跟 & gēn & préposition / conjonction & avec, comme (comparaison) \\
\hline
与 & yǔ & préposition / conjonction & avec, et (soutenu) \\
\hline
比较 & bǐjiào & adverbe / verbe & relativement ; comparer \\
\hline
\end{tabularx}
\end{center}
Point de vigilance : \zh{与} appartient au registre écrit et soutenu ; à l'oral on préfère \zh{跟} ou \zh{和}.
\end{corrige}

\begin{corrige}[Exercice 4 — Associe le pinyin aux caractères]
\begin{enumerate}
\item \zh{跟……一样} → \textbf{b.} \emph{gēn \ldots\ldots yīyàng}
\item \zh{和……不同} → \textbf{c.} \emph{hé \ldots\ldots bùtóng}
\item \zh{与……相同} → \textbf{d.} \emph{yǔ \ldots\ldots xiāngtóng}
\item \zh{跟……不一样} → \textbf{e.} \emph{gēn \ldots\ldots bù yīyàng}
\end{enumerate}
L'option \textbf{a.} \emph{yǔ \ldots\ldots bùtóng} correspondrait à \zh{与……不同}, qui n'était pas proposée : c'est le distracteur.
\end{corrige}

\begin{corrige}[Exercice 5 — Traduis en français]
\begin{enumerate}
\item \zh{喀麦隆的足球跟巴西的足球一样有名。} « Le football camerounais est aussi célèbre que le football brésilien. »
\item \zh{我的想法与你的相同。} « Mon opinion est identique à la tienne. » (\zh{你的} = « la tienne », le nom \zh{想法} est sous-entendu.)
\item \zh{今天的天气和昨天不一样热。} « Il ne fait pas aussi chaud aujourd'hui qu'hier. »
\item \zh{我妹妹跟我不一样，她不喜欢运动。} « Ma petite sœur n'est pas comme moi : elle n'aime pas le sport. »
\end{enumerate}
\end{corrige}

\begin{corrige}[Exercice 6 — Transforme les phrases]
\begin{enumerate}
\item Affirmative : \zh{这本书跟那本书一样有意思。} « Ce livre est aussi intéressant que celui-là. »\\
Négative : \zh{这本书跟那本书不一样有意思。} (\emph{Zhè běn shū gēn nà běn shū bù yīyàng yǒu yìsi.}) « Ce livre n'est pas aussi intéressant que celui-là. »
\item Affirmative : \zh{我的房间与弟弟的房间相同。} « Ma chambre est identique à celle de mon petit frère. »\\
Négative : \zh{我的房间与弟弟的房间不同。} (\emph{Wǒ de fángjiān yǔ dìdi de fángjiān bùtóng.}) « Ma chambre est différente de celle de mon petit frère. »
\item Affirmative : \zh{汉语和法语一样难。} « Le chinois est aussi difficile que le français. »\\
Négative : \zh{汉语和法语不一样难。} (\emph{Hànyǔ hé Fǎyǔ bù yīyàng nán.}) « Le chinois n'est pas aussi difficile que le français. »
\end{enumerate}
Règle appliquée : \zh{一样} se nie par \zh{不一样} ; \zh{相同} se nie par \zh{不同} (jamais \zh{不相同}, forme lourde à éviter à ce niveau).
\end{corrige}

\begin{corrige}[Exercice 7 — Texte à trous : les mots-outils]
Texte complété :\\
\zh{他跟我一样大，我们都是十八岁。} (\emph{Tā gēn wǒ yīyàng dà, wǒmen dōu shì shíbā suì.}) « Il a le même âge que moi, nous avons tous les deux dix-huit ans. »\\
\zh{他的爱好和我的爱好相同：我们都喜欢足球。} (\emph{Tā de àihào hé wǒ de àihào xiāngtóng.}) « Ses goûts et les miens sont identiques. »\\
\zh{可是我们的学校不同。} (\emph{Kěshì wǒmen de xuéxiào bùtóng.}) « Mais nos écoles sont différentes. »\\
\zh{北京与雅温得的天气不一样。} (\emph{Běijīng yǔ Yǎwēndé de tiānqì bù yīyàng.}) « Le temps de Pékin n'est pas le même que celui de Yaoundé. »

Traduction globale : « J'ai un ami chinois, il s'appelle Wang Ming. Il a le même âge que moi, nous avons tous les deux dix-huit ans. Ses goûts sont identiques aux miens : nous aimons tous les deux le football. Mais nos écoles sont différentes : son école est à Pékin, la mienne à Yaoundé. Le temps de Pékin n'est pas le même qu'à Yaoundé : l'hiver, il fait très froid à Pékin, il ne fait pas froid à Yaoundé. »
\end{corrige}

\begin{corrige}[Exercice 8 — Remets les mots dans l'ordre]
\begin{enumerate}
\item \zh{汉语不跟英语一样难。} ou mieux \zh{汉语跟英语不一样难。} (\emph{Hànyǔ gēn Yīngyǔ bù yīyàng nán.}) « Le chinois n'est pas aussi difficile que l'anglais. »
\item \zh{我的意见与老师的相同。} (\emph{Wǒ de yìjiàn yǔ lǎoshī de xiāngtóng.}) « Mon avis est identique à celui du professeur. »
\item \zh{他和我（不）一样高} → \zh{他和我（不）一样高} : la bonne phrase est \zh{他和我（不）一样高} ; avec 不 : \zh{他和我不一样高。} (\emph{Tā hé wǒ bù yīyàng gāo.}) « Il n'est pas aussi grand que moi. »
\item \zh{杜阿拉跟雅温得一样大吗？} (\emph{Dù'ālā gēn Yǎwēndé yīyàng dà ma?}) « Douala est-elle aussi grande que Yaoundé ? »
\end{enumerate}
Point de vigilance : le mot interrogatif \zh{吗} se place toujours en fin de phrase.
\end{corrige}

\begin{corrige}[Exercice 9 — Corrige les phrases fautives]
\begin{enumerate}
\item Faute : \zh{我不跟他一样高。} place \zh{不} avant \zh{跟}, ce qui nie le fait d'être « avec lui » (sens bizarre). En chinois, on nie \zh{一样}, pas la préposition. Correction : \zh{我跟他不一样高。} (\emph{Wǒ gēn tā bù yīyàng gāo.}) « Je ne suis pas aussi grand que lui. »
\item Faute : \zh{我一样跟他喜欢足球。} calque l'ordre français « je, comme lui, aime… ». Le groupe prépositionnel \zh{跟他} doit précéder \zh{一样}. Correction : \zh{我跟他一样喜欢足球。} (\emph{Wǒ gēn tā yīyàng xǐhuan zúqiú.})
\item Faute : \zh{她跟我一样是很高。} ajoute \zh{是} par calque du français « elle est aussi grande ». Après \zh{一样}, l'adjectif seul suffit, sans copule. Correction : \zh{她跟我一样高。} (\emph{Tā gēn wǒ yīyàng gāo.}) « Elle est aussi grande que moi. »
\end{enumerate}
\end{corrige}

\begin{corrige}[Exercice 10 — Compréhension de texte et questions de langue]
\textbf{A. Questions de compréhension}
\begin{enumerate}
\item Le narrateur, Ndong (\zh{恩东}), habite à Douala ; son cousin habite à Yaoundé.
\item La différence porte sur le climat : il fait très chaud à Douala, alors qu'à Yaoundé il ne fait pas trop chaud.
\item Ils ont en commun leur emploi du temps (école à 7 h, retour à 17 h), leur goût pour le football et pour le ndolé.
\end{enumerate}

\textbf{B. Questions de langue}
\begin{enumerate}
\item Exemples relevés : \zh{杜阿拉跟雅温得不一样} (schéma A + 跟 + B + 不一样) ; \zh{雅温得的天气与杜阿拉不同} (schéma A + 与 + B + 不同) ; \zh{我们的生活相同} (A + 相同) ; \zh{我的表哥跟我一样喜欢足球} (A + 跟 + B + 一样 + verbe).
\item Dans \zh{杜阿拉跟雅温得不一样}, \zh{不} nie l'adjectif \zh{一样} (« pas pareil »), c'est donc lui qui est nié, pas la préposition \zh{跟}. Dire \zh{不跟……一样} reviendrait à nier la comparaison elle-même, ce qui est fautif dans ce sens.
\item \zh{我们的生活相同。} → \zh{我们的生活不同。} (\emph{Wǒmen de shēnghuó bùtóng.}) « Nos vies sont différentes. »
\end{enumerate}
\textbf{Barème indicatif (10 pts)} : A 3 × 1 pt ; B1 2 pts ; B2 3 pts ; B3 2 pts.
\textbf{Points de vigilance} : ne pas traduire \zh{一样喜欢足球} par « aiment pareillement le football » de façon lourde ; repérer que \zh{相同} et \zh{不同} s'emploient sans adjectif derrière.
\end{corrige}

\begin{corrige}[Exercice 11 — Thème et version]
\textbf{A. Thème}
\begin{enumerate}
\item \zh{我哥哥跟我爸爸一样高，可是他跟爸爸不一样强壮。} (\emph{Wǒ gēge gēn wǒ bàba yīyàng gāo, kěshì tā gēn bàba bù yīyàng qiángzhuàng.})
\item \zh{汉语跟学生们想的不一样难。} (\emph{Hànyǔ gēn xuéshengmen xiǎng de bù yīyàng nán.})
\item \zh{我的学校与我中国笔友的学校相同。} (\emph{Wǒ de xuéxiào yǔ wǒ Zhōngguó bǐyǒu de xuéxiào xiāngtóng.})
\end{enumerate}

\textbf{B. Version}\\
« Mon école n'est pas de la même taille que la sienne : mon école est très grande, la sienne est relativement petite. Mais nos professeurs sont identiques aux leurs : ils sont tous très sérieux. Lui et moi ne sommes pas pareils : il aime les mathématiques, moi j'aime le chinois. Nous travaillons aussi dur que nos amis chinois. »

\textbf{Barème indicatif (10 pts)} : thème 3 × 2 pts (structure 1 pt, vocabulaire et caractères 1 pt) ; version 4 pts (1 pt par phrase).
\textbf{Points de vigilance} : « aussi… que » = \zh{跟……一样} ; « pas aussi… que » = \zh{跟……不一样} ; ne jamais mettre \zh{是} devant l'adjectif ; \zh{比较小} = « relativement petit » (comparatif atténué sans \zh{比} explicite).
\end{corrige}

\begin{corrige}[Exercice 12 — Expression écrite guidée]
\textbf{Proposition de rédaction modèle} (environ 100 caractères) :

\zh{我的学校跟你的学校不一样大：我的学校很大，有两千个学生。}\\
(\emph{Wǒ de xuéxiào gēn nǐ de xuéxiào bù yīyàng dà : wǒ de xuéxiào hěn dà, yǒu liǎng qiān ge xuésheng.})\\
\zh{我们的教室与你们的教室相同，都很干净。}\\
(\emph{Wǒmen de jiàoshì yǔ nǐmen de jiàoshì xiāngtóng, dōu hěn gānjìng.})\\
\zh{我的老师跟你的老师一样认真。}\\
(\emph{Wǒ de lǎoshī gēn nǐ de lǎoshī yīyàng rènzhēn.})\\
\zh{可是我们的课和你们不同：我们学习法语、英语和汉语。}\\
(\emph{Kěshì wǒmen de kè hé nǐmen bùtóng : wǒmen xuéxí Fǎyǔ, Yīngyǔ hé Hànyǔ.})\\
\zh{我们都一样喜欢运动，我最喜欢足球。}\\
(\emph{Wǒmen dōu yīyàng xǐhuan yùndòng, wǒ zuì xǐhuan zúqiú.})\\
\zh{你的学校怎么样？} (\emph{Nǐ de xuéxiào zěnmeyàng?})

Traduction : « Mon école n'est pas de la même taille que la tienne : elle est très grande, avec deux mille élèves. Nos salles de classe sont identiques aux vôtres, toutes très propres. Mes professeurs sont aussi sérieux que les tiens. Mais nos cours sont différents des vôtres : nous étudions le français, l'anglais et le chinois. Nous aimons tous autant le sport, moi je préfère le football. Et ton école, comment est-elle ? »

\textbf{Barème indicatif (10 pts)} : respect des contraintes (2 structures différentes, vocabulaire scolaire) 3 pts ; correction grammaticale des structures de comparaison 3 pts ; caractères et pinyin 2 pts ; cohérence et longueur 2 pts.
\textbf{Points de vigilance} : vérifier la place de \zh{不} (devant \zh{一样}) ; ne pas écrire \zh{是} + adjectif après \zh{一样} ; varier \zh{跟/和/与} ; terminer par une question pour relancer la correspondance.
\end{corrige}

\newpage

\coursec{Fiche bilan}

\begin{popfigure}
\begin{tikzpicture}[mindmap, grow cyclic, every node/.style={concept, align=center, text=white, font=\small}, concept color=popblue, level 1/.append style={level distance=3.4cm, sibling angle=60}, text width=2.6cm]
\node[concept color=popdark, text width=3cm, font=\small\bfseries] {Le comparatif\\ 一样 / 不一样\\ (Tle A)}
  child[concept color=popblue] { node[align=center]{Structure d'égalité\\ A 跟 B 一样\\ + adjectif} }
  child[concept color=poporange] { node[align=center]{Connecteurs\\ 跟 / 和 / 与\\ (même sens)} }
  child[concept color=popgreen] { node[align=center]{Synonymes\\ 一样 = 相同\\ 不一样 = 不同} }
  child[concept color=poppurple] { node[align=center]{Différence\\ A 跟 B 不一样\\ + adjectif} }
  child[concept color=poppink] { node[align=center]{Négation\\ 不一样 / 不相同\\ jamais de 不 + 一样} }
  child[concept color=popgold] { node[align=center]{Pièges\\ ordre des mots\\ pas de 比 ici} }
  child[concept color=popmuted] { node[align=center]{Emplois\\ décrire, comparer\\ famille, école, pays} };
\end{tikzpicture}
\legende{Carte mentale de la leçon : exprimer l'égalité et la différence avec 跟/和/与 … 一样/不一样.}
\end{popfigure}

\begin{bilan}
\textbf{1. Lexique clé à connaître par cœur}

\begin{center}
\begin{tabularx}{\linewidth}{|X|X|X|X|}\hline
\rowcolor{popblueL} Caractères & Pinyin & Nature & Traduction \\ \hline
跟 & gēn & préposition & avec, comparé à \\ \hline
和 & hé & préposition/conjonction & avec, et \\ \hline
与 & yǔ & préposition (soutenu) & avec \\ \hline
一样 & yíyàng & adjectif & pareil, identique \\ \hline
相同 & xiāngtóng & adjectif & identique, semblable \\ \hline
不同 & bùtóng & adjectif & différent \\ \hline
比较 & bǐjiào & verbe/adverbe & comparer ; assez \\ \hline
差不多 & chàbuduō & adjectif/adverbe & à peu près pareil \\ \hline
\end{tabularx}
\end{center}

\textbf{2. Règles de grammaire à connaître par cœur}

\begin{center}
\begin{tabularx}{\linewidth}{|X|X|X|}\hline
\rowcolor{popblueL} Structure & Exemple (caractères + pinyin) & Traduction \\ \hline
A 跟 B 一样 (+ adj.) & 我跟我弟弟一样高。Wǒ gēn wǒ dìdi yíyàng gāo. & Je suis aussi grand que mon petit frère. \\ \hline
A 和 B 相同 & 他的想法和我的相同。Tā de xiǎngfa hé wǒ de xiāngtóng. & Son idée est identique à la mienne. \\ \hline
A 跟 B 不一样 (+ adj.) & 杜阿拉跟雅温得不一样大。Dùālā gēn Yǎwēndé bù yíyàng dà. & Douala n'est pas aussi grande que Yaoundé. \\ \hline
A 与 B 不同 & 中国的春节与喀麦隆的新年不同。Zhōngguó de Chūnjié yǔ Kāmàilóng de xīnnián bùtóng. & La fête du Printemps chinoise est différente du Nouvel An camerounais. \\ \hline
A 跟 B 差不多 & 我的成绩跟他差不多。Wǒ de chéngji gēn tā chàbuduō. & Mes résultats sont à peu près comme les siens. \\ \hline
\end{tabularx}
\end{center}

\begin{itemize}
\item \cle{跟}, \cle{和}, \cle{与} sont interchangeables dans cette structure ; \zh{与} est plus écrit et soutenu.
\item \cle{一样} et \cle{相同} signifient « pareil » ; la négation est \zh{不一样} ou \zh{不同}, jamais \zh{不一样} avec un autre négatif.
\item L'adjectif se place \textbf{après} \zh{一样} : \zh{A 跟 B 一样 + adjectif}.
\item Ne pas confondre avec le comparatif de supériorité \zh{A 比 B + adjectif} : ici on exprime l'\textbf{égalité} ou la \textbf{différence}, pas « plus que ».
\item Ordre des mots figé : le mot-outil \zh{跟/和/与} se place \textbf{avant} B, jamais après l'adjectif.
\end{itemize}

\textbf{3. Méthodes express}
\begin{enumerate}
\item Identifier ce qu'on compare (A et B) et le point de comparaison (l'adjectif).
\item Choisir : égalité $\rightarrow$ \zh{一样/相同} ; différence $\rightarrow$ \zh{不一样/不同}.
\item Construire : A + 跟/和/与 + B + 一样 (+ adjectif).
\item Vérifier : négation bien placée, pas de \zh{比} dans la phrase, adjectif après \zh{一样}.
\end{enumerate}
\end{bilan}

\begin{aretenir}[Checklist avant le Bac]
\begin{itemize}
\item[$\square$] Je connais les trois connecteurs \zh{跟}, \zh{和}, \zh{与} et leur place dans la phrase.
\item[$\square$] Je sais que \zh{一样} et \zh{相同} expriment l'égalité.
\item[$\square$] Je sais que la négation est \zh{不一样} ou \zh{不同}.
\item[$\square$] Je place correctement l'adjectif après \zh{一样}.
\item[$\square$] Je distingue \zh{A 跟 B 一样} (égalité) de \zh{A 比 B + adj.} (supériorité).
\item[$\square$] Je n'utilise jamais \zh{很} devant \zh{一样}.
\item[$\square$] Je sais former une phrase de différence avec un adjectif : \zh{A 跟 B 不一样 + adj.}
\item[$\square$] Je sais utiliser \zh{差不多} pour « à peu près pareil ».
\item[$\square$] Je respecte l'ordre figé : A + 跟 + B + 一样 (+ adjectif).
\item[$\square$] Je sais traduire du français vers le chinois « aussi … que » et « pas aussi … que ».
\item[$\square$] Je peux comparer deux réalités du Cameroun et de la Chine en trois phrases correctes.
\item[$\square$] Je relis ma phrase pour vérifier les tons du pinyin et la ponctuation chinoise.
\end{itemize}
\end{aretenir}

\findecours

\end{document}
